% Agriculture — Biology Upper Sixth — Cours signé OnBuch+
% Official MINESEC syllabus. Compiler avec : tectonic BIO08-agriculture.tex
\newcommand{\DOCMATIERE}{Biology}
\newcommand{\DOCNIVEAU}{Upper Sixth}
\newcommand{\DOCMODULE}{Upper Sixth — Biology}
\newcommand{\DOCLECON}{8}
\newcommand{\DOCTITRE}{Agriculture}
% OnBuch+ — préambule « cours signés » ENRICHI (compilé avec tectonic / XeLaTeX)
% Système visuel « Pop » d'OnBuch+ : fond crème (jamais blanc), bordures encre
% épaisses, ombres « dures » décalées, titres Archivo Black en capitales.
% Version MATHÉMATIQUES : graphes (pgfplots/tikz), boîtes pédagogiques
% (définition, propriété, méthode, attention, le savais-tu, exercice résolu,
% exercice, TP, bilan), page de garde et signature OnBuch+ ancrée en bas.
%
% Macros attendues dans main.tex AVANT \input{preamble} :
%   \DOCMATIERE  \DOCNIVEAU  \DOCMODULE  \DOCLECON (numéro)  \DOCTITRE
\documentclass[11pt]{article}

\usepackage{fontspec}
\usepackage[english]{babel}
\usepackage[a4paper,margin=2cm,top=3.4cm,bottom=2.8cm,headheight=1.4cm,footskip=1.3cm]{geometry}
\usepackage{amsmath,amssymb,amsfonts}
\usepackage{mathtools} % \xmapsto, \coloneqq... souvent utilisés par les modèles
\usepackage{cancel} % simplifications barrées (bilans de forces)
% \abs{z} (module, valeur absolue) et \norm{u} (norme), utilisés par les modèles
\providecommand{\abs}{}\let\abs\relax\DeclarePairedDelimiter{\abs}{\lvert}{\rvert}
\providecommand{\norm}{}\let\norm\relax\DeclarePairedDelimiter{\norm}{\lVert}{\rVert}
\usepackage[table]{xcolor} % [table] : fournit aussi \rowcolors (alternance de lignes)
\usepackage{pagecolor}
\usepackage{tikz}
\usetikzlibrary{arrows.meta,positioning,calc,shapes.geometric,shapes.misc,shapes.arrows,decorations.pathmorphing,decorations.pathreplacing,decorations.markings,patterns,shadows,mindmap,trees,fit,backgrounds,matrix,intersections,angles,quotes}
\usepackage{pgfplots}
\usepackage[version=4]{mhchem} % équations nucléaires \ce{^{235}_{92}U}
\usepackage[european,siunitx]{circuitikz} % schémas électriques
\pgfplotsset{compat=1.18}
\usepgfplotslibrary{fillbetween,groupplots} % aires entre courbes (calcul intégral)
\usepackage{enumitem}
\usepackage{fancyhdr}
% [most] charge toutes les bibliothèques tcolorbox (listings, minted, raster,
% documentation...) et gonfle inutilement la mémoire TeX sur un document long
% (~300 boîtes) : on ne charge que ce qui est réellement utilisé.
\usepackage[skins,breakable]{tcolorbox}
\usepackage{graphicx}
\usepackage{colortbl}
\usepackage{booktabs}
\usepackage{tabularx}
\usepackage{diagbox} % case d'en-tête diagonale des tableaux à double entrée
% Colonnes extensibles centrée / gauche / droite (C, L, R), souvent utilisées par les modèles
\newcolumntype{C}{>{\centering\arraybackslash}X}
\newcolumntype{L}{>{\raggedright\arraybackslash}X}
\newcolumntype{R}{>{\raggedleft\arraybackslash}X}
\usepackage{array}
\usepackage{multicol}
\usepackage{siunitx}
% parse-numbers=false : accepte aussi \SI{2,5 \times 10^{-3}}{...}
\sisetup{locale=UK,output-decimal-marker={.},group-minimum-digits=4,parse-numbers=false}
\DeclareSIUnit\ohm{\ensuremath{\symup{\Omega}}} % U+2126 absent de la police
\DeclareSIUnit\division{div} % réglages d'oscilloscope (ms/div, V/div)
\DeclareSIUnit\tr{tr} % tours (vitesse de rotation, ex. \SI{1500}{\tr\per\minute})
\usepackage{qrcode}
% \degree : commande fréquemment utilisée par les modèles pour un angle ou une
% température en dehors de \SI{}{} (non fournie par les paquets chargés).
\providecommand{\degree}{^{\circ}}
% \qed (fin de démonstration), utilisé par les modèles sans amsthm
\providecommand{\qed}{\hfill$\square$}
% \barre : double barre des valeurs interdites dans un tableau de variations (tkz-tab)
\providecommand{\barre}{\ensuremath{\|}}
% environnement proof (amsthm non chargé) utilisé par les modèles
\ifdefined\proof\else\newenvironment{proof}[1][Proof]{\par\noindent\textit{#1.}\ }{\qed\par}\fi
\usepackage{lastpage}
\usepackage{hyperref}
\hypersetup{hidelinks}

% ============================================================
% Palette « Pop » OnBuch+ — mêmes hex que la classe P (plus_ui.dart)
% ============================================================
\definecolor{poporange}{HTML}{F59321}
\definecolor{popdark}{HTML}{E07A0C}
\definecolor{popcream}{HTML}{FFF6E8}
\definecolor{popink}{HTML}{1C1714}
\definecolor{poppurple}{HTML}{7A5AE0}
\definecolor{popgreen}{HTML}{2E9E6A}
\definecolor{poppink}{HTML}{E0457B}
\definecolor{popblue}{HTML}{2F7DE1}
\definecolor{popgold}{HTML}{C98A12}
\definecolor{popmuted}{HTML}{8A7F76}
\definecolor{popline}{HTML}{EADFD2}
\definecolor{popgray}{HTML}{8A7F76}
\colorlet{popgrey}{popgray}
% Alias tolérants (le modèle nomme parfois les couleurs différemment)
\colorlet{popwhite}{white}
\colorlet{popblack}{popink}
\colorlet{popred}{poppink}
\colorlet{popyellow}{popgold}
% Teintes claires pour fonds de tableaux / zones
\colorlet{popblueL}{popblue!10!white}
\colorlet{poporangeL}{poporange!14!white}
\colorlet{popgreenL}{popgreen!12!white}
\colorlet{poppurpleL}{poppurple!12!white}
\colorlet{poppinkL}{poppink!12!white}
\colorlet{popgoldL}{popgold!14!white}

\pagecolor{popcream}

% ============================================================
% Typographie
% ============================================================
\defaultfontfeatures{Ligatures=TeX}
\setmainfont{PlusJakartaSans-Medium.ttf}[Path=../../../fonts/,BoldFont=PlusJakartaSans-Bold.ttf]
% Maths en sans-serif (Fira Math) avec chiffres et lettres droites de Plus
% Jakarta, pour que les formules se fondent dans le texte.
\usepackage[math-style=ISO,bold-style=ISO]{unicode-math}
\setmathfont{FiraMath-Regular.otf}[Path=../../../fonts/]
\setmathfont{PlusJakartaSans-Medium.ttf}[Path=../../../fonts/,range={up/num,up/{Latin,latin},\mathup}]
% (icomma retiré : décimales avec point en anglais)
% Caractères mathématiques Unicode absents des polices de texte (Plus Jakarta,
% Archivo Black) : rendus par leur équivalent mathématique, en texte comme en
% formule — sinon ils s'affichent en carré vide (ex. « Arithmétique dans ℤ »).
\usepackage{newunicodechar}
\newunicodechar{ℝ}{\ensuremath{\mathbb{R}}}
\newunicodechar{ℤ}{\ensuremath{\mathbb{Z}}}
\newunicodechar{ℂ}{\ensuremath{\mathbb{C}}}
\newunicodechar{ℕ}{\ensuremath{\mathbb{N}}}
\newunicodechar{ℚ}{\ensuremath{\mathbb{Q}}}
\newunicodechar{𝔻}{\ensuremath{\mathbb{D}}}
\newunicodechar{α}{\ensuremath{\alpha}}
\newunicodechar{β}{\ensuremath{\beta}}
\newunicodechar{γ}{\ensuremath{\gamma}}
\newunicodechar{δ}{\ensuremath{\delta}}
\newunicodechar{ε}{\ensuremath{\varepsilon}}
\newunicodechar{θ}{\ensuremath{\theta}}
\newunicodechar{λ}{\ensuremath{\lambda}}
\newunicodechar{μ}{\ensuremath{\mu}}
\newunicodechar{σ}{\ensuremath{\sigma}}
\newunicodechar{τ}{\ensuremath{\tau}}
\newunicodechar{φ}{\ensuremath{\varphi}}
\newunicodechar{ψ}{\ensuremath{\psi}}
\newunicodechar{ω}{\ensuremath{\omega}}
\newunicodechar{Σ}{\ensuremath{\Sigma}}
% majuscules produites par \MakeUppercase dans les titres (α-aminés -> Α)
\newunicodechar{Α}{\ensuremath{\alpha}}
\newunicodechar{Β}{\ensuremath{\beta}}
\newunicodechar{Γ}{\ensuremath{\Gamma}}
\newunicodechar{Δ}{\ensuremath{\Delta}}
\newunicodechar{Ε}{\ensuremath{\varepsilon}}
\newunicodechar{Θ}{\ensuremath{\Theta}}
\newunicodechar{Λ}{\ensuremath{\Lambda}}
\newunicodechar{Μ}{\ensuremath{\mu}}
\newunicodechar{Τ}{\ensuremath{\tau}}
\newunicodechar{Φ}{\ensuremath{\Phi}}
\newunicodechar{Ψ}{\ensuremath{\Psi}}
\newunicodechar{Ω}{\ensuremath{\Omega}}
\newunicodechar{↦}{\ensuremath{\mapsto}}
\newunicodechar{∈}{\ensuremath{\in}}
\newunicodechar{∉}{\ensuremath{\notin}}
\newunicodechar{∧}{\ensuremath{\wedge}}
\newunicodechar{⇔}{\ensuremath{\Leftrightarrow}}
\newunicodechar{⟺}{\ensuremath{\Longleftrightarrow}}
\newunicodechar{⇒}{\ensuremath{\Rightarrow}}
\newunicodechar{⟹}{\ensuremath{\Longrightarrow}}
\newunicodechar{∘}{\ensuremath{\circ}}
\newunicodechar{∪}{\ensuremath{\cup}}
\newunicodechar{∩}{\ensuremath{\cap}}
\newunicodechar{≡}{\ensuremath{\equiv}}
\newunicodechar{⊂}{\ensuremath{\subset}}
\newunicodechar{⊥}{\ensuremath{\perp}}
\newunicodechar{∃}{\ensuremath{\exists}}
\newunicodechar{∀}{\ensuremath{\forall}}
\newunicodechar{∛}{\ensuremath{\sqrt[3]{\,}}}
\newunicodechar{∜}{\ensuremath{\sqrt[4]{\,}}}
\newunicodechar{⃗}{\ensuremath{{}^{\to}}}
\newunicodechar{ⁿ}{\ifmmode^{n}\else\textsuperscript{\ensuremath{n}}\fi}
\newunicodechar{ˣ}{\ifmmode^{x}\else\textsuperscript{\ensuremath{x}}\fi}
\newunicodechar{ᵇ}{\ifmmode^{b}\else\textsuperscript{\ensuremath{b}}\fi}
\newunicodechar{ʳ}{\ifmmode^{r}\else\textsuperscript{\ensuremath{r}}\fi}
\newunicodechar{ᵘ}{\ifmmode^{u}\else\textsuperscript{\ensuremath{u}}\fi}
\newunicodechar{ʸ}{\ifmmode^{y}\else\textsuperscript{\ensuremath{y}}\fi}
\newunicodechar{ᵃ}{\ifmmode^{a}\else\textsuperscript{\ensuremath{a}}\fi}
\newunicodechar{ᵉ}{\ifmmode^{e}\else\textsuperscript{\ensuremath{e}}\fi}
\newunicodechar{ᵏ}{\ifmmode^{k}\else\textsuperscript{\ensuremath{k}}\fi}
\newunicodechar{ᵖ}{\ifmmode^{p}\else\textsuperscript{\ensuremath{p}}\fi}
\newunicodechar{ᴺ}{\ifmmode^{N}\else\textsuperscript{\ensuremath{N}}\fi}
\newunicodechar{ᶜ}{\ifmmode^{c}\else\textsuperscript{\ensuremath{c}}\fi}
\newunicodechar{ʰ}{\ifmmode^{h}\else\textsuperscript{\ensuremath{h}}\fi}
\newunicodechar{ᵠ}{\ifmmode^{q}\else\textsuperscript{\ensuremath{q}}\fi}
\newunicodechar{ₙ}{\ifmmode_{n}\else\textsubscript{\ensuremath{n}}\fi}
\newunicodechar{ₐ}{\ifmmode_{a}\else\textsubscript{\ensuremath{a}}\fi}
\newunicodechar{ᵢ}{\ifmmode_{i}\else\textsubscript{\ensuremath{i}}\fi}
\newunicodechar{ᵣ}{\ifmmode_{r}\else\textsubscript{\ensuremath{r}}\fi}
\newunicodechar{ₖ}{\ifmmode_{k}\else\textsubscript{\ensuremath{k}}\fi}
\newunicodechar{ᵦ}{\ifmmode_{\beta}\else\textsubscript{\ensuremath{\beta}}\fi}
\newunicodechar{ₓ}{\ifmmode_{x}\else\textsubscript{\ensuremath{x}}\fi}
\newfontfamily\popdisplay{ArchivoBlack-Regular.ttf}[Path=../../../fonts/]
\newfontfamily\poplogo{SpaceGrotesk-Bold.ttf}[Path=../../../fonts/]
\newfontfamily\popsemi{PlusJakartaSans-SemiBold.ttf}[Path=../../../fonts/]
\newfontfamily\popxbold{PlusJakartaSans-ExtraBold.ttf}[Path=../../../fonts/]
\newfontfamily\popxboldls{PlusJakartaSans-ExtraBold.ttf}[Path=../../../fonts/,LetterSpace=3.0]

\setlist[enumerate]{leftmargin=1.7em,itemsep=5pt,topsep=4pt}
\setlist[itemize]{leftmargin=1.7em,itemsep=4pt,topsep=4pt,label={\color{poporange}\textbullet}}
\setlength{\parindent}{0pt}
\setlength{\parskip}{5pt}
\renewcommand{\arraystretch}{1.25}

% pgfplots : style Pop par défaut
\pgfplotsset{
  every axis/.append style={axis line style={popink,line width=1pt},tick style={popink},
    grid style={popline,line width=0.5pt},label style={font=\small},tick label style={font=\footnotesize}},
  popcurve/.style={poporange,line width=1.8pt,no markers,smooth},
}
% Même style disponible dans un \draw TikZ simple (hors pgfplots)
\tikzset{popcurve/.style={poporange,line width=1.8pt}}
\tikzset{popgrid/.style={popline,very thin}}
\colorlet{popcurve}{poporange} % le nom du style est parfois utilisé comme couleur

% ============================================================
% Pill / squiggle
% ============================================================
\newcommand{\poppill}[3]{%
  \tikz[baseline=-0.55ex]\node[draw=popink,line width=1.2pt,rounded corners=2.6pt,%
    fill=#2,inner xsep=6.5pt,inner ysep=3pt]{%
    \popxboldls\fontsize{7}{7}\selectfont\color{#3}\MakeUppercase{#1}};%
}
\newcommand{\popsquiggle}[1]{%
  \tikz[baseline=-0.4ex]{
    \draw[#1,line width=2.6pt,line cap=round]
      plot [smooth] coordinates {(0,0.09) (0.34,0.01) (0.68,0.21) (1.02,0.01) (1.36,0.10)};
  }%
}
% Mot-clé surligné (marqueur orange sous le texte)
\newcommand{\cle}[1]{{\popxbold\color{popdark}#1}}

% ============================================================
% En-tête / pied
% ============================================================
\pagestyle{fancy}
\fancyhf{}
\renewcommand{\headrulewidth}{0pt}
\renewcommand{\footrulewidth}{0pt}
\lhead{\raisebox{-0.15ex}{\poplogo\fontsize{13}{13}\selectfont\color{popink}OnBuch\color{poporange}+}}
\rhead{\poppill{\DOCMATIERE\ \textbullet\ \DOCNIVEAU\ \textbullet\ Chapter \DOCLECON}{white}{popink}}
\lfoot{\popxbold\fontsize{7.6}{7.6}\selectfont\color{popmuted}\MakeUppercase{Signed course OnBuch+}\ \textbullet\ {\popsemi\DOCTITRE}}
\rfoot{\tikz[baseline=-0.6ex]\node[rounded corners=8.5pt,fill=popink,minimum height=17pt,minimum width=17pt,inner xsep=5pt,inner sep=0]{\popxbold\fontsize{8}{8}\selectfont\color{popcream}\thepage\,/\,\pageref*{LastPage}};}

% ============================================================
% Titres
% ============================================================
\newcommand{\chaptitre}[2]{%
  \begin{tcolorbox}[enhanced,colback=poporange,colframe=popink,boxrule=2.6pt,arc=11pt,
      drop shadow=popink,left=15pt,right=15pt,top=12pt,bottom=11pt,before skip=3pt,after skip=18pt]
    \poppill{#2}{popink}{popcream}\par\vspace{9pt}
    {\raggedright\hyphenpenalty=10000\exhyphenpenalty=10000\popdisplay\fontsize{22}{24}\selectfont\color{popcream}\MakeUppercase{#1}\par}
  \end{tcolorbox}
}
% Section numérotée : pastille orange avec le numéro + titre Archivo
\newcounter{popsec}
\newcommand{\coursec}[1]{%
  \stepcounter{popsec}\setcounter{popsub}{0}%
  \par\vspace{16pt}\noindent
  \begin{minipage}{\linewidth}
  \tikz[baseline=-0.7ex]\node[draw=popink,line width=1.6pt,fill=poporange,rounded corners=4pt,minimum size=22pt,inner sep=0]{\popdisplay\fontsize{12}{12}\selectfont\color{popcream}\thepopsec};%
  \hspace{8pt}{\popdisplay\fontsize{15}{17}\selectfont\color{popink}\MakeUppercase{#1}}\par
  \vspace{4pt}\noindent{\color{poporange}\rule{36pt}{3.6pt}}
  \end{minipage}\par\nopagebreak\vspace{8pt}
}
\newcounter{popsub}
\newcommand{\courssub}[1]{%
  \stepcounter{popsub}%
  \par\vspace{9pt}\noindent{\popxbold\fontsize{11.5}{13}\selectfont\color{popink}{\color{poporange}\thepopsec.\thepopsub}\hspace{6pt}#1}\par\nopagebreak\vspace{4pt}
}

% ============================================================
% Boîtes pédagogiques — même recette : fond blanc, bordure épaisse colorée,
% ombre dure encre, badge plein. Titre optionnel [#1].
% ============================================================
\tcbset{popbase/.style={breakable,enhanced,colback=white,boxrule=2.3pt,arc=9pt,
  shadow={3pt}{-3pt}{0pt}{popink},left=14pt,right=14pt,top=11pt,bottom=11pt,before skip=11pt,after skip=15pt,
  fontupper=\popsemi\fontsize{10.6}{14.5}\selectfont\color{popink},
  fontlower=\popsemi\fontsize{10.6}{14.5}\selectfont\color{popink}}}
% Coche dessinée (le glyphe ✓ n'existe pas dans les polices du document)
\newcommand{\popcheck}[1]{\tikz[baseline=-0.5ex]\draw[#1,line width=1.6pt,line cap=round,line join=round](-0.13,0.0)--(-0.04,-0.09)--(0.13,0.1);}
\providecommand{\coche}{\popcheck{popgreen}}
\providecommand{\resultat}[1]{\par\smallskip\noindent\fcolorbox{popgreen}{popgreenL}{\parbox{\dimexpr\linewidth-2\fboxsep-2\fboxrule\relax}{\textbf{Result.} #1}}\par\smallskip}
\newcommand{\popboxhead}[3]{\poppill{#1}{#2}{white}\ifx\relax#3\relax\else\hspace{6pt}{\popxbold\fontsize{10.6}{12}\selectfont\color{popink}#3}\fi\par\vspace{7pt}}

\newtcolorbox{aretenir}[1][]{popbase,colframe=popgreen,before upper={\popboxhead{Key points}{popgreen}{#1}}}
\newtcolorbox{exemplebox}[1][]{popbase,colframe=poporange,before upper={\popboxhead{Exemple}{poporange}{#1}}}
\newtcolorbox{definition}[1][]{popbase,colframe=popblue,before upper={\popboxhead{Definition}{popblue}{#1}}}
\newtcolorbox{propriete}[1][]{popbase,colframe=poppurple,before upper={\popboxhead{Property}{poppurple}{#1}}}
\newtcolorbox{methode}[1][]{popbase,colframe=popgold,before upper={\popboxhead{Method}{popgold}{#1}}}
\newtcolorbox{attention}[1][]{popbase,colframe=poppink,before upper={\popboxhead{Warning}{poppink}{#1}}}
\newtcolorbox{experience}[1][]{popbase,colframe=popblue,colback=popblueL,before upper={\popboxhead{Experiment}{popblue}{#1}}}
% « Le savais-tu ? » : carte violette pleine (contexte, culture, Cameroun)
\newtcolorbox{savaistu}[1][]{popbase,colback=poppurple,colframe=popink,
  fontupper=\popsemi\fontsize{10.4}{14.2}\selectfont\color{white},
  before upper={\poppill{Did you know?}{white}{poppurple}\ifx\relax#1\relax\else\hspace{6pt}{\popxbold\color{white}#1}\fi\par\vspace{7pt}}}
% Exercice résolu : énoncé en haut, \tcblower puis la solution sur fond crème
\newcounter{popexo}\newcounter{popexr}
\newtcolorbox{exoresolu}[1][]{popbase,colframe=poporange,colbacklower=popcream,
  segmentation style={popink,line width=1.2pt,dashed},
  before upper={\stepcounter{popexr}\popboxhead{Worked example \thepopexr}{poporange}{#1}},
  before lower={\poppill{Solution}{popink}{popcream}\par\vspace{6pt}}}
% Exercice (non résolu en place) — la correction est dans la partie Corrigés
\newtcolorbox{exercice}[1][]{popbase,colframe=popink,
  before upper={\stepcounter{popexo}\popboxhead{Exercise \thepopexo}{popink}{#1}}}
% Corrigé
\newtcolorbox{corrige}[1][]{popbase,colframe=popgreen,colback=popgreenL,
  before upper={\popboxhead{Solution}{popgreen}{#1}}}
% Objectifs / prérequis / bilan
\newtcolorbox{objectifs}{popbase,colframe=popink,colback=poporangeL,
  before upper={\popboxhead{Chapter objectives}{popink}{}}}
\newtcolorbox{prerequis}{popbase,colframe=popink,colback=popblueL,
  before upper={\popboxhead{Prerequisites}{popblue}{}}}
\newtcolorbox{bilan}{popbase,colframe=popink,colback=white,boxrule=2.8pt,
  before upper={\popboxhead{Fiche bilan}{poporange}{L'essentiel à connaître pour l'examen}}}
% Figure encadrée avec légende
\newtcolorbox{popfigure}[1][]{enhanced,breakable=false,colback=white,colframe=popline,boxrule=1.4pt,arc=8pt,
  left=10pt,right=10pt,top=10pt,bottom=8pt,before skip=10pt,after skip=12pt,halign=center,
  lower separated=false,halign lower=center,
  fontlower=\popsemi\fontsize{9}{11.5}\selectfont\color{popmuted},
  #1}
\newcommand{\legende}[1]{\par\vspace{6pt}{\popsemi\fontsize{9}{11.5}\selectfont\color{popmuted}#1}}

% ============================================================
% Signature OnBuch+
% ============================================================
\newcommand{\poplogoinline}[1]{{\poplogo\fontsize{#1}{#1}\selectfont\color{popink}OnBuch\color{poporange}+}}
\newcommand{\signatureonbuch}{%
  \begin{tcolorbox}[enhanced,colback=popink,colframe=popink,boxrule=2.2pt,arc=10pt,
      drop shadow=poporange,left=14pt,right=14pt,top=10pt,bottom=10pt,before skip=0pt,after skip=0pt]
    \tikz[baseline=-0.7ex]\node[circle,fill=popgreen,draw=popcream,line width=1.3pt,minimum size=22pt,inner sep=0]{\popcheck{white}};%
    \hspace{9pt}\begin{minipage}[c]{0.62\linewidth}
      {\popxbold\fontsize{12}{13}\selectfont\color{popcream}Signed\ \textbullet\ {\poplogo OnBuch\color{poporange}+}}\par\vspace{2pt}
      {\raggedright\popsemi\fontsize{8}{10.5}\selectfont\color{popline}\MakeUppercase{OnBuch+ teaching team}\\\MakeUppercase{Follows the official MINESEC syllabus}\par}
    \end{minipage}\hfill
    {\poplogo\fontsize{22}{22}\selectfont\color{popcream}OnBuch\color{poporange}+}
  \end{tcolorbox}%
}

% ============================================================
% Page de garde
% #1 titre · #2 matière · #3 niveau · #4 module · #5 n° leçon · #6 durée ·
% #7 description / objectifs (texte ou itemize)
% ============================================================
\newcommand{\pagegarde}[7]{%
  \thispagestyle{empty}
  \newgeometry{a4paper,margin=1.7cm,top=1.6cm,bottom=1.4cm}%
  \begin{tikzpicture}[remember picture,overlay]
    \fill[poporange!18] ([xshift=-3.2cm,yshift=-3.6cm]current page.north east) circle (4.6cm);
    \fill[poppurple!14] ([xshift=2.4cm,yshift=7.6cm]current page.south west) circle (3.8cm);
  \end{tikzpicture}%
  {\poplogo\fontsize{30}{30}\selectfont\color{popink}OnBuch\color{poporange}+}\hfill
  \raisebox{0.6ex}{\poppill{Signed course \textbullet\ Premium}{popink}{popcream}}\par
  \vspace{1pt}\popsquiggle{poppurple}\par
  \vspace{8pt}
  \begin{tcolorbox}[enhanced,colback=poporange,colframe=popink,boxrule=2.8pt,arc=13pt,
      drop shadow=popink,left=18pt,right=18pt,top=12pt,bottom=13pt,after skip=10pt]
    \poppill{#2 \textbullet\ #3}{popink}{popcream}\hspace{5pt}\poppill{#4}{white}{popink}\par\vspace{8pt}
    {\popdisplay\fontsize{14}{14}\selectfont\color{popink}CHAPTER #5}\par\vspace{4pt}
    {\raggedright\hyphenpenalty=10000\exhyphenpenalty=10000\popdisplay\fontsize{25}{27}\selectfont\color{popcream}\MakeUppercase{#1}\par}
    \vspace{8pt}
    {\popxbold\fontsize{9.5}{9.5}\selectfont\color{popink}\MakeUppercase{Suggested time: #6}}
  \end{tcolorbox}
  \begin{tcolorbox}[enhanced,colback=white,colframe=popink,boxrule=2.2pt,arc=10pt,
      drop shadow=popink,left=16pt,right=16pt,top=10pt,bottom=10pt,after skip=8pt]
    \poppill{In this chapter}{poppurple}{white}\par\vspace{5pt}
    \popsemi\fontsize{9.3}{12.6}\selectfont\color{popink}#7
  \end{tcolorbox}
  \noindent\begin{minipage}[t]{0.487\linewidth}
    \begin{tcolorbox}[enhanced,colback=popgreen,colframe=popink,boxrule=2.2pt,arc=10pt,
        drop shadow=popink,left=11pt,right=11pt,top=10pt,bottom=10pt,height=3.1cm,valign=center]
      \tcbox[colback=white,colframe=popink,boxrule=1.3pt,arc=5pt,boxsep=0pt,nobeforeafter,
        left=3pt,right=3pt,top=3pt,bottom=3pt,valign=center]{\qrcode[height=1.9cm]{https://play.google.com/store/apps/details?id=com.luvvix.onbuch&hl=en}}%
      \hspace{8pt}\begin{minipage}[c]{0.5\linewidth}
        {\popdisplay\fontsize{11}{12.5}\selectfont\color{white}\MakeUppercase{Download OnBuch}}\par\vspace{4pt}
        {\raggedright\popsemi\fontsize{8.4}{11}\selectfont\color{white}Free \textbullet\ Google Play\\AI tutor, past papers, results\par}
      \end{minipage}
    \end{tcolorbox}
  \end{minipage}\hfill
  \begin{minipage}[t]{0.487\linewidth}
    \begin{tcolorbox}[enhanced,colback=poppurple,colframe=popink,boxrule=2.2pt,arc=10pt,
        drop shadow=popink,left=11pt,right=11pt,top=10pt,bottom=10pt,height=3.1cm,valign=center]
      \tikz[baseline=-0.7ex]\node[circle,fill=white,draw=popink,line width=1.3pt,minimum size=1.6cm,inner sep=0]{\popdisplay\fontsize{24}{24}\selectfont\color{poppurple}+};%
      \hspace{8pt}\begin{minipage}[c]{0.6\linewidth}
        {\popdisplay\fontsize{11}{12.5}\selectfont\color{white}\MakeUppercase{Go OnBuch+}}\par\vspace{4pt}
        {\raggedright\popsemi\fontsize{8.4}{11}\selectfont\color{white}All signed courses, unlimited AI tutor, PDF sheets — OnBuch+ tab in the app\par}
      \end{minipage}
    \end{tcolorbox}
  \end{minipage}
  \vspace{8pt}
  \popcontenu
  \vfill
  \signatureonbuch
  \restoregeometry
  \newpage
}

% Rangée « Ce cours contient » (page de garde)
\newcommand{\poptile}[3]{% #1 couleur #2 gros texte #3 légende
  \begin{tcolorbox}[enhanced,colback=#1,colframe=popink,boxrule=2pt,arc=9pt,shadow={2.5pt}{-2.5pt}{0pt}{popink},
      left=6pt,right=6pt,top=9pt,bottom=9pt,halign=center,valign=center,height=1.75cm,width=\linewidth,nobeforeafter]
    {\popdisplay\fontsize{12.5}{13}\selectfont\color{white}\MakeUppercase{#2}}\par\vspace{3pt}
    {\centering\popsemi\fontsize{7.8}{9.5}\selectfont\color{white}#3\par}
  \end{tcolorbox}}
\newcommand{\popcontenu}{%
  \par\noindent{\popxboldls\fontsize{7.5}{7.5}\selectfont\color{popmuted}THIS SIGNED COURSE CONTAINS}\par\vspace{7pt}
  \noindent\begin{minipage}[t]{0.235\linewidth}\poptile{popblue}{Course}{complete, illustrated, step by step}\end{minipage}\hfill
  \begin{minipage}[t]{0.235\linewidth}\poptile{poporange}{Methods}{and worked examples}\end{minipage}\hfill
  \begin{minipage}[t]{0.235\linewidth}\poptile{poppink}{Exercises}{exam style + solutions}\end{minipage}\hfill
  \begin{minipage}[t]{0.235\linewidth}\poptile{popgreen}{Summary sheet}{the essentials to remember}\end{minipage}\par}

% Sommaire
\newtcolorbox{sommaire}{enhanced,colback=white,colframe=popink,boxrule=2.2pt,arc=10pt,
  drop shadow=popink,left=18pt,right=18pt,top=13pt,bottom=13pt,before skip=6pt,after skip=20pt,
  before upper={\poppill{Contents}{poppurple}{white}\par\vspace{9pt}\popsemi\fontsize{10.6}{14.5}\selectfont\color{popink}}}

% Signature de fin de cours — poussée tout en bas de la dernière page
\newcommand{\findecours}{%
  \par\vfill\nopagebreak
  \begin{center}\popsquiggle{poporange}\end{center}\vspace{4pt}
  \signatureonbuch
}
\AtBeginDocument{\def\llbracket{\mathopen{[\mkern-3.5mu[}}\def\rrbracket{\mathclose{]\mkern-3.5mu]}}} % intervalles d'entiers (glyphe ⟦ absent de la police)
% Glyphes absents de Fira Math : redéfinis à partir de glyphes disponibles
\AtBeginDocument{%
  \def\setminus{\mathbin{\backslash}}\let\smallsetminus\setminus
  \def\vdots{\mathord{\vbox{\baselineskip3.2pt\lineskiplimit0pt\kern4pt\hbox{.}\hbox{.}\hbox{.}}}}%
  \def\ddots{\mathinner{\mkern1mu\raise6.5pt\hbox{.}\mkern2mu\raise3.6pt\hbox{.}\mkern2mu\raise0.7pt\hbox{.}\mkern1mu}}%
  \def\triangle{\mathord{\scalebox{0.9}{$\symup{\Delta}$}}}%
  \def\nabla{\mathord{\raisebox{\depth}{\scalebox{1}[-1]{$\symup{\Delta}$}}}}%
  \def\checkmark{\popcheck{popdark}}%
  \def\star{\mathbin{\ast}}\def\bigstar{\mathord{\ast}}%
  \def\diamond{\mathbin{\scalebox{0.6}{$\Diamond$}}}%
  \def\bigcup{\mathop{\mathchoice{\raisebox{-0.3em}{\scalebox{1.7}{$\cup$}}}{\scalebox{1.25}{$\cup$}}{\cup}{\cup}}\displaylimits}%
  \def\bigcap{\mathop{\mathchoice{\raisebox{-0.3em}{\scalebox{1.7}{$\cap$}}}{\scalebox{1.25}{$\cap$}}{\cap}{\cap}}\displaylimits}%
  \def\lesseqqgtr{\mathrel{\lessgtr}}\def\bigoplus{\mathop{\oplus}\displaylimits}%
}
\providecommand{\ssi}{if and only if} % abréviation fréquente
\AtBeginDocument{\providecommand{\wideparen}{\overparen}} % arc de cercle

%% FIGFIX-BEGIN v5 — correctifs globaux des figures (docs/onbuch-generation/tools/figfix)
\usepackage{adjustbox}\usepackage{etoolbox}\tcbuselibrary{raster}
% toute figure TikZ/pgfplots plus large que la ligne est réduite à la largeur disponible
\BeforeBeginEnvironment{tikzpicture}{\begin{adjustbox}{max width=\linewidth}}
\AfterEndEnvironment{tikzpicture}{\end{adjustbox}}
% légendes pgfplots lisibles par-dessus les courbes
\pgfplotsset{every axis/.append style={legend style={fill=white,fill opacity=0.92,text opacity=1,draw=black!15,inner sep=3pt}}}
% étiquettes d'arêtes (syntaxe edge["..."]) avec fond blanc
\tikzset{every edge quotes/.append style={fill=white,inner sep=1.2pt}}
%% FIGFIX-END
\begin{document}

\pagegarde{Agriculture}{Biology}{Upper Sixth}{Upper Sixth — Biology}{8}{5 periods}{This chapter explores the scientific foundations of agriculture: the types and properties of soils, the nutritional requirements of plants and animals, and the techniques used to cultivate crops and raise livestock. It equips learners with the knowledge needed to improve agricultural productivity and to tackle GCE Advanced Level questions on agriculture with confidence.
\vspace{4pt}
\begin{multicols}{2}\small
\begin{itemize}
\item By the end of this chapter I can describe the main types of soils and relate their properties to agricultural use.
\item By the end of this chapter I can explain soil formation, soil profile and the factors affecting soil fertility.
\item By the end of this chapter I can state the mineral elements required by plants and describe deficiency symptoms and methods of correcting them.
\item By the end of this chapter I can distinguish between food and non-food crops and describe the cultivation techniques of selected Cameroonian crops.
\item By the end of this chapter I can describe the nutritional requirements of farm animals and the classes of feedstuffs.
\item By the end of this chapter I can explain the principles of ration formulation and feeding regimes for livestock.
\end{itemize}\end{multicols}}

\begin{sommaire}
\begin{multicols}{2}
\begin{enumerate}
\item Soils: Types, Formation and Fertility
\item Plant Nutritional Requirements
\item Cultivation Techniques of Food and Non-Food Crops
\item Animal Nutrition and Animal Production
\item Integration activity
\item Methods and worked examples
\item Exercises
\item Solutions to the exercises
\item Summary sheet
\end{enumerate}
\end{multicols}
\end{sommaire}

\begin{objectifs}
\begin{itemize}
\item By the end of this chapter I can describe the main types of soils and relate their properties to agricultural use.
\item By the end of this chapter I can explain soil formation, soil profile and the factors affecting soil fertility.
\item By the end of this chapter I can state the mineral elements required by plants and describe deficiency symptoms and methods of correcting them.
\item By the end of this chapter I can distinguish between food and non-food crops and describe the cultivation techniques of selected Cameroonian crops.
\item By the end of this chapter I can describe the nutritional requirements of farm animals and the classes of feedstuffs.
\item By the end of this chapter I can explain the principles of ration formulation and feeding regimes for livestock.
\item By the end of this chapter I can describe the systems and techniques of animal production, including housing, breeding, health and record keeping.
\item By the end of this chapter I can apply agricultural knowledge to solve practical farm problems and answer GCE-style structured and essay questions.
\end{itemize}
\end{objectifs}

\begin{prerequis}
\begin{itemize}
\item Plant nutrition and photosynthesis (Lower Sixth): autotrophic nutrition, mineral absorption by roots
\item Basic ecology: ecosystems, nutrient cycling (nitrogen and carbon cycles)
\item Chemical composition of living matter: carbohydrates, proteins, lipids, vitamins and minerals
\item Digestion in mammals and the structure of the alimentary canal (including ruminants)
\item Reproduction in plants and animals (Lower Sixth)
\end{itemize}
\end{prerequis}

\newpage

\coursec{Soils: Types, Formation and Fertility}

In the village of Bafut, two farmers plant maize on neighbouring plots during the same week; at harvest, one fills ten bags while the other barely fills three, despite using the same seed variety. Meanwhile, a poultry keeper in Bamenda loses half her broilers before week six, while her neighbour's birds reach market weight on schedule. What hidden factors in the soil, the feed and the farming methods explain such dramatic differences? This chapter takes you into the science of agriculture --- understanding soils, feeding crops and animals correctly, and mastering production techniques --- so that farming in Cameroon becomes a predictable, profitable science rather than a gamble. We begin at the very foundation of it all: the soil beneath our feet.

\courssub{What is Soil and Why Does it Matter?}

Walk into any farm around Buea, Ngaoundéré or Bafoussam and scoop up a handful of earth. What you hold is not simply ``dirt''. It is a living, complex mixture of rock fragments, decaying organisms, water, air and countless microbes. Every crop a farmer plants --- maize, cassava, cocoa, plantain --- depends on this thin layer of material for anchorage, water and mineral nutrients.

\begin{definition}[Soil]
\cle{Soil} is the uppermost layer of the Earth's crust, formed by the weathering of rocks and the decomposition of organic matter, consisting of mineral particles, organic matter (humus), water, air and living organisms, and capable of supporting plant growth.
\end{definition}

Soil is the \cle{basis of agriculture} because it performs four essential functions for the farmer:
\begin{itemize}
\item \textbf{Anchorage:} roots grip soil particles, holding the plant upright against wind and rain.
\item \textbf{Water supply:} soil stores rain and irrigation water and releases it to roots.
\item \textbf{Mineral nutrition:} soil supplies essential elements such as nitrogen, phosphorus and potassium.
\item \textbf{Habitat:} soil hosts organisms (earthworms, bacteria, fungi) that recycle nutrients and maintain fertility.
\end{itemize}

\begin{savaistu}[Soil is a non-renewable resource on a human timescale]
It can take several hundred years for just a few centimetres of topsoil to form naturally, yet a single season of careless farming on a steep slope in the North West Region can wash it away. This is why soil conservation is a matter of national food security.
\end{savaistu}

\courssub{Soil Formation (Pedogenesis)}

Soil does not appear from nowhere. It forms slowly from \cle{parent rock} through a process called \cle{pedogenesis}, driven by weathering and the action of living organisms.

\textbf{Weathering of parent rock}\par
Weathering is the breakdown of solid rock into smaller and smaller particles. Three types act together:

\begin{itemize}
\item \textbf{Physical (mechanical) weathering:} rock is broken into fragments without any change in chemical composition. In Cameroon this includes the alternate heating and cooling of rocks under the hot sun of the Far North (thermal expansion and contraction), and the wedging action of water and growing roots in rock cracks.
\item \textbf{Chemical weathering:} the minerals in rock are chemically changed. Rainwater, slightly acidified by dissolved carbon dioxide, dissolves limestone (solution); oxygen oxidises iron minerals, giving tropical soils their red colour; water reacts with feldspar minerals to form clay (hydrolysis). Chemical weathering is fastest in hot, wet climates --- which is why tropical soils in the South of Cameroon are deeply weathered.
\item \textbf{Biological weathering:} living organisms break down rock. Plant roots exert pressure in cracks; burrowing animals mix and fragment material; lichens and bacteria secrete acids that dissolve minerals.
\end{itemize}

\textbf{Factors controlling soil formation}\par
The nature of the soil produced depends on five factors, easily remembered as \textbf{CLORPT}:
\begin{itemize}
\item \textbf{Cl}imate: rainfall and temperature control the speed of weathering and leaching.
\item \textbf{O}rganisms: vegetation, microbes and animals add organic matter and mix the soil.
\item \textbf{R}elief (topography): steep slopes lose soil by erosion; valleys accumulate deep soils.
\item \textbf{P}arent material: the original rock determines mineral content (e.g.\ basalt gives fertile volcanic soils).
\item \textbf{T}ime: the longer the process, the deeper and more developed the soil.
\end{itemize}

\courssub{The Soil Profile}

If you dig a pit about one metre deep on undisturbed land, the wall of the pit reveals distinct horizontal layers called \cle{horizons}. The vertical sequence of horizons is called the \cle{soil profile}.

\begin{popfigure}
\begin{tikzpicture}[scale=1]
% Bedrock
\fill[gray!60] (0,0) rectangle (8,1.2);
\draw[gray, thick] (0.4,0.3) -- (1.4,0.9); \draw[gray, thick] (2.2,0.2) -- (3.2,1.0);
\draw[gray, thick] (4.0,0.3) -- (5.0,0.9); \draw[gray, thick] (5.8,0.2) -- (6.8,1.0);
% C horizon
\fill[popgold!50] (0,1.2) rectangle (8,2.6);
\foreach \x in {0.5,1.5,2.5,3.5,4.5,5.5,6.5,7.5} \fill[gray!70] (\x,1.7) circle (0.08);
\foreach \x in {1.0,2.0,3.0,4.0,5.0,6.0,7.0} \fill[gray!70] (\x,2.2) circle (0.06);
% B horizon
\fill[poporange!45] (0,2.6) rectangle (8,4.4);
% A horizon
\fill[popdark!55] (0,4.4) rectangle (8,5.8);
% O horizon
\fill[popgreen!45] (0,5.8) rectangle (8,6.4);
\foreach \x in {0.4,1.2,2.0,2.8,3.6,4.4,5.2,6.0,6.8,7.6} \draw[popgreen, thick] (\x,6.4) -- (\x,6.8);
% Labels
\node[right, align=left] at (8.2,6.1) {\textbf{O horizon}: organic litter};
\node[right, align=left] at (8.2,5.1) {\textbf{A horizon}: topsoil, humus-rich};
\node[right, align=left] at (8.2,3.5) {\textbf{B horizon}: subsoil, minerals accumulate};
\node[right, align=left] at (8.2,1.9) {\textbf{C horizon}: weathered parent rock};
\node[right, align=left] at (8.2,0.6) {\textbf{Bedrock}: unweathered rock};
\draw[->, thick] (8.15,6.1) -- (7.95,6.1);
\draw[->, thick] (8.15,5.1) -- (7.95,5.1);
\draw[->, thick] (8.15,3.5) -- (7.95,3.5);
\draw[->, thick] (8.15,1.9) -- (7.95,1.9);
\draw[->, thick] (8.15,0.6) -- (7.95,0.6);
\end{tikzpicture}
\legende{A typical soil profile showing the horizons O, A, B, C and the bedrock.}
\end{popfigure}

\begin{center}
\begin{tabularx}{\linewidth}{|l|X|}
\hline
\rowcolor{popblueL} \textbf{Horizon} & \textbf{Characteristics} \\
\hline
O (organic) & Surface layer of dead leaves and decomposing remains; source of humus. \\
\hline
A (topsoil) & Dark, rich in humus mixed with mineral particles; most biological activity; most roots found here; loses minerals by leaching (eluviation). \\
\hline
B (subsoil) & Lighter in colour; accumulates clay, iron and aluminium oxides washed down from A (illuviation); fewer organisms. \\
\hline
C & Partly weathered parent rock fragments; little organic matter. \\
\hline
Bedrock (R) & Solid, unweathered rock from which the soil formed. \\
\hline
\end{tabularx}
\end{center}

\courssub{Components of Soil}

A fertile soil is a four-part mixture. By volume, an ideal mineral soil contains roughly 45\% mineral particles, 5\% organic matter, 25\% water and 25\% air.

\begin{itemize}
\item \textbf{Mineral particles:} weathered rock fragments (sand, silt, clay) supplying mineral nutrients.
\item \textbf{Organic matter (humus):} dark, decomposed plant and animal remains. \cle{Humus} improves structure, holds water and releases nutrients slowly.
\item \textbf{Soil water:} a dilute solution of mineral salts; the medium through which roots absorb nutrients.
\item \textbf{Soil air:} occupies the pore spaces; supplies oxygen for root respiration and for aerobic microbes.
\item \textbf{Soil organisms:} earthworms, termites, bacteria and fungi decompose organic matter, fix nitrogen and aerate the soil.
\end{itemize}

\courssub{Soil Texture}

\begin{definition}[Soil texture]
\cle{Soil texture} is the relative proportion by mass of the different-sized mineral particles --- \textbf{sand} (0.05--2 mm), \textbf{silt} (0.002--0.05 mm) and \textbf{clay} (less than 0.002 mm) --- present in a soil.
\end{definition}

Texture controls almost everything a farmer cares about: how fast water drains, how much water is stored, how easily roots penetrate, and how well nutrients are held. The \cle{soil texture triangle} is used to classify any soil once the percentages of sand, silt and clay are known.

\begin{popfigure}
\begin{tikzpicture}[scale=0.09]
% Triangle outline
\draw[thick] (0,0) -- (100,0) -- (50,86.6) -- cycle;
% Grid lines (selected)
\draw[gray!50] (20,0) -- (60,69.28);
\draw[gray!50] (40,0) -- (70,51.96);
\draw[gray!50] (60,0) -- (80,34.64);
\draw[gray!50] (80,0) -- (90,17.32);
\draw[gray!50] (10,17.32) -- (90,17.32);
\draw[gray!50] (20,34.64) -- (80,34.64);
\draw[gray!50] (30,51.96) -- (70,51.96);
\draw[gray!50] (40,69.28) -- (60,69.28);
\draw[gray!50] (10,17.32) -- (20,0);
\draw[gray!50] (30,51.96) -- (40,0);
\draw[gray!50] (20,34.64) -- (60,0);
% Loam zone (approx 40% sand, 40% silt, 20% clay region)
\fill[popgreen, opacity=0.5] (40,17.32) -- (52.5,17.32) -- (45,30.31) -- (37.5,30.31) -- cycle;
\node[font=\small] at (44,23.5) {\textbf{LOAM}};
% Vertex labels
\node[below] at (0,0) {100\% sand};
\node[below] at (100,0) {100\% silt};
\node[above] at (50,86.6) {100\% clay};
\node[font=\small] at (50,-8) {sand \% $\longleftarrow$ \quad \quad $\longrightarrow$ silt \%};
\end{tikzpicture}
\legende{The soil texture triangle. Reading the percentages of sand, silt and clay locates the soil class; the shaded central zone is loam.}
\end{popfigure}

\begin{methode}[Determining soil texture by the sedimentation (jar) method]
\begin{enumerate}
\item Collect a sample of dry topsoil, remove stones and crush lumps.
\item Half-fill a tall transparent jar (or measuring cylinder) with the soil, then add water and a spoonful of salt or dispersing agent.
\item Shake vigorously for about two minutes to separate all particles.
\item Stand the jar undisturbed. Sand settles within about one minute, silt within a few hours, and clay may take a day or more, forming distinct layers.
\item Measure the thickness of each layer and calculate its percentage of the total settled depth.
\item Plot the three percentages on the texture triangle to name the soil class.
\end{enumerate}
\end{methode}

Texture can also be estimated \textbf{by feel}: a moist sandy soil feels gritty and will not form a ball; a clay soil feels sticky and can be rolled into a long, flexible ribbon; a loam forms a ball that breaks easily.

\courssub{Main Soil Types and their Distribution in Cameroon}

\begin{center}
\begin{tabularx}{\linewidth}{|l|X|X|}
\hline
\rowcolor{popblueL} \textbf{Soil type} & \textbf{Properties} & \textbf{Occurrence in Cameroon} \\
\hline
Sandy & Large particles; drains fast, warms quickly, low nutrient and water retention & Parts of the Far North (sahelian zone) \\
\hline
Clay & Tiny particles; holds water and nutrients well but poorly aerated, hard when dry, waterlogged when wet & Low-lying plains, e.g.\ Logone floodplain \\
\hline
Loamy & Balanced sand--silt--clay mixture with humus; fertile and easy to work & Western highlands, many farmed areas \\
\hline
Lateritic & Red, rich in iron and aluminium oxides; heavily leached, low fertility unless manured & Humid South, Centre and East \\
\hline
Alluvial & Deposited by rivers; deep, fertile, regularly renewed by floods & Benue and Logone valleys, coastal plains \\
\hline
Volcanic & From weathered basalt; dark, deep, very fertile & Mount Cameroon and the Bamenda highlands \\
\hline
Peaty & Very rich in organic matter; acidic and often waterlogged & Waterlogged depressions and swamps \\
\hline
\end{tabularx}
\end{center}

\begin{propriete}[Loamy soil: the ideal agricultural soil]
A \cle{loam} (roughly 40\% sand, 40\% silt, 20\% clay, plus humus) is the best general-purpose agricultural soil because it combines \textbf{good drainage} (from sand), \textbf{good water and nutrient retention} (from clay and humus) and \textbf{good aeration and easy root penetration} (from its crumb structure). Proof is not examinable; you must simply be able to justify the superiority of loam in terms of drainage, aeration and nutrient retention.
\end{propriete}

\courssub{Soil Structure}

\begin{definition}[Soil structure]
\cle{Soil structure} is the way in which soil particles are grouped and bound together into larger units called \textbf{aggregates} (or peds), and the arrangement of the pore spaces between them.
\end{definition}

\begin{attention}[Texture $\neq$ structure]
A very common GCE mistake is to confuse the two. \textbf{Texture} is about the \emph{sizes} of individual particles (sand, silt, clay) and cannot be changed by cultivation. \textbf{Structure} is about how those particles are \emph{arranged into aggregates}; it can be improved by adding humus or destroyed by cultivating wet soil. A clay soil (texture) can have an excellent crumb structure.
\end{attention}

The main structural types are:
\begin{itemize}
\item \textbf{Crumb:} small, rounded, porous aggregates; ideal for seedbeds --- excellent aeration, drainage and root growth.
\item \textbf{Blocky:} larger, roughly cube-shaped aggregates; moderately good, common in subsoils.
\item \textbf{Platy:} thin horizontal plates; impedes drainage and root penetration.
\item \textbf{Prismatic:} vertical columns; restricts root spread and water movement.
\end{itemize}

\courssub{Physical and Chemical Properties of Soil}

\textbf{Porosity} is the percentage of soil volume occupied by pore spaces; it determines how much air and water the soil can hold. \textbf{Water-holding capacity} is the amount of water a soil retains against gravity --- highest in clay and humus-rich soils, lowest in sands. \textbf{Capillarity} is the upward movement of water through fine pores, which can bring stored water up to roots but also causes water loss by evaporation from bare soil.

\begin{popfigure}
\begin{tikzpicture}
\begin{axis}[
ybar, bar width=0.7cm,
width=11cm, height=7cm,
ymin=0, ymax=100,
ylabel={Relative value (\%)},
symbolic x coords={Sandy, Loamy, Clay},
xtick=data,
legend style={at={(0.5,1.02)}, anchor=south, legend columns=2},
nodes near coords, every node near coord/.append style={font=\small}
]
\addplot[fill=popblue] coordinates {(Sandy,25) (Loamy,60) (Clay,85)};
\addplot[fill=poporange] coordinates {(Sandy,90) (Loamy,60) (Clay,25)};
\legend{Water-holding capacity, Drainage rate}
\end{axis}
\end{tikzpicture}
\legende{Comparison of water-holding capacity and drainage rate in sandy, loamy and clay soils (relative values). Loam offers the best balance.}
\end{popfigure}

\textbf{Soil pH} measures acidity or alkalinity. Most Cameroonian crops grow best between pH 5.5 and 7.0. Heavy tropical rainfall leaches away basic cations (Ca$^{2+}$, Mg$^{2+}$, K$^+$), leaving soils acidic. Acidity is corrected by \cle{liming} --- adding agricultural lime (calcium carbonate), which neutralises excess H$^+$ ions and also supplies calcium.

\textbf{Cation exchange capacity (CEC)} is the soil's ability to hold positively charged nutrient ions (K$^+$, Ca$^{2+}$, Mg$^{2+}$, NH$_4^+$) on the surfaces of clay and humus particles, preventing them from being leached while keeping them available to roots. Clay and humus have high CEC; sand has very low CEC --- one more reason sandy soils are infertile.

\begin{exoresolu}[Choosing a crop for a given soil]
A farmer near Maroua has a light, gritty soil that dries out two days after rain and gives low yields of maize despite heavy fertiliser use. (a) Identify the soil texture. (b) Explain the poor response to fertiliser. (c) Suggest one suitable crop and one soil improvement measure.
\tcblower
\textbf{(a)} Rapid drainage and gritty feel indicate a \textbf{sandy soil}.\\
\textbf{(b)} Sand has very low cation exchange capacity and low water-holding capacity. Soluble fertiliser ions (e.g.\ K$^+$, NO$_3^-$) are rapidly \textbf{leached} below the rooting zone by rain, so the plants cannot use them and the fertiliser is wasted.\\
\textbf{(c)} A drought-tolerant crop such as \textbf{groundnuts, millet or cowpea} suits sandy soil. The soil can be improved by adding \textbf{organic manure or compost}, which raises water-holding capacity and CEC, and by applying fertiliser in small, frequent doses.
\end{exoresolu}

\courssub{Soil Fertility and its Maintenance}

\begin{definition}[Soil fertility]
\cle{Soil fertility} is the ability of a soil to supply plants with adequate water, air and mineral nutrients in the right proportions for healthy growth and good yields.
\end{definition}

Fertility is not permanent: continuous cropping, leaching and erosion exhaust it. Good farmers maintain fertility by:

\begin{itemize}
\item \textbf{Crop rotation:} growing different crops in sequence on the same plot (e.g.\ maize $\rightarrow$ groundnuts $\rightarrow$ cassava). Legumes restore nitrogen through their root-nodule bacteria, and pests specific to one crop are starved out.
\item \textbf{Fallowing:} resting the land under natural vegetation for some years so that humus and nutrients accumulate again --- the principle of traditional shifting cultivation.
\item \textbf{Organic manures:} farmyard manure, compost and green manure add humus, improve structure and release nutrients slowly.
\item \textbf{Inorganic (chemical) fertilisers:} supply concentrated nutrients quickly (e.g.\ NPK 20-10-10 for maize), but add no humus and can acidify soil if misused.
\item \textbf{Cover cropping and mulching:} covering the soil with living plants (e.g.\ \emph{Pueraria}) or dead material (grass, plantain leaves) protects against erosion, conserves moisture, suppresses weeds and adds organic matter.
\end{itemize}

\textbf{Soil erosion} is the removal of topsoil by running water or wind --- the fastest way to destroy fertility, and a serious problem on the steep cultivated slopes of the North West and South West Regions. \textbf{Conservation methods} include:
\begin{itemize}
\item \textbf{Contour ploughing:} ploughing along lines of equal height so each furrow checks run-off.
\item \textbf{Terracing:} cutting slopes into level steps to slow water flow (seen on the hills around Dschang).
\item \textbf{Windbreaks:} rows of trees or tall crops that reduce wind speed and protect soil in the dry North.
\item \textbf{Strip cropping, cover crops and mulching,} which keep the soil permanently covered.
\end{itemize}

\begin{aretenir}[Key points]
\begin{itemize}
\item Soil forms by weathering of parent rock plus organic activity, controlled by climate, organisms, relief, parent material and time.
\item The profile shows horizons O, A, B, C above bedrock; the A horizon is the fertile topsoil.
\item Texture = particle size proportions; structure = arrangement of aggregates. Do not confuse them.
\item Loam is the ideal soil: balanced drainage, aeration and nutrient retention.
\item Fertility is maintained by rotation, fallowing, manures, mulching and erosion control.
\end{itemize}
\end{aretenir}

\begin{exemplebox}[Exam tip]
GCE questions frequently ask you to \textbf{relate soil properties to crop choice} (e.g.\ ``Explain why rice is grown in the clay soils of the Logone plain while groundnuts thrive in the sandy soils of the Far North''). Always structure your justification around the same three pillars: \textbf{drainage}, \textbf{aeration} and \textbf{nutrient content} (including CEC and pH). A property named without its consequence for the crop earns little credit.
\end{exemplebox}

\begin{exercice}[Test your understanding]
\begin{enumerate}
\item Define soil texture and soil structure, giving one example of each.
\item A soil sample settles in a jar into layers of 6 cm sand, 3 cm silt and 1 cm clay. Calculate the percentage of each fraction and name the likely soil class.
\item Explain three ways in which adding farmyard manure improves a sandy soil.
\item A farmer on a steep slope near Bamenda loses topsoil every rainy season. Describe two conservation methods he could adopt, explaining how each works.
\end{enumerate}
\end{exercice}

\begin{corrige}[Exercise 1]
\textbf{1.} Texture: the proportions of sand, silt and clay particles in a soil (e.g.\ a loam with 40\% sand, 40\% silt, 20\% clay). Structure: the arrangement of particles into aggregates (e.g.\ crumb structure).\\
\textbf{2.} Total depth $= 6+3+1 = 10$ cm. Sand $= 60\%$, silt $= 30\%$, clay $= 10\%$. This is a \textbf{sandy loam} --- free-draining but reasonably workable.\\
\textbf{3.} Manure (i) increases water-holding capacity, (ii) raises cation exchange capacity so nutrients are retained, and (iii) improves crumb structure and aeration as it decomposes to humus.\\
\textbf{4.} \textbf{Contour ploughing:} furrows along contour lines trap run-off and let water soak in instead of flowing downhill. \textbf{Terracing:} level steps reduce slope length and water speed, so soil is deposited rather than carried away. (Windbreaks or mulching are also acceptable with correct explanations.)
\end{corrige}

\coursec{Plant Nutritional Requirements}

In the previous section we saw that a fertile soil is not merely a pile of particles: it is a living, chemically active medium whose texture, structure, organic matter content and pH determine how well it can support a crop. But \emph{why} does a maize plant in a red lateritic soil of the West Region turn pale and stunted, while the same variety thrives a few kilometres away on a well-manured plot? The answer lies in \textbf{plant nutrition}: the soil is the pantry, and the plant must draw from it a precise menu of chemical elements. This section studies that menu --- which elements a plant needs, in what quantities, what happens when one is missing, and how the farmer can correct shortages with fertilisers.

\courssub{Autotrophic Nutrition: The Basis of Crop Production}

Green plants are \cle{autotrophs}: they build their own organic matter from simple inorganic raw materials. The fundamental process is \textbf{photosynthesis}, summarised by:

\[ 6\ce{CO2} + 6\ce{H2O} \xrightarrow[\text{chlorophyll}]{\text{light}} \ce{C6H12O6} + 6\ce{O2} \]

Carbon dioxide provides the carbon skeleton, water provides hydrogen and electrons, and light provides the energy. The glucose produced is then converted into starch (the yam tuber, the cassava root, the maize grain we harvest), cellulose (stems, cotton fibre), proteins and oils.

\begin{attention}[What photosynthesis does NOT provide]
Photosynthesis supplies carbon, hydrogen and oxygen --- about 95\% of a plant's dry mass. The remaining 5\%, the \textbf{mineral fraction}, must come from the soil solution through the roots. A plant can photosynthesise perfectly and still die if a single essential mineral is missing. This is why soil fertility (Section 1) and mineral nutrition are inseparable.
\end{attention}

Even with all minerals present, the \textbf{yield} of a crop is limited by whichever factor is in shortest supply: light intensity, \ce{CO2} concentration, water availability, temperature, or mineral nutrients. This is the \cle{law of the minimum} (associated with the chemist Liebig): growth is controlled not by the total of resources available, but by the \emph{scarcest} resource --- just as the water level in a leaking barrel can rise no higher than the shortest stave. In most Cameroonian farms, after water, the limiting factors are mineral nutrients --- especially nitrogen, phosphorus and potassium.

\courssub{Essential Mineral Elements}

An element is \textbf{essential} if it satisfies three \cle{criteria of essentiality}:
\begin{enumerate}
\item The plant cannot complete its life cycle (germination $\rightarrow$ growth $\rightarrow$ flowering $\rightarrow$ seed production) without it.
\item Its role is \emph{specific} and cannot be replaced by any other element.
\item It is directly involved in plant metabolism (a constituent of a molecule or an enzyme activator), not merely improving growth indirectly.
\end{enumerate}

Essential elements are divided into two groups according to the quantities required.

\begin{definition}[Macronutrient and micronutrient]
A \cle{macronutrient} is an essential element required by the plant in relatively \textbf{large amounts} (typically more than 1 g per kg of dry matter). The mineral macronutrients absorbed from the soil are: \textbf{nitrogen (N), phosphorus (P), potassium (K), calcium (Ca), magnesium (Mg) and sulphur (S)}. (Carbon, hydrogen and oxygen, obtained from \ce{CO2} and \ce{H2O}, are also macronutrients.)

A \cle{micronutrient} (trace element) is an essential element required in \textbf{very small amounts} (milligrams per kg of dry matter): \textbf{iron (Fe), manganese (Mn), boron (B), zinc (Zn), copper (Cu), molybdenum (Mo) and chlorine (Cl)}. Despite the tiny quantities, a deficiency of any one of them is just as damaging as a deficiency of a macronutrient.
\end{definition}

\courssub{Functions and Deficiency Symptoms of the Major Elements}

Each element has precise roles, and its absence produces characteristic, diagnosable symptoms. A useful diagnostic clue is \textbf{element mobility} within the plant: mobile elements (N, P, K, Mg) are moved from old leaves to young growing points when scarce, so symptoms appear on \emph{older} leaves first; immobile elements (Ca, Fe) cannot be relocated, so symptoms appear on \emph{young} leaves first.

\begin{center}
\begin{tabularx}{\linewidth}{|l|X|X|}
\hline
\rowcolor{popblueL} \textbf{Element} & \textbf{Main functions} & \textbf{Deficiency symptoms} \\
\hline
Nitrogen (N) & Constituent of proteins, nucleic acids, chlorophyll & General yellowing (\cle{chlorosis}) of older leaves; stunted growth; thin, weak stems \\
\hline
Phosphorus (P) & ATP, nucleic acids, phospholipids; energy transfer; root development & Poor root growth; delayed maturity; dark green or \textbf{purplish} leaves (anthocyanin accumulation) \\
\hline
Potassium (K) & Enzyme activation; stomatal opening; osmoregulation; disease resistance & \textbf{Leaf scorch}: yellowing then browning of leaf \emph{margins} of older leaves; weak stems, lodging in cereals \\
\hline
Magnesium (Mg) & Central atom of the chlorophyll molecule & \textbf{Interveinal chlorosis} of older leaves (veins stay green, tissue between them yellows) \\
\hline
Calcium (Ca) & Cell wall cement (calcium pectate); membrane stability & Death of growing points and root tips; blossom-end rot in tomato; symptoms on young tissues \\
\hline
Sulphur (S) & Constituent of some amino acids (cysteine, methionine) & General chlorosis of \emph{young} leaves (resembles N deficiency but on young leaves) \\
\hline
Iron (Fe) & Chlorophyll synthesis (cofactor); electron transport & Interveinal chlorosis of \emph{young} leaves; common in alkaline/calcareous soils \\
\hline
\end{tabularx}
\end{center}

\begin{propriete}[NPK: the elements most often limiting yield]
Nitrogen, phosphorus and potassium are the three elements removed from the soil in the \textbf{largest quantities} by harvested crops, and they are therefore the ones most often deficient in cultivated soils. They form the basis of compound fertilisers labelled \textbf{NPK}, whose three numbers give the percentage by mass of N, \ce{P2O5} and \ce{K2O} respectively. For example, a bag marked \textbf{20-10-10} contains 20\% nitrogen, 10\% phosphorus pentoxide and 10\% potassium oxide by mass. (Knowledge of this property is examinable; no proof is required.)
\end{propriete}

\begin{popfigure}
\begin{center}
\begin{tikzpicture}[scale=1.0]
% Leaf outline
\draw[thick, popgreen, fill=popgreenL] (0,0) .. controls (1.6,1.4) and (3.4,1.6) .. (5,0)
   .. controls (3.4,-1.6) and (1.6,-1.4) .. (0,0) -- cycle;
% Midrib and veins
\draw[popgreen, thick] (0,0) -- (5,0);
\draw[popgreen] (1,0) -- (1.7,0.75); \draw[popgreen] (1,0) -- (1.7,-0.75);
\draw[popgreen] (2.2,0) -- (2.9,0.95); \draw[popgreen] (2.2,0) -- (2.9,-0.95);
\draw[popgreen] (3.4,0) -- (4.0,0.7); \draw[popgreen] (3.4,0) -- (4.0,-0.7);
% Petiole
\draw[popgreen, thick] (0,0) -- (-0.7,0);
% Marginal scorch (K)
\draw[very thick, poporange, dashed] (0.3,0.35) .. controls (1.8,1.25) and (3.4,1.35) .. (4.8,0.12);
% Labels
\node[align=center, font=\small, text=popdark] at (2.5,2.3) {\textbf{K deficiency:} marginal scorch\\ (yellow-brown edges, older leaves)};
\draw[->, poporange, thick] (2.5,1.95) -- (2.6,1.15);
\node[align=center, font=\small, text=popdark] at (6.3,0.6) {\textbf{Mg deficiency:} interveinal\\ chlorosis, veins stay green};
\draw[->, popdark, thick] (5.6,0.35) -- (3.1,0.45);
\node[align=center, font=\small, text=popdark] at (2.5,-2.3) {\textbf{N deficiency:} uniform pale\\ yellow of whole older leaf};
\draw[->, popgold, thick] (2.5,-1.95) -- (2.5,-0.75);
\end{tikzpicture}
\end{center}
\legende{Stylised leaf showing the three classic chlorosis patterns: uniform yellowing (N), interveinal yellowing with green veins (Mg), and marginal scorch (K).}
\end{popfigure}

\begin{attention}[Not all yellow leaves mean nitrogen deficiency]
The most common diagnostic error --- in the field and in the GCE exam --- is to attribute every yellowing to lack of nitrogen. \textbf{Magnesium} deficiency causes interveinal chlorosis of older leaves, and \textbf{iron} deficiency causes interveinal chlorosis of young leaves. Always note \emph{which} leaves are affected (old or young) and \emph{where} the yellowing occurs (whole leaf, between veins, or margins) before concluding.
\end{attention}

\courssub{Absorption of Minerals by Roots}

Minerals are absorbed from the soil solution as \textbf{ions} (\ce{NO3-}, \ce{NH4+}, \ce{H2PO4-}, \ce{K+}, \ce{Ca^{2+}}, \ce{Mg^{2+}}\ldots), mainly by the \cle{root hairs} just behind the root tip. Root hairs enormously increase the absorptive surface: a single maize plant may carry millions of them.

Absorption occurs by two mechanisms:
\begin{itemize}
\item \textbf{Passive uptake} (diffusion, mass flow with the transpiration stream) when the ion concentration in the soil solution exceeds that in the root cells.
\item \textbf{Active transport}, the dominant mechanism: carrier proteins in the root-cell membranes pump ions \emph{against} their concentration gradient, using energy from ATP supplied by respiration. This is why waterlogged, airless soils (poor in \ce{O2} for respiration) reduce mineral uptake even when minerals are abundant.
\end{itemize}

Two further points are essential:

\textbf{Mycorrhizal associations.} The roots of most crop plants form a mutualistic symbiosis with soil fungi (\cle{mycorrhizae}). The fungal hyphae extend far beyond the root hairs, greatly increasing the absorbing surface and improving uptake of poorly mobile ions such as phosphate; in return the fungus receives carbohydrates from the plant.

\textbf{Effect of soil pH.} Nutrient availability depends strongly on pH (link with Section 1). In strongly acid soils (pH $< 5$, common in the humid forest zone of the South and Centre), phosphorus becomes fixed as insoluble iron and aluminium phosphates, and calcium and magnesium are leached away. In alkaline soils (pH $> 7.5$), iron, manganese, zinc and boron become unavailable, producing chlorosis even when the elements are present in the soil. Maximum overall availability occurs around pH 6.0--7.0, which is why liming acid soils is a standard corrective practice.

\courssub{The Nitrogen Cycle and Biological Nitrogen Fixation}

Nitrogen is the element needed in the greatest quantity, yet although air is about 78\% \ce{N2}, plants cannot use gaseous nitrogen directly. They depend on the \cle{nitrogen cycle}, driven largely by micro-organisms:

\begin{enumerate}
\item \textbf{Nitrogen fixation}: conversion of atmospheric \ce{N2} into ammonium compounds, either industrially (Haber process), by lightning, or biologically by bacteria such as free-living \emph{Azotobacter} and symbiotic \emph{Rhizobium}.
\item \textbf{Ammonification}: decomposition of organic nitrogen (proteins, urea in manure and dead matter) into ammonium ions \ce{NH4+} by saprotrophic bacteria and fungi.
\item \textbf{Nitrification}: oxidation of \ce{NH4+} to nitrites \ce{NO2-} (by \emph{Nitrosomonas}) then to nitrates \ce{NO3-} (by \emph{Nitrobacter}). Nitrate is the main form absorbed by crop roots.
\item \textbf{Plant uptake and assimilation}: \ce{NO3-} is absorbed and built into amino acids and proteins, which pass along food chains.
\item \textbf{Denitrification}: under waterlogged, anaerobic conditions, denitrifying bacteria convert \ce{NO3-} back to \ce{N2}, a loss of fertility.
\end{enumerate}

\begin{popfigure}
\begin{center}
\begin{tikzpicture}[scale=1.0, every node/.style={font=\small}]
% Boxes
\node[draw, thick, rounded corners, fill=popblueL, text=popdark, align=center, minimum width=2.6cm] (n2) at (0,3.4) {Atmospheric \ce{N2}};
\node[draw, thick, rounded corners, fill=popgreenL, text=popdark, align=center, minimum width=2.6cm] (nh4) at (5,3.4) {Ammonium \ce{NH4+}};
\node[draw, thick, rounded corners, fill=popgreenL, text=popdark, align=center, minimum width=2.6cm] (no3) at (5,0) {Nitrates \ce{NO3-}};
\node[draw, thick, rounded corners, fill=popblueL, text=popdark, align=center, minimum width=2.6cm] (org) at (0,0) {Plant \& animal\\ protein};
% Arrows
\draw[->, very thick, poppurple] (n2) -- node[above, align=center]{Fixation\\ (\emph{Rhizobium}, lightning)} (nh4);
\draw[->, very thick, poporange] (nh4) -- node[right, align=center]{Nitrification\\ (\emph{Nitrosomonas},\\ \emph{Nitrobacter})} (no3);
\draw[->, very thick, popgreen] (no3) -- node[below, align=center]{Root uptake \&\\ assimilation} (org);
\draw[->, very thick, popgold] (org) -- node[left, align=center]{Ammonification\\ (decomposers)} (nh4);
\draw[->, very thick, poppink, dashed] (no3.north west) to[bend right=25] node[above, sloped, align=center]{Denitrification (anaerobic)} (n2.south east);
\end{tikzpicture}
\end{center}
\legende{The nitrogen cycle: fixation, ammonification, nitrification, plant uptake and denitrification.}
\end{popfigure}

\begin{savaistu}[Rhizobium and the Cameroonian mixed farm]
Legumes such as groundnut, cowpea (ni\'eb\'e), soya and beans host \emph{Rhizobium} bacteria in swellings on their roots called \cle{root nodules}. The bacterium fixes atmospheric \ce{N2} and passes ammonium compounds to the plant; the plant supplies the bacterium with carbohydrates. When the legume residues are ploughed in, the soil is enriched in nitrogen --- which is why Cameroonian farmers traditionally intercrop maize with cowpea or rotate groundnut with cereals. This natural fertilisation reduces the need for expensive nitrogen fertiliser and is a classic GCE example of applied symbiosis.
\end{savaistu}

\courssub{Fertilisers: Organic and Inorganic}

\begin{definition}[Fertiliser]
A \cle{fertiliser} is any material, organic or inorganic, added to the soil (or to a nutrient solution) to supply one or more essential plant nutrients and so maintain or increase crop yield.
\end{definition}

\begin{center}
\begin{tabularx}{\linewidth}{|X|X|X|}
\hline
\rowcolor{popblueL} \textbf{Type} & \textbf{Advantages} & \textbf{Disadvantages} \\
\hline
\textbf{Organic}: farmyard manure, compost, green manure (ploughed-in legumes), crop residues & Release nutrients slowly; improve soil structure, water retention and humus content; feed soil organisms; cheap and locally available & Bulky to transport; nutrient content low and variable; slow action; may carry weed seeds or pathogens if not well rotted \\
\hline
\textbf{Inorganic (mineral)}: NPK compounds, urea (46\% N), ammonium sulphate (21\% N + S), single superphosphate & Concentrated and fast-acting; precise, known composition; easy to store and apply; doses can match crop needs & Expensive; no improvement of soil structure; easily leached (nitrates); risk of over-use: eutrophication, soil acidification (ammonium sulphate), salt damage \\
\hline
\end{tabularx}
\end{center}

\textbf{Correct application methods.}
\begin{itemize}
\item \textbf{Broadcasting}: spreading evenly over the whole field, usually before planting.
\item \textbf{Placement (localised) application}: putting the fertiliser in the planting hole or in a band a few centimetres beside and below the seed --- more efficient, and avoids scorching the seed.
\item \textbf{Top-dressing}: applying nitrogen (e.g.\ urea) to a growing crop, e.g.\ maize at knee-height; should be done on moist soil and lightly covered to limit volatilisation losses.
\item \textbf{Timing and dose}: match the crop's demand curve; split nitrogen applications rather than applying one heavy dose.
\end{itemize}

\begin{exoresolu}[Reading an NPK label and calculating a dose]
A farmer buys a 50 kg bag of compound fertiliser labelled \textbf{15-15-15}. (a) What mass of nitrogen does the bag contain? (b) The recommended nitrogen dose for his maize crop is 60 kg of N per hectare. How many bags must he apply per hectare?
\tcblower
\textbf{(a)} The first number of the label gives the percentage of nitrogen by mass:
\[ m(\mathrm{N}) = 50\ \mathrm{kg} \times \frac{15}{100} = 7.5\ \mathrm{kg\ of\ N\ per\ bag}. \]
(The bag also contains 7.5 kg of \ce{P2O5} and 7.5 kg of \ce{K2O}.)

\textbf{(b)} Number of bags required:
\[ n = \frac{60\ \mathrm{kg\ N\,ha^{-1}}}{7.5\ \mathrm{kg\ N\ per\ bag}} = 8\ \mathrm{bags\ per\ hectare}. \]
The farmer must therefore apply 8 bags (400 kg) of 15-15-15 per hectare to supply 60 kg of nitrogen.
\end{exoresolu}

\begin{popfigure}
\begin{center}
\begin{tikzpicture}
\begin{axis}[
  width=0.85\linewidth, height=6.5cm,
  xlabel={Fertiliser applied (kg ha$^{-1}$)},
  ylabel={Crop yield (t ha$^{-1}$)},
  xmin=0, xmax=420, ymin=0, ymax=10,
  xtick={0,100,200,300,400},
  ytick={0,2,4,6,8,10},
  grid=both, grid style={popline!40},
  axis line style={popdark}, label style={popdark}, tick label style={popdark, font=\small},
  legend style={at={(0.97,0.4)}, anchor=east, font=\small, draw=popline}
]
\addplot[very thick, poporange, domain=0:400, samples=80] {2 + 8*(1-exp(-x/90))};
\addlegendentry{Yield response}
\addplot[dashed, thick, popmuted, domain=0:400, samples=2] {2 + 0.08*x};
\addlegendentry{Constant return (hypothetical)}
\draw[dashed, popgreen, thick] (axis cs:0,10) -- (axis cs:400,10);
\node[font=\small, text=popdark, anchor=south west] at (axis cs:8,9.1) {Yield ceiling (other factors now limiting)};
\node[font=\small, text=popdark, align=center] at (axis cs:300,5.6) {Diminishing\\ returns};
\end{axis}
\end{tikzpicture}
\end{center}
\legende{Yield versus fertiliser rate: each additional unit of fertiliser gives a smaller yield increase (law of diminishing returns) until a ceiling is reached where another factor becomes limiting.}
\end{popfigure}

The curve above illustrates the \cle{law of diminishing returns}: beyond the economic optimum, extra fertiliser costs more than the extra yield is worth, and the excess leaches into rivers.

\begin{attention}[Dangers of over-fertilisation]
Applying too much mineral fertiliser does not give unlimited yield. Excess nitrate leaches into streams and lakes, causing \cle{eutrophication}: algal blooms, deoxygenation and death of aquatic life. Continuous use of ammonium sulphate \textbf{acidifies} the soil; heavy urea or salt-based fertiliser doses can ``burn'' seedlings. Fertiliser is a medicine: the right dose cures, the wrong dose poisons.
\end{attention}

\courssub{Experimental Evidence: Water Culture (Hydroponics)}

How do we know that, say, magnesium is essential? By growing plants in a solution containing \emph{everything except} magnesium and observing the result. This is the \textbf{culture solution} (water culture) technique, and \cle{hydroponics} is its commercial extension.

\begin{methode}[Setting up a culture-solution experiment to test essentiality]
\begin{enumerate}
\item Prepare a \textbf{complete culture solution} containing known salts supplying all essential elements (e.g.\ potassium nitrate, calcium phosphate, magnesium sulphate, an iron salt and a trace-element mixture) in distilled water.
\item Prepare a series of \textbf{deficient solutions}, each identical to the complete solution but lacking \emph{one} element (e.g.\ replace magnesium sulphate by potassium sulphate to remove Mg while keeping S).
\item Place seedlings of the same species, age and size (e.g.\ young maize or tomato) in each solution, supporting them so that only the roots are immersed; wrap the containers in foil to exclude light and prevent algal growth.
\item \textbf{Aerate} the solutions continuously (aquarium pump) so that roots can respire; renew the solutions regularly to keep composition and pH stable.
\item Grow all plants under identical light, temperature and humidity for several weeks.
\item \textbf{Observe and record}: plants in the complete solution grow normally; each deficient solution produces its characteristic symptom (e.g.\ interveinal chlorosis of old leaves without Mg). \textbf{Confirm} by re-adding the missing element: recovery proves essentiality.
\end{enumerate}
\end{methode}

\begin{exoresolu}[Diagnosing a field deficiency]
A farmer near Bafoussam notices that the \emph{older} leaves of his maize turn yellow \emph{between the veins while the veins remain green}; the young leaves look normal. (a) Name the likely deficient element and justify. (b) State the element's function. (c) Suggest two corrections.
\tcblower
\textbf{(a)} The symptom is \textbf{interveinal chlorosis on older leaves}. Because the element concerned is \emph{mobile} in the plant, it is withdrawn from old leaves to supply young ones, so old leaves suffer first. Interveinal chlorosis of old leaves is characteristic of \textbf{magnesium deficiency} (nitrogen would yellow the whole leaf uniformly; iron would attack young leaves).

\textbf{(b)} Magnesium is the \textbf{central atom of the chlorophyll molecule}; without it chlorophyll cannot be formed, hence the yellowing and reduced photosynthesis.

\textbf{(c)} Corrections: apply a magnesium-containing fertiliser (e.g.\ magnesium sulphate / kieserite, or dolomitic lime if the soil is also acid); and add organic matter (compost or farmyard manure), which releases Mg slowly and improves its retention in the soil. Checking and correcting soil acidity by liming also improves Mg availability.
\end{exoresolu}

\begin{exercice}[Deficiency symptoms table]
Copy and complete a table with three columns --- \emph{Element / Function / Deficiency symptom} --- for N, P, K, Mg, Ca and Fe. Then answer: a tomato grower observes death of the shoot tips and rotting of the blossom end of young fruits. Which element is deficient, and why do the symptoms appear on the youngest tissues?
\end{exercice}

\begin{corrige}[Exercise 1]
Table: see the summary table in this section (N: protein/chlorophyll, uniform chlorosis of old leaves; P: ATP/nucleic acids, poor roots and purpling; K: enzyme activation/stomata, marginal scorch of old leaves; Mg: chlorophyll centre, interveinal chlorosis of old leaves; Ca: cell walls, death of growing points; Fe: chlorophyll synthesis, interveinal chlorosis of young leaves).

The tomato grower's plants lack \textbf{calcium}. Calcium is \emph{immobile} in the plant (it is fixed in cell walls as calcium pectate), so it cannot be moved from old tissues to new ones; the youngest tissues --- shoot tips and developing fruits --- are therefore the first to show deficiency, giving death of growing points and blossom-end rot.
\end{corrige}

\begin{exercice}[Fertiliser calculation]
A bag of urea is labelled 46\% N and has a mass of 50 kg. (a) Calculate the mass of nitrogen in one bag. (b) A yam grower wishes to supply 92 kg of nitrogen per hectare. How many bags of urea does he need per hectare? (c) Give one reason why he should split this dose into two or three applications rather than applying it all at planting.
\end{exercice}

\begin{corrige}[Exercise 2]
\textbf{(a)} $m(\mathrm{N}) = 50 \times \dfrac{46}{100} = 23$ kg of N per bag.

\textbf{(b)} $n = \dfrac{92}{23} = 4$ bags per hectare.

\textbf{(c)} Nitrogen from urea is rapidly converted to nitrate, which is easily \textbf{leached} by heavy rain before the crop can absorb it; splitting the dose matches supply to the crop's demand, reduces leaching losses (and the risk of eutrophication) and avoids scorching young plants with a single concentrated dose.
\end{corrige}

\begin{aretenir}[Key points for the GCE]
\begin{itemize}
\item Plants are autotrophs: \ce{CO2} and water give C, H, O; all other elements come from the soil solution.
\item Macronutrients: N, P, K, Ca, Mg, S. Micronutrients: Fe, Mn, B, Zn, Cu, Mo, Cl. Essentiality is proved by culture-solution experiments.
\item Learn deficiency symptoms \textbf{in table form} (element -- function -- symptom): this is a frequent GCE short-answer item. Remember mobility: mobile elements $\rightarrow$ old leaves first; immobile $\rightarrow$ young leaves first.
\item N, P and K most often limit yield $\Rightarrow$ NPK fertilisers; the label numbers give \% N, \% \ce{P2O5}, \% \ce{K2O}; yield follows the law of diminishing returns.
\item \emph{Rhizobium} in legume nodules fixes atmospheric nitrogen --- the biological basis of intercropping and rotation in Cameroonian mixed farming.
\item Over-fertilisation causes eutrophication and soil acidification: correct diagnosis and correct dose are the farmer's (and the candidate's) watchwords.
\end{itemize}
\end{aretenir}

\coursec{Cultivation Techniques of Food and Non-Food Crops}

In the previous section we established that a crop can only yield well if the soil supplies it with the right mineral elements in the right amounts. But even the most fertile soil in the Centre Region will give a poor harvest if the farmer plants the wrong variety, at the wrong time, at the wrong spacing, and never weeds. Feeding the plant is necessary; \emph{managing} the plant is equally essential. This section takes you through the complete \cle{cultivation} sequence — from the choice of crop and the preparation of the land to harvesting, storage and marketing — using Cameroonian crops as our working examples.

\courssub{Classification of Crops}

Farmers and agricultural officers group crops according to their main use.

\begin{definition}[Food crop, cash crop]
A \cle{food crop} is a crop grown primarily for human (or animal) consumption: maize, cassava, plantain, yam, cocoyam, groundnut, cowpea, rice, millet, sorghum. A \cle{cash crop} (or non-food/industrial crop) is grown primarily for \emph{sale} — for export or for industry — rather than for direct consumption: cocoa, coffee, cotton, oil palm, rubber, tobacco, banana (export), tea, sugar cane. Note that the distinction is about \emph{purpose}, not biology: a farmer may sell part of his maize and eat part of his cocoa-bean income.
\end{definition}

\begin{center}
\begin{tabularx}{\linewidth}{|l|X|X|}
\hline
\rowcolor{popblueL} \textbf{Criterion} & \textbf{Food crops} & \textbf{Cash crops} \\
\hline
Main purpose & Feeding the family / local markets & Sale, export, agro-industry \\
\hline
Examples in Cameroon & Maize, cassava, plantain, yam, groundnut, cowpea, rice, millet, sorghum & Cocoa, coffee, cotton, oil palm, rubber, tobacco, banana (export), tea, sugar cane \\
\hline
Typical scale & Smallholdings (0.5--3 ha) & Smallholdings \emph{or} large plantations (CDC, SOCAPALM, HEVECAM) \\
\hline
Life cycle & Mostly annual or short-cycle & Mostly perennial tree crops \\
\hline
\end{tabularx}
\end{center}

\courssub{Farming Systems in Cameroon}

A \cle{farming system} is the overall way a farmer organises land, labour and crops.

\begin{itemize}
\item \textbf{Shifting cultivation}: a piece of forest or bush is cleared, burned, cropped for 2--3 years, then abandoned when fertility falls; the farmer moves to a new plot. Simple but destructive when population pressure shortens the fallow cycle below the threshold needed for soil recovery.
\item \textbf{Bush fallowing}: an improvement on shifting cultivation — plots are cultivated in rotation and left to \emph{fallow} (rest under natural vegetation) for several years so that soil fertility is restored before re-cropping. Common in the humid forest zone where land is still relatively abundant.
\item \textbf{Subsistence farming}: production aimed mainly at feeding the farmer's household, with small surpluses sold locally. Typical of food-crop farms in the West and North-West Regions.
\item \textbf{Plantation agriculture}: large estates growing a single cash crop (monoculture) with hired labour and processing facilities, e.g.\ CDC banana and rubber estates in the South-West, SOCAPALM oil palm estates in the Littoral and South-West, HEVECAM rubber in the Centre and South.
\item \textbf{Mixed farming}: crops and livestock are kept on the same farm; crop residues feed the animals and manure fertilises the fields — common in the Adamawa and North-West highlands.
\item \textbf{Market gardening (maraîchage)}: intensive production of vegetables (tomato, cabbage, carrot, onion, green beans, huckleberry) on small, well-watered plots near towns, e.g.\ the valleys around Bafoussam, Santa, Dschang and the Wouri lowlands near Douala.
\item \textbf{Agroforestry systems}: deliberate integration of trees with crops and/or livestock — e.g.\ cocoa under shade trees, coffee under \emph{Albizia} or \emph{Leucaena}, or \emph{Acacia} alley cropping in the Sudan-Sahel zone.
\end{itemize}

\courssub{Land Preparation}

Land preparation creates a suitable \cle{seedbed}: loose, aerated, weed-free soil in which roots can penetrate and water can infiltrate.

\begin{enumerate}
\item \textbf{Clearing}: removal of trees, shrubs and previous crop residues (cutlass, axe, chainsaw, or bulldozer on plantations). In forest zones clearing is done at the end of the rainy season (November--January); debris may be burned (adds ash/potash but destroys organic matter and soil organisms) or left as mulch (slash-and-mulch).
\item \textbf{Ploughing}: the primary tillage — turning over the topsoil (15--25 cm depth) with a hoe, animal-drawn plough or tractor to bury weeds and loosen the soil. On steep slopes, contour ploughing or minimum tillage reduces erosion.
\item \textbf{Harrowing}: secondary tillage — breaking the clods left by ploughing into a fine tilth and levelling the surface. Done with a harrow, rake or hoe.
\item \textbf{Ridging and bed making}: soil is heaped into ridges (yam, groundnut, potato, cassava in wet areas) or raised beds (vegetables, nursery) to improve drainage and give tubers room to swell. In the Sudan-Sahel, tied ridges (with cross-ties) conserve runoff water.
\end{enumerate}

\textbf{Timing with the seasons.} In the humid forest zone (bimodal rainfall), the main planting follows the onset of the heavy rains (March--April) with a second season around August--September. In the savanna zones (unimodal rainfall), land is prepared at the end of the dry season (February--April) and planting waits for the first reliable rains (May--June). Planting too early exposes seeds to drought; too late shortens the growing season and increases pest pressure.

\begin{popfigure}
\begin{tikzpicture}[font=\small]
\tikzset{box/.style={draw=popblue, fill=popblueL, rounded corners, text width=3.1cm, align=center, minimum height=0.9cm},
arr/.style={->, thick, popdark}}
\node[box] (a) at (0,0) {Land preparation (Feb--Mar)};
\node[box] (b) at (4.6,0) {Seed selection \& planting (Mar--Apr)};
\node[box] (c) at (9.2,0) {Maintenance: weeding, fertiliser, earthing-up};
\node[box] (d) at (9.2,-2.2) {Harvesting (Jul--Aug / dry season)};
\node[box] (e) at (4.6,-2.2) {Drying \& shelling};
\node[box] (f) at (0,-2.2) {Storage / marketing};
\draw[arr] (a) -- (b);
\draw[arr] (b) -- (c);
\draw[arr] (c) -- (d);
\draw[arr] (d) -- (e);
\draw[arr] (e) -- (f);
\draw[arr, dashed, popgreen] (f.west) .. controls (-1.6,-1.1) .. (a.west) node[midway, left, text=popgreen, align=center, text width=1.6cm]{\scriptsize seed kept for next season};
\end{tikzpicture}
\legende{Production cycle of maize, from land preparation to storage. Part of the harvest is kept as seed for the next season.}
\end{popfigure}

\courssub{Propagation Methods}

Crops are multiplied either by \cle{seeds} (sexual propagation) or by plant parts (asexual/vegetative propagation).

\begin{definition}[Vegetative propagation]
\cle{Vegetative propagation} is the production of new plants from vegetative parts — stems (cuttings, suckers), underground organs (tubers, rhizomes, corms, setts) — or by artificial joining of plant parts (grafting, budding, layering) — without seeds. The offspring are genetically identical to the parent (a clone), so desirable characters are preserved; the drawback is that diseases of the parent are transmitted and there is no genetic variation for adaptation.
\end{definition}

\begin{center}
\begin{tabularx}{\linewidth}{|l|X|X|}
\hline
\rowcolor{popblueL} \textbf{Method} & \textbf{Principle} & \textbf{Cameroonian examples} \\
\hline
Seed (sexual) & Sowing of fertilised ovules & Maize, cowpea, groundnut, rice, millet, sorghum, cocoa, oil palm (nursery) \\
\hline
Stem cutting & A stem piece (20--30 cm, 5--7 nodes) roots and sprouts & Cassava, sugar cane, sweet potato vines, \emph{Gmelina}, \emph{Leucaena} \\
\hline
Tuber / sett & Underground organ with buds (``eyes'') & Yam (setts 200--500 g), potato (seed tubers), cocoyam (corms/cormels) \\
\hline
Sucker & Shoot from the base of the mother plant & Plantain, banana, pineapple, sisal \\
\hline
Rhizome division & Underground stem with buds & Ginger, turmeric \\
\hline
Grafting & A scion is joined onto a rooted stock (whip, cleft, side) & Mango, citrus, avocado, improved rubber (GT1, PB series) \\
\hline
Budding & A single bud is inserted under the stock bark (T-bud, patch bud) & Rubber, citrus, cocoa (improved clones from IRAD/SODECAO) \\
\hline
Layering & Stem rooted while still attached to mother & Citrus, mango, some ornamentals \\
\hline
\end{tabularx}
\end{center}

\begin{popfigure}
\begin{tikzpicture}[font=\small, scale=1]
% Cutting
\draw[thick, popgreen] (0,0) -- (0,1.8);
\draw[thick, popgreen] (0,1.2) -- (0.5,1.7);
\draw[thick, popgreen] (0,0.8) -- (-0.45,1.3);
\foreach \x in {-0.25,-0.1,0.05,0.2} \draw[popdark] (\x,0) -- (\x*1.6,-0.55);
\node[align=center] at (0,-1.0) {\textbf{Stem cutting}\\ \scriptsize stem piece roots (cassava)};
% Grafting
\begin{scope}[xshift=4.6cm]
\draw[thick, popdark] (0,-0.6) -- (0,0.5);
\draw[thick, poporange] (0,0.5) -- (0,1.8);
\draw[thick, poporange] (0,1.3) -- (0.5,1.8);
\draw[popink, thick] (-0.12,0.5) -- (0.12,0.5);
\foreach \x in {-0.2,0,0.2} \draw[popdark] (\x,-0.6) -- (\x*1.5,-1.1);
\node[align=center] at (0,-1.55) {\textbf{Grafting}\\ \scriptsize scion on rooted stock};
\node[popink, right, align=left] at (0.2,0.5) {\scriptsize union};
\end{scope}
% Budding
\begin{scope}[xshift=9.2cm]
\draw[thick, popdark] (0,-0.6) -- (0,1.6);
\fill[poporange] (0.06,0.7) ellipse (0.09 and 0.2);
\draw[popink] (0.02,0.35) -- (0.02,1.05);
\foreach \x in {-0.2,0,0.2} \draw[popdark] (\x,-0.6) -- (\x*1.5,-1.1);
\node[align=center] at (0,-1.55) {\textbf{Budding}\\ \scriptsize single bud in T-cut};
\node[popink, right, align=left] at (0.18,0.7) {\scriptsize bud};
\end{scope}
\end{tikzpicture}
\legende{Three vegetative propagation methods: stem cutting, grafting (scion joined to stock) and budding (a single bud inserted in a T-shaped cut).}
\end{popfigure}

\textbf{Seed selection and treatment.} Only bold, clean, undamaged seeds of a good variety should be planted. Diseased or weevil-bored grains give weak seedlings. Seeds may be \emph{treated} before sowing: dressing with fungicide/insecticide against soil pests (e.g.\ maize seed dressed with thiram + imidacloprid), or soaking (e.g.\ rice for 24 h, maize for 6--12 h) to speed germination. For legumes, inoculation with \emph{Rhizobium} strains specific to the crop enhances nitrogen fixation.

\begin{methode}[Carrying out a seed germination test]
\begin{enumerate}
\item Take a representative sample of the seed lot and count out 100 seeds (or 4 replicates of 25).
\item Place them on moist cotton wool, blotting paper or sand in a covered dish; keep warm (about $25$--$30\ ^{\circ}\mathrm{C}$) and moist, in the dark or diffuse light.
\item After the normal germination period for the species (e.g.\ 5--7 days for maize, 7--10 days for cocoa), count the seeds that have produced a normal radicle and plumule (discard abnormal seedlings).
\item Calculate the percentage germination:
\[ \text{Germination (\%)} = \frac{\text{number of normal seedlings}}{\text{number of seeds sown}} \times 100 \]
\item Decide: a lot giving less than about $80\ \%$ germination (for most cereals and legumes) should be rejected or sown at a higher rate. For cocoa, minimum is $90\ \%$.
\end{enumerate}
\end{methode}

\begin{exoresolu}[Germination test and seed requirement]
A farmer in Ngaoundéré tests 50 maize seeds; 43 germinate normally. He wants about $40\,000$ plants per hectare. (a) Calculate the percentage germination. (b) How many viable seeds per hectare must he sow, assuming one seed per stand?
\tcblower
(a) $\text{Germination} = \dfrac{43}{50}\times 100 = 86\ \%$. The lot is acceptable (above $80\ \%$ for maize).

(b) Only $86\ \%$ of sown seeds will give a plant, so:
\[ \text{Seeds to sow} = \frac{40\,000}{0.86} \approx 46\,512 \text{ seeds per hectare.} \]
He should sow about $46\,500$ seeds per hectare (roughly $12$ kg at about $3\,900$ seeds/kg) to obtain the target stand.
\end{exoresolu}

\courssub{Planting Practices}

\textbf{Spacing and depth.} Each crop has an optimum spacing determined by its growth habit, root system and target yield. Sowing depth must suit seed size and soil type: small seeds shallow (1--2 cm), large seeds or cuttings deeper (5--10 cm for maize; cassava stakes buried two-thirds of their length, buds upward).

\begin{center}
\begin{tabular}{|l|c|c|c|}
\hline
\rowcolor{popblueL} \textbf{Crop} & \textbf{Spacing (cm)} & \textbf{Depth (cm)} & \textbf{Plant population (stands/ha)} \\
\hline
Maize (grain) & 75 $\times$ 25 (1 seed/stand) & 5--7 & 53\,333 \\
Maize (green) & 50 $\times$ 25 (2 seeds/stand) & 5--7 & 160\,000 \\
Cassava & 100 $\times$ 100 & stake 2/3 buried & 10\,000 \\
Yam & 100 $\times$ 100 & sett 5--10 & 10\,000 \\
Plantain / Banana & 300 $\times$ 300 & sucker in 30 cm hole & 1\,111 \\
Cocoa & 300 $\times$ 300 & seedling in 40 $\times$ 40 $\times$ 40 cm hole & 1\,111 \\
Oil palm & 900 $\times$ 900 $\times$ 900 (triangular) & seedling in 60 cm cube hole & 143 \\
Groundnut & 50 $\times$ 15 (1 seed/stand) & 5--6 & 133\,333 \\
Rice (lowland) & 20 $\times$ 20 (3--4 seedlings/hill) & 2--3 (transplanted) & 250\,000 hills \\
Cotton & 90 $\times$ 30 (2--3 seeds/stand) & 3--5 & 74\,000 \\
\hline
\end{tabular}
\end{center}

\begin{propriete}[Spacing and yield]
Correct spacing maximises yield per hectare because it minimises \cle{competition} between plants for light, water and mineral nutrients while using the land fully. Too dense a stand gives many small, etiolated plants and low yield per plant; too wide a stand wastes land and lets weeds invade. (Not examinable as a proof, but you must be able to \emph{explain} it in terms of intraspecific competition and light interception.)
\end{propriete}

\textbf{Nursery techniques.} Some crops are first raised in a \cle{nursery} — a small, sheltered, well-watered seedbed — then transplanted. This allows close care (watering, shading, protection from pests) during the delicate seedling stage and gives a head start on weeds.

\begin{itemize}
\item \textbf{Cocoa}: seeds sown in shaded nursery beds or polythene bags (4 $\times$ 6 cm); transplanted at 4--6 months (4--6 leaves) into field at $3 \times 3$ m.
\item \textbf{Oil palm}: pre-nursery (3--4 months in small bags) $\rightarrow$ main nursery (12--15 months in large bags) $\rightarrow$ field at $9 \times 9 \times 9$ m triangular.
\item \textbf{Vegetables} (tomato, cabbage, onion, pepper): raised in nursery beds (1 m wide, raised 15 cm, shaded) and transplanted at 3--5 weeks (4--6 true leaves).
\item \textbf{Coffee}: seeds in nursery beds, transplanted at 6--9 months.
\end{itemize}

\begin{attention}[Planting cuttings the right way up]
A classic and costly mistake: cassava stakes and yam setts planted \textbf{upside down} or \textbf{too deep} sprout poorly or rot. Cassava stakes must be planted with the buds (nodes) pointing upward; yam setts with the sprouting end up, about 5--10 cm deep. Reversed or deeply buried cuttings give weak, late sprouts and a sharply reduced yield.
\end{attention}

\begin{savaistu}[Improved planting material in Cameroon]
IRAD (Institute of Agricultural Research for Development) and SODECAO (Cocoa Development Corporation) distribute improved varieties: maize (CMS 8704, CMS 9015, TZEE-Y), cassava (TMS 30572, TMS 4(2)1425 resistant to mosaic), cocoa (SNKA 12, SNKA 13, SNKA 14 — high yield, black pod tolerant), oil palm (Tenera hybrids from Dami and La Mé). Always source seed from certified multipliers to ensure genetic purity and phytosanitary quality.
\end{savaistu}

\courssub{Crop Maintenance}

After planting, the farmer must protect the crop's access to light, water and nutrients:

\begin{itemize}
\item \textbf{Weeding}: removes competitors; usually 2--3 weedings per season for annual crops (e.g.\ maize at 2--3 and 6--7 weeks after planting). Methods: hand hoe, cutlass, animal-drawn cultivator, or selective herbicides (pre-emergence: atrazine + metolachlor for maize; post-emergence: nicosulfuron for grasses).
\item \textbf{Mulching}: covering the soil with dry grass, leaves, crop residues or plastic film to conserve moisture, suppress weeds, moderate soil temperature and add organic matter (standard practice in plantain, cocoa young plantations and market-garden plots).
\item \textbf{Watering / irrigation}: essential for nurseries, market gardening and dry-season cropping. Methods: watering can, hose, sprinkler, drip irrigation (increasingly used for high-value vegetables), flood irrigation (irrigated rice at Yagoua, Maga, SEMRY schemes).
\item \textbf{Staking}: supporting climbing or heavy-bearing plants (yam vines — 2 m stakes, tomato — stakes or trellis, passion fruit) to keep fruits off the soil and improve air circulation.
\item \textbf{Pruning}: removing dead, diseased or excess branches (cocoa — chupons and low branches; coffee — multiple stem pruning, capping at 1.5--2 m; fruit trees — formative and maintenance pruning) to improve aeration, light penetration and yield.
\item \textbf{Thinning}: removing excess seedlings so the remaining ones have the correct spacing (maize: thin to 1 plant/stand at 2 weeks; vegetables: thin to final spacing at 2--3 true leaves).
\item \textbf{Earthing-up}: heaping soil around the base of plants (maize at 6--7 weeks to support stems and cover brace roots; groundnut at flowering to encourage peg penetration; potato to prevent tuber greening; cassava in waterlogged areas).
\item \textbf{Fertiliser application}: basal (NPK at planting) and top-dressing (urea at knee-high for maize, at flowering for groundnut). Rates based on soil test; typical: maize 200 kg/ha NPK 20-10-10 + 100 kg/ha urea; cocoa 300 kg/ha NPK 12-12-17+2MgO+Te per year split in two applications.
\end{itemize}

\courssub{Pests, Diseases and Weeds of Major Crops}

A \cle{pest} is any organism that damages the crop; control methods fall into three groups:

\begin{itemize}
\item \textbf{Cultural control}: practices that prevent or reduce attack — crop rotation, early planting, resistant/tolerant varieties, destruction of crop residues, clean planting material, proper spacing, field hygiene, trap crops (e.g.\ \emph{Desmodium} for stem borer push-pull).
\item \textbf{Biological control}: use of natural enemies — e.g.\ parasitoid wasp \emph{Anagyrus lopezi} against cassava mealybug (classical success in 1980s), \emph{Bacillus thuringiensis} (Bt) against caterpillars, \emph{Trichoderma} spp. against soil-borne fungi, neem (\emph{Azadirachta indica}) extracts as botanical insecticide.
\item \textbf{Chemical control}: pesticides (insecticides, fungicides, herbicides, nematicides). Effective but costly, polluting and dangerous if misused; always a last resort, applied at the correct dose with protective equipment (PPE), respecting pre-harvest intervals (PHI).
\end{itemize}

\begin{definition}[Integrated pest management (IPM)]
\cle{Integrated pest management} is the combined, rational use of cultural, biological and (minimal, well-targeted) chemical methods to keep pest populations below the economic damage threshold, while protecting the farmer, the consumer and the environment. Key steps: monitoring $\rightarrow$ threshold decision $\rightarrow$ selective intervention $\rightarrow$ evaluation.
\end{definition}

\begin{center}
\begin{tabularx}{\linewidth}{|l|X|X|}
\hline
\rowcolor{popblueL} \textbf{Crop} & \textbf{Important enemies} & \textbf{Main IPM components} \\
\hline
Maize & Stem borers (\emph{Busseola}, \emph{Chilo}), fall armyworm (\emph{Spodoptera frugiperda}); maize streak virus (MSV); striga (\emph{Striga hermonthica}) parasitic weed & Early planting, resistant varieties (TZEE-Y, CMS 9015), rotation with legumes, push-pull (\emph{Desmodium}/\emph{Napier}), targeted Bt or emamectin benzoate for FAW \\
\hline
Cassava & Cassava mealybug (\emph{Phenacoccus manihoti}), green mite (\emph{Mononychellus tanajoa}); cassava mosaic disease (CMD), cassava brown streak disease (CBSD) & Clean cuttings (certified), resistant varieties (TMS 30572, TME 419), biological control (mealybug), roguing infected plants \\
\hline
Cocoa & Capsid bugs (\emph{Sahlbergella}, \emph{Distantiella}); black pod disease (\emph{Phytophthora megakarya}, \emph{P. palmivora}); swollen shoot virus (CSSV) & Regular harvesting of diseased pods (sanitation), pruning for aeration, shade management, resistant clones (SNKA series), copper-based fungicide sprays (8--10 rounds/yr), vector control for CSSV \\
\hline
Groundnut & Aphids (\emph{Aphis craccivora}) vector of rosette virus; leaf spots (\emph{Cercospora} early/late); termites & Resistant varieties (ICGV-SM 90704, RMP 12), early planting (escape aphids), seed treatment, fungicide (chlorothalonil) for leaf spot, crop rotation \\
\hline
Plantain / Banana & Banana weevil (\emph{Cosmopolites sordidus}), nematodes (\emph{Radopholus}, \emph{Pratylenchus}); black Sigatoka (\emph{Mycosphaerella fijiensis}) & Clean planting material (boiling water treatment 50°C 20 min or tissue culture), mulching, nematicide (if economic), resistant hybrids (FHIA, PITA), deleafing for Sigatoka \\
\hline
Cotton & Bollworms (\emph{Helicoverpa}, \emph{Diparopsis}), aphids, jassids; bacterial blight & Early planting, resistant varieties (IRMA 1243, FK 37), threshold-based sprays (pyrethroids + organophosphates rotation), destruction of ratoon cotton \\
\hline
\end{tabularx}
\end{center}

\courssub{Harvesting, Processing and Storage}

\textbf{Maturity indices} tell the farmer when to harvest. Harvesting too early gives immature, watery produce that rots; too late invites field pests, shattering losses and quality decline.

\begin{center}
\begin{tabularx}{\linewidth}{|l|X|}
\hline
\rowcolor{popblueL} \textbf{Crop} & \textbf{Maturity indices} \\
\hline
Maize (grain) & Husks dry and brown, grains hard (dent stage), moisture content 20--25\% \\
Cassava & 9--18 months after planting (variety dependent), tubers full size, leaves yellowing \\
Yam & 7--10 months, vines yellowing and drying, tuber skin set \\
Plantain / Banana & Fingers plump, angularities disappear, colour change (green to light green/yellow), 12--15 weeks after flowering \\
Cocoa & Pods turn yellow/orange (depending on variety), 5--6 months after pollination \\
Groundnut & Leaves yellow, pods well-filled with veined interior, 100--130 days after planting \\
Rice & Grains hard, 80--85\% straw coloured, moisture 20--22\% \\
Cotton & Bolls fully open, lint fluffy, 150--180 days after planting \\
\hline
\end{tabularx}
\end{center}

\textbf{Harvesting techniques}: hand-picking (cocoa pods with cutlass or hooked knife, taking care not to damage the flower cushion), pulling (groundnut, cotton), digging (yam, cassava, potato — careful to avoid tuber damage), cutting (maize cobs, rice panicles), mechanical combine (large-scale rice, maize in some schemes).

\textbf{Processing and drying}: 
\begin{itemize}
\item \textbf{Maize}: shelled (by hand, stick, or sheller) and dried on clean surfaces (mats, tarpaulins, raised platforms) to $12$--$13\ \%$ moisture; winnowed to remove chaff.
\item \textbf{Cassava}: highly perishable (48--72 h); processed into gari (fermented, roasted granules), flour (dried chips milled), water-fufu (fermented paste), or chips (dried). Detoxification (cyanogenic glucosides) by fermentation, boiling, roasting.
\item \textbf{Cocoa}: beans + pulp removed from pods, fermented 5--6 days in heaps (covered with banana leaves) or wooden boxes (turned every 2 days) — fermentation develops chocolate flavour precursors, kills embryo, reduces bitterness; then sun-dried on raised platforms to $7$--$8\ \%$ moisture.
\item \textbf{Coffee}: wet processing (pulping $\rightarrow$ fermentation 12--36 h $\rightarrow$ washing $\rightarrow$ drying) for arabica; dry processing (whole cherry dried) for robusta.
\item \textbf{Oil palm}: fresh fruit bunches (FFB) processed within 24 h (sterilisation $\rightarrow$ stripping $\rightarrow$ digestion $\rightarrow$ pressing $\rightarrow$ clarification) to crude palm oil (CPO) and palm kernel oil (PKO).
\end{itemize}

\textbf{Storage}: dried grains are kept in granaries (traditional raised cribs), jute or polythene bags on pallets (not directly on floor), or airtight metal silos / hermetic bags (PICS bags), in a dry, ventilated store. The great enemies of stored produce are \cle{storage pests}:

\begin{itemize}
\item \textbf{Primary insects}: weevils (\emph{Sitophilus zeamais} in maize, \emph{S. oryzae} in rice) and lesser grain borer (\emph{Rhyzopertha dominica}) bore into whole grains. \textbf{Secondary insects}: flour beetles (\emph{Tribolium}), saw-toothed grain beetle (\emph{Oryzaephilus}) feed on damaged grain/dust.
\item \textbf{Pulse beetles}: bruchids (\emph{Callosobruchus maculatus} in cowpea, \emph{Acanthoscelides obtectus} in beans) lay eggs on pods/seeds; larvae develop inside.
\item \textbf{Control}: thorough drying ($<13\%$ mc), clean stores (hygiene), insecticide dusting (actellic dust, pirimiphos-methyl) or fumigation (phosphine — certified operators only), airtight storage (hermetic bags, metal silos), use of ash, neem leaves or \emph{Ocimum} in traditional stores.
\item \textbf{Moulds} (\emph{Aspergillus flavus}, \emph{A. parasiticus}) produce dangerous aflatoxins in groundnut, maize, cottonseed when stored wet ($>14\%$ mc). Control: dry properly before storage, ventilate, avoid re-wetting, sort mouldy grains, use aflasafe (biocontrol \emph{A. flavus} atoxigenic strains) in field.
\end{itemize}

\begin{attention}[Aflatoxin — silent killer]
Aflatoxins are potent carcinogens (liver cancer) and immunosuppressants. They are invisible, odourless and heat-stable (not destroyed by cooking). Groundnut and maize in Cameroon frequently exceed the EU limit (4 ppb total aflatoxins). Prevention starts in the field: avoid drought stress at pod-filling, harvest promptly, dry fast on clean surfaces, store dry. Never blend mouldy grain with clean grain — it contaminates the whole lot.
\end{attention}

\courssub{Case Studies: Maize and Cocoa}

\textbf{Maize (food crop, \emph{Zea mays}).} 
\begin{enumerate}
\item \textbf{Site selection}: well-drained loam, pH 5.5--7.0, rainfall 600--1200 mm.
\item \textbf{Land preparation}: clear, plough, harrow at end of dry season (Feb--Mar in bimodal, Apr--May in unimodal). Conservation agriculture: minimum tillage + cover crop (\emph{Mucuna}) in some zones.
\item \textbf{Planting}: at first effective rains. Selected seed (germination $> 80\ \%$, dressed) sown $75 \times 25$ cm, 5--7 cm deep, 1 seed/stand (grain) or 2 seeds/stand (green maize). Basal fertiliser: NPK 20-10-10 at 200 kg/ha in band beside row.
\item \textbf{Maintenance}: 1st weeding at 2--3 weeks; 2nd weeding + earthing-up + top-dress urea (100 kg/ha) at 6--7 weeks (knee-high). Monitor for fall armyworm (scout 20 plants/ha, spray if $>20\%$ plants with fresh damage).
\item \textbf{Harvest}: 90--120 days (grain), 60--70 days (green). Cobs harvested when husks dry, shelled, dried to $12$--$13\%$ mc.
\item \textbf{Storage}: in PICS bags or treated bags on pallets; monitor for weevils.
\item \textbf{Marketing}: sold to traders, cooperatives, or used for poultry feed (major demand from poultry industry in Yaoundé, Douala, Bafoussam).
\end{enumerate}

\textbf{Cocoa (cash crop, \emph{Theobroma cacao}).} Perennial crop of the humid forest zone (rainfall 1500--2500 mm, well-distributed).
\begin{enumerate}
\item \textbf{Nursery}: seeds from approved seed gardens (SODECAO/IRAD) sown in polythene bags (4 $\times$ 6 cm) under 50--60\% shade; watered daily; transplanted at 4--6 months (4--6 hardened leaves).
\item \textbf{Field establishment}: holes $40 \times 40 \times 40$ cm at $3 \times 3$ m (1111 trees/ha), lined with topsoil + 5 kg compost + 100 g NPK. Plant at onset of rains (Mar--Apr or Aug--Sep). Temporary shade: plantain (1:1) or \emph{Leucaena}; permanent shade: \emph{Albizia}, \emph{Terminalia}, \emph{Milicia} (iroko) at 12--15 m spacing.
\item \textbf{Young plantation (0--3 yr)}: weeding (ring weeding 1 m radius), mulching, shade management, fertiliser (NPK 12-12-17+2MgO at 100 g/tree/yr split), pest monitoring.
\item \textbf{Mature plantation (3+ yr)}: pruning (remove chupons, dead wood, low branches — once/yr after main harvest), sanitation harvesting (remove black pods weekly), fertiliser (300 kg/ha/yr NPK 12-12-17+2MgO split Mar/Apr and Sep/Oct), capsid control (8--10 sprays/yr with approved insecticides, rotation of modes of action), black pod control (copper hydroxide/fosetyl-Al sprays at pod set).
\item \textbf{Harvest}: main crop Oct--Mar, light crop May--Aug. Pods cut with knife/hook (no damage to cushion), opened with wooden mallet, beans + pulp heaped/boxed for fermentation.
\item \textbf{Fermentation \& drying}: 5--6 days fermentation (turn day 2 and 4), sun-dried on raised platforms to $7$--$8\%$ mc (beans rattle, no mould).
\item \textbf{Marketing}: sold to licensed buying agents (LBAs) or cooperatives; quality checked (bean count/100g $<100$, moisture $<8\%$, mould $<4\%$, slaty $<8\%$). Price set by National Cocoa and Coffee Board (NCCB) based on world market.
\end{enumerate}

\begin{popfigure}
\begin{tikzpicture}[font=\scriptsize]
\begin{axis}[
width=14cm, height=5.6cm,
xmin=1, xmax=12.9,
ymin=0, ymax=5,
xtick={1,2,3,4,5,6,7,8,9,10,11,12},
xticklabels={J,F,M,A,M,J,J,A,S,O,N,D},
ytick={0.5,1.5,2.5,3.5,4.5},
yticklabels={Forest: harvest cocoa, Forest: plant maize, Forest: clear land, Savanna: plant maize, Savanna: clear land},
axis y line*=left, axis x line*=bottom,
xlabel={Month}, ylabel={},
x label style={font=\small},
every tick label/.style={font=\scriptsize},
]
% Savanna clear land: Feb-Apr
\addplot[draw=none, fill=popgreen, fill opacity=0.7] coordinates {(2,4.1) (4,4.1) (4,4.9) (2,4.9)} \closedcycle;
% Savanna plant maize: May-Jun
\addplot[draw=none, fill=poporange, fill opacity=0.7] coordinates {(5,3.1) (6,3.1) (6,3.9) (5,3.9)} \closedcycle;
% Forest clear land: Dec-Feb
\addplot[draw=none, fill=popgreen, fill opacity=0.7] coordinates {(1,2.1) (2,2.1) (2,2.9) (1,2.9)} \closedcycle;
\addplot[draw=none, fill=popgreen, fill opacity=0.7] coordinates {(12,2.1) (12.9,2.1) (12.9,2.9) (12,2.9)} \closedcycle;
% Forest plant maize: Mar-Apr (main) and Aug-Sep (second)
\addplot[draw=none, fill=poporange, fill opacity=0.7] coordinates {(3,1.1) (4,1.1) (4,1.9) (3,1.9)} \closedcycle;
\addplot[draw=none, fill=poporange, fill opacity=0.5] coordinates {(8,1.1) (9,1.1) (9,1.9) (8,1.9)} \closedcycle;
% Forest harvest cocoa: Oct-Dec (main) and May-Jul (light)
\addplot[draw=none, fill=popblue, fill opacity=0.7] coordinates {(10,0.1) (12,0.1) (12,0.9) (10,0.9)} \closedcycle;
\addplot[draw=none, fill=popblue, fill opacity=0.5] coordinates {(5,0.1) (7,0.1) (7,0.9) (5,0.9)} \closedcycle;
\end{axis}
\end{tikzpicture}
\legende{Seasonal calendar of cropping activities in Cameroon. Dark bars: main season; pale bars: secondary season. The humid forest zone has two rainy seasons (two maize plantings possible); the savanna zone has one.}
\end{popfigure}

\begin{aretenir}[Key points]
\begin{itemize}
\item Food crops feed people; cash crops are grown for sale — the distinction is one of purpose, not biology.
\item The cultivation sequence is always: land preparation $\rightarrow$ propagation/planting $\rightarrow$ maintenance $\rightarrow$ protection $\rightarrow$ harvesting $\rightarrow$ processing/storage.
\item Vegetative propagation gives uniform clones but transmits diseases; seeds give variation and are easier to store and transport.
\item Correct spacing, depth and orientation of planting material determine germination, sprouting and final yield.
\item IPM combines cultural, biological and chemical control; proper drying is the key to safe storage and aflatoxin prevention.
\item Improved varieties (IRAD, SODECAO, IITA, CIRAD) and certified planting material are essential for productivity and disease management.
\end{itemize}
\end{aretenir}

\begin{attention}[Exam tip — structuring a crop essay]
For any GCE essay asking you to ``describe the cultivation of a named crop'', organise your answer in the chronological order of the production cycle: \textbf{land preparation $\rightarrow$ planting (propagation material, spacing, depth) $\rightarrow$ maintenance (weeding, fertiliser, pest and disease control) $\rightarrow$ harvesting (maturity index, technique) $\rightarrow$ processing and storage}. Add the climatic/soil requirements at the start and marketing at the end for full marks. Always \emph{name} a real crop (maize, cocoa, cassava, plantain) and give real figures (spacing, duration, fertiliser rates, plant population).
\end{attention}

\begin{exercice}[Spacing and plant population]
A farmer plants cassava at a spacing of $1\ \mathrm{m} \times 1\ \mathrm{m}$ on a square field of side $100$ m. (a) How many stands does he plant? (b) If each stand yields on average 4 kg of tubers, what is the expected yield in tonnes per hectare? (c) Explain why doubling the number of stands by planting at $0.5 \times 0.5$ m would probably \emph{reduce} the yield per stand.
\end{exercice}

\begin{corrige}[Exercise 1]
(a) Area of field $= 100 \times 100 = 10\,000\ \mathrm{m^2} = 1$ ha. Each stand occupies $1\ \mathrm{m^2}$, so he plants $10\,000$ stands.

(b) Yield $= 10\,000 \times 4 = 40\,000$ kg $= 40$ tonnes per hectare.

(c) At $0.5 \times 0.5$ m, each stand has only $0.25\ \mathrm{m^2}$. Intraspecific competition for light, water and mineral nutrients becomes severe: the plants shade one another, roots overlap, and each plant produces fewer and smaller tubers. Yield per stand falls, and total yield per hectare usually falls too — this illustrates why correct spacing, not maximum density, maximises yield.
\end{corrige}

\begin{exercice}[Seed rate calculation]
A farmer wants to plant 1 hectare of maize at $75 \times 25$ cm spacing (1 seed per stand). The seed lot has a germination percentage of $85\%$ and a 1000-grain weight of $320$ g. Calculate: (a) the number of stands per hectare; (b) the number of seeds to sow per hectare; (c) the weight of seed required in kg.
\end{exercice}

\begin{corrige}[Exercise 2]
(a) Stands per hectare $= \dfrac{10\,000\ \mathrm{m^2}}{0.75 \times 0.25\ \mathrm{m^2}} = \dfrac{10\,000}{0.1875} = 53\,333$ stands/ha.

(b) Seeds to sow $= \dfrac{53\,333}{0.85} = 62\,745$ seeds/ha.

(c) Weight $= \dfrac{62\,745 \times 320\ \mathrm{g}}{1000} = 20.08$ kg $\approx 20.1$ kg/ha.
\end{corrige}

\begin{exercice}[Integrated pest management]
A cocoa farmer in the Centre Region observes black pod disease on 15\% of pods in July. He has been spraying a copper fungicide every 3 weeks since May. (a) Why might the disease still be spreading despite spraying? (b) Propose an IPM strategy to reduce black pod incidence for the next season, listing at least four cultural/biological measures and one chemical adjustment.
\end{exercice}

\begin{corrige}[Exercise 3]
(a) Possible reasons: (i) poor spray coverage (dense canopy, no pruning); (ii) wrong timing (spraying after infection established); (iii) resistance development (same mode of action used repeatedly); (iv) high inoculum from mummified pods on trees/ground; (v) excessive shade/humidity.

(b) IPM strategy for next season:
\begin{enumerate}
\item \textbf{Sanitation}: complete removal and burial/burning of all mummified pods (source of inoculum) before flowering.
\item \textbf{Pruning}: structural pruning after main harvest to open canopy, improve air flow and spray penetration.
\item \textbf{Shade management}: thin permanent shade trees to reduce humidity below 80\% in canopy.
\item \textbf{Resistant clones}: gradually replant gaps with SNKA 12/13/14 (black pod tolerant).
\item \textbf{Chemical adjustment}: rotate copper (M1) with fosetyl-Al (33) or metalaxyl-M (4) + mancozeb (M3) to manage resistance; apply at pod set (first cherelles) and repeat at 3-week intervals during high humidity; use motorised mist-blower for better coverage.
\end{enumerate}
\end{corrige}

\coursec{Animal Nutrition and Animal Production}

In the previous section we studied how crops are cultivated to feed people and to supply raw materials. But crops also feed animals, and animals in turn supply meat, milk, eggs, hides, draught power and manure. From the cattle herds of the Adamawa Plateau to the poultry farms around Bafoussam and the pigsties of the West Region, animal production is a pillar of the Cameroonian rural economy. Just as a crop only yields well when the soil supplies the right nutrients, a farm animal only produces well when its diet supplies the right nutrients in the right proportions. This section therefore answers two linked questions: \emph{what do farm animals need to eat?} and \emph{how do we organise their production efficiently and profitably?}

\courssub{Classes of Nutrients Required by Farm Animals}

Observe a goat tethered by the roadside in Bamenda. It eats grass all day, drinks water, and yet grows, produces milk and kids once a year. Hidden inside that grass are six classes of \cle{nutrients}, each with a distinct role.

\begin{definition}[Nutrients required by farm animals]
A \cle{nutrient} is any chemical substance in feed that is used by the animal's body for energy, growth, repair or regulation. Farm animals require six classes:
\begin{itemize}
\item \textbf{Carbohydrates}: the main \cle{energy} source (starches, sugars, cellulose). Energy fuels movement, body heat, milk synthesis and egg formation. Excess is stored as fat.
\item \textbf{Proteins}: supply \cle{amino acids} for growth of muscle, wool, feathers, milk and eggs, and for enzymes and hormones. Essential for young, pregnant and lactating animals. Ten amino acids are essential for non-ruminants; ruminants can synthesise them via rumen microbes if nitrogen is available.
\item \textbf{Lipids (fats and oils)}: concentrated energy (about $2.25$ times the energy of carbohydrates per gram), carriers of fat-soluble vitamins (A, D, E, K), and sources of essential fatty acids (linoleic, linolenic).
\item \textbf{Vitamins}: needed in tiny amounts for regulation; fat-soluble (A, D, E, K) and water-soluble (B-complex, C). Examples: vitamin A (vision, reproduction, epithelial integrity), vitamin D (calcium/phosphorus metabolism, bone formation), vitamin E (antioxidant), B-vitamins (coenzymes in metabolism).
\item \textbf{Minerals}: macro-minerals (calcium, phosphorus, sodium, chlorine, magnesium, potassium, sulphur) and micro-minerals (iron, copper, zinc, manganese, iodine, cobalt, selenium). Calcium and phosphorus for bones, eggshells and milk; sodium and chlorine (common salt) for osmotic balance; iron and copper for haemoglobin; iodine for thyroid hormones.
\item \textbf{Water}: the most abundant and most neglected nutrient. A lactating cow may drink $40$--$60$ litres per day; without water, feed intake and production collapse within hours. Water regulates temperature, transports nutrients, and participates in hydrolysis reactions.
\end{itemize}
\end{definition}

\begin{attention}[Water is a nutrient]
Many farmers in the dry season of the Far North provide feed but forget water. A cow can survive weeks without feed but only a few days without water. Always provide clean water \emph{ad libitum} (freely available). Water quality matters: high salinity or microbial contamination reduces intake and causes disease.
\end{attention}

\courssub{Feedstuffs: Roughages, Concentrates, Supplements and Additives}

Feeds are grouped according to their fibre content and nutrient density.

\begin{definition}[Roughage and concentrate]
A \cle{roughage} (or forage) is a bulky feed with high crude fibre content (more than about $18\%$ of dry matter) and low digestible energy per kilogram: grasses (e.g. \emph{Panicum maximum}, \emph{Brachiaria}), legume fodder (e.g. \emph{Stylosanthes}, \emph{Calliandra}, \emph{Leucaena}), hay (dried forage) and silage (forage fermented and stored in the absence of air).
A \cle{concentrate} is a feed with low fibre and high nutrient density, fed in small quantities: energy concentrates (maize, sorghum, cassava flour, wheat bran, rice bran) and protein concentrates (cottonseed cake, groundnut cake, soya bean meal, fish meal, bone meal, blood meal, brewer's dried grains).
\end{definition}

\begin{definition}[Ration and balanced diet]
A \cle{ration} is the total quantity of feed offered to an animal in 24 hours. A \cle{balanced ration} (balanced diet) is a ration that supplies all the nutrients required by the animal, in the correct amounts and proportions, for its physiological state (maintenance, growth, pregnancy, lactation or egg production), at minimum cost. It must be palatable, non-toxic, and provide adequate dry matter intake.
\end{definition}

\textbf{Supplements and additives.} A \cle{supplement} corrects a specific deficiency: mineral licks (salt, calcium, phosphorus) for grazing cattle, or bone meal / oyster shell added to poultry feed for eggshell formation. \cle{Additives} are non-nutritive substances added in small amounts: vitamin premixes, antibiotics or coccidiostats (disease prevention), enzymes (phytase, xylanase) and growth promoters. Their misuse (e.g. antibiotic residues in meat, antimicrobial resistance) is a public health concern; withdrawal periods must be respected.

\begin{savaistu}[Local feed resources in Cameroon]
Cameroon has abundant agro-industrial by-products for animal feed: brewer's dried grains (from breweries in Douala and Yaoundé), cottonseed cake (from SODECOTON in the North), palm kernel cake (from palm oil mills in the South-West and Littoral), wheat bran (from flour mills), and cassava peels (from processing units). Research at IRAD (Institute of Agricultural Research for Development) evaluates their nutritive value and inclusion levels.
\end{savaistu}

\courssub{Digestion in Ruminants versus Non-Ruminants}

Why can a cow thrive on grass alone while a pig or a chicken cannot? The answer lies in the stomach.

\begin{propriete}[Microbial fermentation in the rumen]
Ruminants (cattle, goats, sheep) possess a four-chambered stomach. The largest chamber, the \cle{rumen}, hosts billions of micro-organisms (bacteria, protozoa, fungi) which ferment cellulose and other fibres, producing \cle{volatile fatty acids} (VFAs: acetate, propionate, butyrate) that the animal absorbs through the rumen wall and uses for energy. The microbes also synthesise microbial protein (from non-protein nitrogen like urea) and B-vitamins, which the animal digests later in the abomasum and small intestine. Non-ruminants (pigs, poultry, humans) lack a functional rumen and \textbf{cannot digest cellulose efficiently}; they therefore require concentrated, low-fibre, high-quality feeds.
\end{propriete}

\begin{popfigure}
\begin{tikzpicture}[scale=0.95, every node/.style={font=\small}]
% Oesophagus
\draw[thick, popink] (0.4,3.4) -- (1.6,2.6);
\node[popink, anchor=south] at (0.4,3.4) {Oesophagus};
% Rumen (big sac)
\draw[thick, popblue, fill=popblueL] (1.2,2.4) ellipse (2.4 and 1.5);
\node[popblue] at (1.0,2.4) {\textbf{RUMEN}};
\node[popblue, font=\scriptsize] at (1.0,1.9) {Fermentation vat};
% Reticulum
\draw[thick, popgreen, fill=popgreenL] (3.6,2.6) circle (0.75);
\node[popgreen, font=\scriptsize] at (3.6,2.6) {Reticulum};
% Omasum
\draw[thick, poporange, fill=popgold!30] (4.9,2.4) circle (0.7);
\node[poporange, font=\scriptsize] at (4.9,2.4) {Omasum};
% Abomasum
\draw[thick, poppurple, fill=poppink!25] (6.3,1.5) ellipse (1.3 and 0.7);
\node[poppurple, font=\scriptsize] at (6.3,1.5) {Abomasum};
% Small intestine
\draw[thick, popdark] (7.6,1.5) -- (8.6,1.5) -- (8.6,0.6);
\node[popdark, anchor=west, font=\scriptsize] at (8.7,0.6) {Small intestine};
% Flow arrows
\draw[->, thick, popdark] (1.6,2.6) -- (2.9,2.6);
\draw[->, thick, popdark] (4.3,2.5) -- (4.2,2.4);
\draw[->, thick, popdark] (5.6,2.2) -- (5.9,1.9);
% Regurgitation arrow
\draw[->, thick, dashed, poppink] (2.2,3.6) to[bend left=30] (0.9,3.3);
\node[poppink, font=\scriptsize, align=center, text width=2.6cm] at (1.6,4.2) {Regurgitation\\(chewing the cud)};
\end{tikzpicture}
\legende{The four chambers of the ruminant stomach and the flow of digesta: oesophagus $\rightarrow$ rumen $\rightarrow$ reticulum $\rightarrow$ omasum $\rightarrow$ abomasum (the ``true stomach'' with gastric juice) $\rightarrow$ small intestine.}
\end{popfigure}

Feed passes from the oesophagus into the \textbf{rumen} and \textbf{reticulum} (often called the reticulorumen), where microbes ferment fibre; coarse material is regurgitated and re-chewed (rumination). It then flows to the \textbf{omasum} (water absorption, many folds) and the \textbf{abomasum}, the true stomach where acid (HCl) and enzymes (pepsin) act as in non-ruminants.

\textbf{Practical implications.} Cattle and goats can be raised on pasture and crop residues (maize stover, groundnut haulms) with only small concentrate supplements. Pigs and poultry must be fed concentrates rich in energy and good-quality protein, because they cannot extract enough from fibre.

\begin{attention}[Do not overfeed concentrates to ruminants]
A common and costly mistake: suddenly giving cattle or goats large amounts of maize or bran. Rapid fermentation produces excess gas trapped in the rumen (\cle{bloat}) and excess lactic acid (\cle{ruminal acidosis}), both potentially fatal. Always introduce concentrates gradually (over 2--3 weeks) and ensure roughage forms at least about $40\%$ of the ration dry matter to keep the rumen functioning.
\end{attention}

\courssub{Nutritional Requirements by Physiological State}

An animal's needs change with what its body is doing. Every ration is built in two layers: a \cle{maintenance ration} (to keep the animal alive, warm and moving without gaining or losing weight) plus a \cle{production ration} added on top.

\begin{center}
\begin{tabularx}{\linewidth}{|l|X|}
\hline
\rowcolor{popblueL} \textbf{Physiological state} & \textbf{Extra requirements above maintenance} \\
\hline
Maintenance (adult, idle) & Energy and protein just to maintain body weight and temperature \\
\hline
Growth (young animals) & High protein (essential amino acids) and minerals (Ca, P) for muscle and bone \\
\hline
Pregnancy (last third) & Extra energy, protein and minerals for the foetus; rapid foetal growth in last trimester \\
\hline
Lactation & Very high demand: milk exports protein, fat, lactose, calcium and water daily; peak yield at 4--8 weeks post-partum \\
\hline
Egg production (layers) & High calcium (shells $\approx$ 2 g Ca/egg), protein and energy; about $110$--$120$ g feed/bird/day for modern hybrids \\
\hline
\end{tabularx}
\end{center}

\begin{exemplebox}[Nutrient partitioning in a dairy cow]
A 500 kg Holstein cow producing 25 kg milk/day (3.5\% fat, 3.2\% protein) requires approximately: 160 MJ metabolisable energy, 1.8 kg crude protein, 90 g calcium, 65 g phosphorus. Maintenance alone needs ~60 MJ ME and 600 g CP. The production ration must supply the difference. Underfeeding in early lactation causes excessive body condition loss and ketosis.
\end{exemplebox}

\courssub{Ration Formulation: the Pearson Square}

To formulate a balanced ration we compare the \cle{dry matter} (DM, feed minus water), \cle{crude protein} (CP) and energy content of available feeds, then mix them in calculated proportions. The simplest tool is the \cle{Pearson square}, used to blend two feeds to a target protein percentage. It assumes linear mixing and is valid for one nutrient (usually CP) at a time; energy and minerals must be checked separately.

\begin{methode}[Using the Pearson square]
\begin{enumerate}
\item Write the target CP percentage (on dry matter basis) in the centre of a square.
\item Write the CP of the two feeds at the left corners (high-protein feed top, low-protein feed bottom).
\item Subtract diagonally (larger minus smaller) and write the differences at the right corners.
\item The right-hand values give the parts of each feed in the mixture: the value opposite a feed is the proportion of that feed.
\item Convert parts to percentages of the total.
\item Verify by calculating the weighted average CP of the mixture.
\end{enumerate}
\end{methode}

\begin{exoresolu}[Formulating a 16\% CP broiler ration]
A farmer near Dschang wants to mix maize ($9\%$ CP) and soya bean meal ($44\%$ CP) to obtain a $16\%$ CP ration. In what proportions should he mix them?
\tcblower
\textbf{Pearson square:} parts of maize $= 44 - 16 = 28$; parts of soya bean meal $= 16 - 9 = 7$. Total parts $= 28 + 7 = 35$.
\begin{align*}
\text{Maize} &= \frac{28}{35}\times 100 = 80\% \\
\text{Soya bean meal} &= \frac{7}{35}\times 100 = 20\%
\end{align*}
\textbf{Check:} $(80\times 9 + 20\times 44)/100 = (720 + 880)/100 = 16\%$ CP. The farmer mixes $80$ kg of maize with $20$ kg of soya bean meal for every $100$ kg of feed (then adds minerals and vitamins).
\end{exoresolu}

\begin{popfigure}
\begin{tikzpicture}[scale=1.0, every node/.style={font=\small}]
\node[popblue] (a) at (0,2) {Soya 44\%};
\node[popblue] (b) at (0,0) {Maize 9\%};
\node[popdark, font=\large\bfseries] (c) at (2,1) {16\%};
\node[poporange] (d) at (4,2) {7 parts soya};
\node[poporange] (e) at (4,0) {28 parts maize};
\draw[thick, popmuted] (a) -- (e);
\draw[thick, popmuted] (b) -- (d);
\draw[thick, popmuted] (a) -- (d);
\draw[thick, popmuted] (b) -- (e);
\node[popmuted, font=\scriptsize] at (2,-0.9) {Subtract diagonally: $44-16=28$ ; $16-9=7$};
\end{tikzpicture}
\legende{Pearson square for a 16\% crude protein ration: 7 parts soya bean meal to 28 parts maize, i.e. 20\% soya and 80\% maize.}
\end{popfigure}

\begin{exoresolu}[Formulating a 17\% CP dairy supplement with three ingredients]
A farmer wants a 17\% CP supplement using cottonseed cake (41\% CP), maize (9\% CP) and wheat bran (15\% CP). Use Pearson square twice: first combine maize and wheat bran to make a 12\% CP energy mix, then combine that with cottonseed cake.
\tcblower
\textbf{Step 1:} Maize (9\%) + Wheat bran (15\%) $\rightarrow$ target 12\%. Parts: maize = 15-12 = 3; wheat bran = 12-9 = 3. Mix = 50\% maize, 50\% wheat bran. CP = 12\%.
\textbf{Step 2:} Energy mix (12\%) + Cottonseed cake (41\%) $\rightarrow$ target 17\%. Parts: energy mix = 41-17 = 24; cottonseed cake = 17-12 = 5. Total = 29.
Final mix: Energy mix = 24/29 $\approx$ 82.8\%; Cottonseed cake = 5/29 $\approx$ 17.2\%.
Within energy mix: Maize = 41.4\%, Wheat bran = 41.4\%, Cottonseed cake = 17.2\%.
Check: $(41.4\times9 + 41.4\times15 + 17.2\times41)/100 = 17\%$.
\end{exoresolu}

\courssub{Systems of Animal Production}

\begin{center}
\begin{tabularx}{\linewidth}{|l|X|X|}
\hline
\rowcolor{popblueL} \textbf{System} & \textbf{Features} & \textbf{Advantages / limitations in Cameroon} \\
\hline
Extensive & Free range, pastoralism, transhumance; animals find their own feed; low inputs; indigenous breeds & Low cost; uses marginal land; but low productivity, high disease exposure, farmer--grazer conflicts (North West, Adamawa), seasonal feed shortage \\
\hline
Semi-intensive & Animals graze by day, housed at night; some supplementary feeding and health care; crossbred or improved breeds & Good compromise for village poultry and goats; moderate cost and output; manure collected for crops \\
\hline
Intensive & Animals permanently housed; all feed, water and health care provided; high stocking density; exotic/hybrid breeds & High output per animal (broilers, layers, pigs); efficient feed conversion; but high capital, feed cost (70--80\% of total cost) and disease risk; requires reliable electricity and water \\
\hline
\end{tabularx}
\end{center}

\begin{savaistu}[Transhumance in the Adamawa]
The Adamawa Plateau hosts the largest cattle population in Cameroon. Pastoralists (Mbororo) practice transhumance: moving herds between dry-season lowlands (river valleys) and wet-season highlands. This traditional system is under pressure from crop expansion, climate variability, and land tenure insecurity. Improved corridors and water points are policy priorities.
\end{savaistu}

\courssub{Housing and Equipment}

Good housing protects animals from predators, rain, heat and thieves, and makes feeding, cleaning and disease control easy.

\begin{itemize}
\item \textbf{Poultry house}: orientation east--west to avoid direct sun; open sides with wire mesh for \cle{ventilation}; deep litter (wood shavings, sawdust, rice husks) kept dry; density about $8$--$10$ broilers or $4$--$5$ layers per m$^2$; perches (for layers), nest boxes (1 per 4--5 hens), feeders and drinkers (1 per 50 birds for tube feeders, 1 per 80 for bell drinkers).
\item \textbf{Pigsty}: concrete floor sloping to a drain for easy washing; separate pens for boar, pregnant sows, farrowing sows (with a creep area protecting piglets from crushing) and fatteners; shade and wallow or sprinklers, because pigs cannot sweat and are heat-stressed above $25^\circ$C.
\item \textbf{Cattle shed}: roofed, well-drained floor (earth or concrete with grooves), feeding trough and water trough, milking area; space about $3.5$--$4$ m$^2$ per adult cow; calf pens separate.
\end{itemize}

\begin{propriete}[Hygiene and ventilation]
Damp, dirty, poorly ventilated housing concentrates ammonia and pathogens and is the single greatest cause of respiratory disease and parasite build-up in intensive systems. Dry litter, regular removal of dung, disinfection and free air movement are non-negotiable. (Not examinable as a proof, but frequently examined as short-answer reasoning.)
\end{propriete}

\begin{popfigure}
\begin{tikzpicture}[scale=0.8, every node/.style={font=\small}]
% Poultry house plan
\draw[thick, popdark] (0,0) rectangle (10,6);
% Door
\draw[thick, popblue, fill=popblueL] (0,2.5) rectangle (1,3.5);
\node[popblue, font=\scriptsize, rotate=90] at (0.5,3) {Door + Footbath};
% Ventilation openings
\foreach \x in {1.5,3.5,5.5,7.5,9.5} {
  \draw[thick, popgreen, fill=popgreenL] (\x,5.8) rectangle (\x+0.5,6);
  \draw[thick, popgreen, fill=popgreenL] (\x,0) rectangle (\x+0.5,0.2);
}
\node[popgreen, font=\scriptsize] at (5,6.3) {Wire mesh ventilation};
% Feeders
\foreach \x in {2,5,8} {
  \draw[thick, poporange, fill=popgold!30] (\x,1.5) rectangle (\x+1,2);
  \node[poporange, font=\scriptsize] at (\x+0.5,1.75) {Feeder};
}
% Drinkers
\foreach \x in {2,5,8} {
  \draw[thick, popblue, fill=popblueL] (\x,3.5) rectangle (\x+1,4);
  \node[popblue, font=\scriptsize] at (\x+0.5,3.75) {Drinker};
}
% Orientation
\node[poppurple, font=\scriptsize] at (5,-0.5) {East $\leftarrow$ \textbf{Length} $\rightarrow$ West};
\node[poppurple, font=\scriptsize, rotate=90] at (-0.5,3) {North $\uparrow$ \textbf{Width} $\downarrow$ South};
% Dimensions
\node[popdark, font=\scriptsize] at (5,6.5) {Length $\approx$ 10 m, Width $\approx$ 6 m for 500 broilers (10 birds/m$^2$ = 50 m$^2$)};
\end{tikzpicture}
\legende{Plan of a poultry house for 500 broilers: east--west orientation, wire-mesh ventilation on both long sides, evenly spaced feeders and drinkers, door with footbath. Deep litter floor (wood shavings).}
\end{popfigure}

\courssub{Breeding and Reproduction}

Improvement begins with \cle{selection} of breeding stock: choose parents with good growth rate, milk or egg yield, disease resistance and body conformation, from reliable records. \textbf{Mating systems} include natural mating (one bull or boar for a group of females) and \cle{artificial insemination} (AI), in which semen from a superior male is collected, stored (fresh or frozen) and deposited in the female by a technician; AI spreads superior genes cheaply and avoids keeping dangerous males. \cle{Crossbreeding} (e.g. Holstein $\times$ Gudali for dairy, Large White $\times$ local pig for growth) exploits heterosis (hybrid vigour).

Gestation lengths to remember: cow about $283$ days, goat and ewe about $150$ days, sow about $114$ days (``3 months, 3 weeks, 3 days''), hen incubation $21$ days. Care of the young includes colostrum (first milk, rich in antibodies) within the first $6$ hours, warmth (brooder for chicks, heat lamp for piglets), and creep feed for early weaning.

\begin{exemplebox}[Artificial insemination in Cameroon]
AI is promoted by MINEPIA (Ministry of Livestock, Fisheries and Animal Industries) for dairy improvement. The Centre for Artificial Insemination (CIAC) in Ngaoundéré produces and distributes semen from improved dairy breeds (Holstein, Montbéliarde) and beef breeds. Farmers in the West and North West use AI to upgrade local cattle without keeping a bull.
\end{exemplebox}

\courssub{Animal Health: Diseases, Parasites and Prevention}

\begin{center}
\begin{tabularx}{\linewidth}{|l|X|X|}
\hline
\rowcolor{popblueL} \textbf{Animal} & \textbf{Common diseases / parasites} & \textbf{Prevention} \\
\hline
Poultry & Newcastle disease, fowl pox, coccidiosis, Gumboro (IBD); lice, worms (Ascaridia) & Vaccination (Newcastle: Hitchner B1 at day 7, La Sota at day 21; Gumboro at day 14), coccidiostats in feed, clean dry litter, deworming \\
\hline
Cattle & Trypanosomiasis (tsetse fly), tick-borne diseases (anaplasmosis, babesiosis, heartwater), foot-and-mouth disease (FMD); ticks, worms, liver fluke & Trypanocidal drugs (prophylactic/isometamidium), dipping/spraying against ticks (acaricides), strategic deworming (albendazole, ivermectin), FMD vaccination in risk zones \\
\hline
Pigs & African swine fever (ASF), worms, mange (Sarcoptes) & Strict biosecurity (fencing, no swill feeding without boiling 30 min), deworming, acaricide spray for mange; no vaccine for ASF \\
\hline
Goats/sheep & Peste des petits ruminants (PPR), worms (Haemonchus), mange, foot rot & PPR vaccination (at 3 months, then annually), deworming (FAMACHA-guided for Haemonchus), hoof trimming, clean pens \\
\hline
\end{tabularx}
\end{center}

Prevention rests on three pillars: \cle{vaccination} (following a calendar, e.g. Newcastle vaccine in chicks), regular \cle{deworming} and tick control, and \cle{biosecurity} (quarantine of new animals for 21 days, footbaths at entrances, restricting visitors, isolating sick animals, proper carcass disposal).

\begin{attention}[Exam tip]
GCE questions on animal production very often ask for a \textbf{labelled housing plan} (poultry house or pigsty) or a \textbf{vaccination schedule} presented as a table. Practise drawing a simple, clearly labelled plan and tabulating: age of animal -- vaccine -- route of administration (e.g. eye drop, drinking water, intramuscular).
\end{attention}

\courssub{Records, Economics and Improvement}

Without records, a farmer cannot know whether the enterprise is making a profit. Essential \cle{farm records} include: inventory records (stock of animals, feed, equipment), production records (milk per cow per day, eggs per day per hen house), breeding records (service dates, birth dates, parentage), health records (treatments, vaccinations, mortality), and financial records (variable costs: feed, drugs, vaccines; fixed costs: housing, equipment; revenue: sales of animals, products, manure).

Key \cle{production indices}:
\begin{align*}
\text{Average daily gain (ADG)} &= \frac{\text{final weight} - \text{initial weight}}{\text{number of days}} \quad (\text{kg/day or g/day}) \\
\text{Feed conversion ratio (FCR)} &= \frac{\text{feed consumed (kg)}}{\text{weight gained (kg)}} \quad (\text{dimensionless; lower is better}) \\
\text{Egg production (\%)} &= \frac{\text{eggs collected per day}}{\text{number of hens present}}\times 100 \\
\text{Mortality rate (\%)} &= \frac{\text{number of deaths}}{\text{number at start}}\times 100
\end{align*}

\begin{exoresolu}[Calculating ADG and FCR for broilers]
A batch of broilers weighs $0.2$ kg at 2 weeks and $2.0$ kg at 8 weeks (42 days between weighings). Over this period each bird consumed $3.36$ kg of feed. Calculate the ADG and the FCR.
\tcblower
\begin{align*}
\text{ADG} &= \frac{2.0 - 0.2}{42} = \frac{1.8}{42} \approx 0.043\ \text{kg/day} = 43\ \text{g/day} \\
\text{FCR} &= \frac{3.36}{1.8} \approx 1.87
\end{align*}
An FCR of $1.87$ means $1.87$ kg of feed produced $1$ kg of live weight gain — a good performance for broilers (lower FCR = better efficiency). Target for modern broilers: FCR $< 1.7$ at 35 days.
\end{exoresolu}

\begin{popfigure}
\begin{tikzpicture}
\begin{axis}[
width=11cm, height=6.5cm,
xlabel={Age (weeks)}, ylabel={Live weight (kg)},
xmin=0, xmax=9, ymin=0, ymax=2.8,
xtick={0,1,2,3,4,5,6,7,8},
grid=both, grid style={popline!40},
legend style={at={(0.03,0.97)}, anchor=north west, font=\scriptsize},
]
\addplot[smooth, very thick, popblue] coordinates {(0,0.04)(1,0.15)(2,0.35)(3,0.65)(4,1.0)(5,1.4)(6,1.8)(7,2.2)(8,2.5)};
\addlegendentry{Good quality feed (16--21\% CP)}
\addplot[smooth, very thick, poporange, dashed] coordinates {(0,0.04)(1,0.12)(2,0.25)(3,0.45)(4,0.7)(5,0.95)(6,1.2)(7,1.45)(8,1.65)};
\addlegendentry{Poor quality feed (low protein)}
\end{axis}
\end{tikzpicture}
\legende{Growth curves of broiler chickens: live weight versus age. A balanced ration (high crude protein and energy) gives faster growth and earlier market weight; poor feed delays slaughter age and raises costs.}
\end{popfigure}

\textbf{Marketing.} Profit depends as much on selling as on producing. Farmers should time sales to periods of high demand (end-of-year festivities, Ramadan for sheep and goats, Easter for poultry), add value where possible (eggs graded and packed, chickens dressed and chilled), and keep transport costs low by grouping sales through cooperatives or farmer associations. Knowledge of live weight vs. carcass weight (dressing percentage: broilers $\approx 70\%$, pigs $\approx 75\%$, cattle $\approx 50\%$) helps negotiate prices.

\begin{aretenir}[Key points]
\begin{itemize}
\item Six nutrient classes: carbohydrates, proteins, lipids, vitamins, minerals, water.
\item Roughages (fibre-rich) versus concentrates (nutrient-dense); ruminants digest fibre thanks to rumen microbes, non-ruminants cannot.
\item Rations are built for maintenance plus production; the Pearson square balances two feeds to a target protein level (check energy and minerals separately).
\item Production systems range from extensive (pastoralism) to intensive (housed, high input); each has trade-offs in Cameroon.
\item Success rests on housing hygiene, selection and breeding (AI, crossbreeding), vaccination and deworming, and record-keeping with indices such as ADG, FCR and egg production percentage.
\end{itemize}
\end{aretenir}

\begin{exercice}[Practice]
\begin{enumerate}
\item A farmer mixes cottonseed cake ($41\%$ CP) with maize ($9\%$ CP) to obtain a $17\%$ CP supplement for dairy cows. Using the Pearson square, find the proportions of each feed.
\item A layer flock of $200$ hens produces $150$ eggs per day. Calculate the percentage egg production and comment.
\item Draw a labelled plan of a poultry house for $500$ broilers, indicating orientation, ventilation openings, feeders and drinkers.
\item Explain why a sudden change from pasture to a maize-rich ration can kill a goat.
\item A piglet weighs $1.5$ kg at weaning (28 days) and $30$ kg at 112 days. It consumed $85$ kg of feed over this period. Calculate ADG and FCR for the growing-finishing phase.
\end{enumerate}
\end{exercice}

\begin{corrige}[Exercise 1]
\begin{enumerate}
\item Pearson square: parts of cottonseed cake $= 17 - 9 = 8$; parts of maize $= 41 - 17 = 24$; total $= 32$. Cottonseed cake $= 8/32 \times 100 = 25\%$; maize $= 24/32 \times 100 = 75\%$. Check: $(25\times 41 + 75\times 9)/100 = (1025 + 675)/100 = 17\%$.
\item Egg production $= 150/200 \times 100 = 75\%$. This is a good laying rate for a well-managed flock at peak production (target $> 80\%$ for modern hybrids at peak).
\item The plan should show: east--west orientation, wire-mesh side walls for ventilation, roof overhang, deep litter floor, feeders and drinkers distributed evenly, a door and a footbath at the entrance; stocking density about $10$ birds/m$^2$ gives $50$ m$^2$ of floor space (e.g. $10 \times 5$ m).
\item Sudden excess of maize causes rapid microbial fermentation in the rumen: gas accumulates (bloat) and lactic acid lowers rumen pH (acidosis), damaging the rumen wall and killing rumen microbes; both conditions can be fatal. Concentrates must be introduced gradually with adequate roughage.
\item ADG $= (30 - 1.5) / (112 - 28) = 28.5 / 84 \approx 0.339$ kg/day $= 339$ g/day. FCR $= 85 / 28.5 \approx 2.98$. This FCR is acceptable for growing-finishing pigs (target $< 3.0$).
\end{enumerate}
\end{corrige}

\newpage

\coursec{Integration activity}

\coursec{Agriculture}

\courssub{Integration Activity: The Nkeng Family Mixed Farm Project}

\begin{aretenir}[Aims of the Activity]
This integration activity brings together all concepts from Chapter BIO08 (Lessons 88--92). By the end, you should be able to:
\begin{itemize}
\item assess the properties of a given soil and propose scientifically justified measures to improve its \cle{fertility};
\item design a rational \cle{land-use plan} including \cle{crop rotation} and soil conservation on sloping land;
\item construct a \cle{cropping calendar} adapted to the bimodal rainfall regime of the Cameroonian highlands;
\item apply the \cle{Pearson square} method to formulate a balanced poultry ration from local feedstuffs;
\item plan \cle{housing}, a \cle{vaccination} and \cle{deworming} schedule, and a \cle{record book} for small livestock;
\item integrate plant and animal production into a sustainable \cle{mixed farming} system and defend your choices orally.
\end{itemize}
\end{aretenir}

\begin{attention}[Skills Assessed]
This activity mobilises the GCE Advanced Level competences of \textbf{analysing data}, \textbf{applying knowledge to a novel situation}, \textbf{carrying out calculations} and \textbf{communicating scientific arguments}. These are exactly the skills tested in Paper 2 (structured questions) and Paper 3 (practical/problem-solving).
\end{attention}

\courssub{Situation}

\begin{experience}[The Nkeng Family Farm Project]
The Nkeng family lives near \textbf{Dschang} in the West Region of Cameroon, at an altitude of about \SI{1400}{m}. The climate is of the highland type, with a long rainy season from \textbf{mid-March to mid-November} (peak rainfall July--September) and a dry season from \textbf{mid-November to mid-March}. Mean temperatures range from \SIrange{18}{24}{\ensuremath{{}^{\circ}\mathrm{C}}}.

The family owns \textbf{2 hectares} of gently sloping land. A soil survey and simple field tests gave the following results:

\begin{center}
\begin{tabularx}{\linewidth}{|l|X|}
\hline
\rowcolor{popblueL} \textbf{Property} & \textbf{Observation} \\
\hline
Texture & Lateritic sandy loam; reddish colour \\
\hline
pH (water) & 5.2 (acidic) \\
\hline
Organic matter & Low; topsoil only about \SI{10}{cm} deep \\
\hline
Drainage & Good on the slope; a small marshy patch at the bottom \\
\hline
Erosion & Rills visible after heavy storms on the upper slope \\
\hline
\end{tabularx}
\end{center}

The family wishes to produce, for home consumption and for sale at the Dschang market:
\begin{itemize}
\item \textbf{maize} and \textbf{cassava} (staple food crops);
\item \textbf{vegetables} (cabbage, carrots, tomatoes) on the more fertile lowland patch;
\item \textbf{100 broilers} and \textbf{50 layers} per production cycle;
\item \textbf{2 dairy goats} for milk.
\end{itemize}

Available local feedstuffs and their approximate crude protein (CP) contents:
\begin{center}
\begin{tabularx}{\linewidth}{|l|X|}
\hline
\rowcolor{popblueL} \textbf{Feedstuff} & \textbf{Crude protein (\%)} \\
\hline
Maize grain (energy feed) & 9 \\
\hline
Cottonseed cake (protein concentrate) & 41 \\
\hline
\end{tabularx}
\end{center}

A broiler finisher ration should contain about \textbf{20\,\% CP}. The family has access to poultry manure, kitchen waste, and can grow \emph{Calliandra} and leucaena as fodder shrubs along the boundaries.
\end{experience}

\begin{savaistu}[Cameroon Context]
The West Region (Dschang, Bafoussam) is a major agricultural zone with volcanic and lateritic soils. The bimodal rainfall pattern (though here described as a long single season at 1400 m) allows two maize crops per year. \emph{Calliandra calothyrsus} and \emph{Leucaena leucocephala} are widely promoted by IRAD and NGOs for alley cropping and fodder banks because they fix nitrogen, tolerate acid soils, and provide high-protein browse for ruminants.
\end{savaistu}

\courssub{Tasks}

Work in groups of four to six. Answer the following tasks in order; each one builds on the previous.

\begin{enumerate}
\item \textbf{Soil assessment.} From the soil data, identify \textbf{two strengths} and \textbf{three limitations} of this soil for crop production. For each limitation, propose one corrective measure and justify it biologically or chemically.

\item \textbf{Land-use plan.} Draw a sketch plan of the 2 ha farm, dividing it into at least four plots. Position the crops, the poultry house, the goat shed and the fodder shrubs. Justify the position of each element (think of slope, drainage, wind, hygiene and labour).

\item \textbf{Crop rotation.} Propose a \textbf{three-year rotation} for the main upland plot, including a legume. Explain two advantages of your rotation.

\item \textbf{Cropping calendar.} Using the rainfall data, construct a seasonal calendar (table) showing land preparation, planting, weeding/fertiliser application and harvesting for maize, cassava and one vegetable.

\item \textbf{Feed formulation.} Using the \textbf{Pearson square}, calculate the proportions of maize grain and cottonseed cake needed to prepare \SI{100}{kg} of a broiler finisher ration at 20\,\% CP. Show all working.

\item \textbf{Animal husbandry plan.} For the poultry and the goats, propose:
\begin{enumerate}
\item one housing requirement and the reason for it;
\item a simple health calendar (vaccination of chicks against Newcastle disease; deworming of goats at the start of the rains);
\item three records the family should keep, and why.
\end{enumerate}

\item \textbf{Integration and sustainability.} Show with a simple flow diagram how wastes and products circulate between the crop and animal components (e.g.\ manure $\rightarrow$ crops; crop residues $\rightarrow$ animals). State one reason why this integration improves sustainability.

\item \textbf{Presentation.} Prepare a 10-minute group presentation defending your plan. The class will vote for the most \textbf{sustainable and profitable} design.
\end{enumerate}

\courssub{Model Answers}

\textbf{Task 1 — Soil Assessment.}\par
\textbf{Strengths:}
\begin{enumerate}
\item The sandy loam texture is \cle{well aerated} and easy to work, suitable for root crops like cassava which require loose soil for tuber expansion.
\item Drainage is good on the slope, preventing waterlogging of maize roots (maize is sensitive to anaerobic conditions).
\end{enumerate}

\textbf{Limitations and Corrections:}
\begin{itemize}
\item \emph{Acidity (pH 5.2):} At this pH, aluminium and manganese become soluble and toxic to roots, while phosphorus availability is sharply reduced due to fixation by iron/aluminium oxides. \textbf{Correction:} Apply agricultural lime (\ce{CaCO3}) at a rate determined by a lime requirement test (typically \SIrange{2}{4}{t/ha} for such soils) to raise the pH towards 6.0--6.5. Liming also supplies \ce{Ca^{2+}} and \ce{Mg^{2+}}, improves microbial activity, and enhances P availability.
\item \emph{Low organic matter:} Poor soil structure, low water-holding capacity, low cation exchange capacity (CEC), and limited nutrient supply. \textbf{Correction:} Incorporate compost, farmyard manure from the poultry and goats, and practise green manuring with a legume (e.g., \emph{Mucuna} or \emph{Crotalaria}) during the rotation. This builds humus, improves aggregation, and increases CEC.
\item \emph{Erosion on the slope:} Loss of the thin fertile topsoil (only \SI{10}{cm} deep) by surface run-off during heavy storms. \textbf{Correction:} Construct contour ridges (tied ridges) and plant grass strips (\emph{Pennisetum purpureum}) or fodder shrubs (\emph{Calliandra}) along contour lines to slow run-off, increase infiltration, and trap sediment.
\end{itemize}

\textbf{Task 2 — Land-Use Plan.}\par

\begin{popfigure}
\begin{tikzpicture}[scale=1.0, font=\small]
\draw[thick, popdark] (0,0) rectangle (10,6);
% Plot A - upper slope
\draw[fill=popgreenL] (0,3.2) rectangle (4.8,6);
\node[align=center, text width=4.2cm] at (2.4,4.6) {Plot A (upper slope, $\sim$1 ha)\\ Maize / legume rotation\\ Contour ridges + tied ridges};
% Plot B - mid slope
\draw[fill=popblueL] (5.2,3.2) rectangle (10,6);
\node[align=center, text width=4.2cm] at (7.6,4.6) {Plot B (mid slope, $\sim$1 ha)\\ Cassava + fodder shrubs\\ (\emph{Calliandra}) on contour lines};
% Plot C - lowland
\draw[fill=popgreenL!70!popblueL] (0,0) rectangle (4.8,2.8);
\node[align=center, text width=4.2cm] at (2.4,1.4) {Plot C (lowland, moist, $\sim$0.5 ha)\\ Vegetables: cabbage, carrots, tomatoes\\ Raised beds + organic manure};
% Farmstead - well drained, near homestead
\draw[fill=popblueL!60] (5.2,0) rectangle (10,2.8);
\node[align=center, text width=4.2cm] at (7.6,1.9) {Poultry house (100 broilers + 50 layers)\\ Goat shed (2 dairy goats)\\ East--west orientation, raised floor};
\node[align=center, text width=4.2cm] at (7.6,0.7) {Compost pit \\ Manure storage};
% Slope arrow
\draw[->, thick, poporange] (0.3,5.8) -- (9.7,0.3) node[midway, above, sloped, font=\scriptsize]{Slope direction (upper $\to$ lower)};
% Wind direction indicator (optional)
\draw[->, thick, poppurple] (10.2,3) -- (10.2,5) node[midway, right, font=\scriptsize]{Prevailing wind};
\end{tikzpicture}
\legende{Model land-use plan for the 2 ha mixed farm near Dschang (not to scale).}
\end{popfigure}

\textbf{Justifications:}
\begin{itemize}
\item \textbf{Buildings (poultry house, goat shed, compost pit)} are placed on well-drained ground (mid-slope, away from the marshy patch) near the homestead for easy surveillance, security, and to minimise labour for carrying feed and manure. The east--west orientation reduces direct solar radiation on the long walls, preventing overheating. They are positioned \emph{downwind} of the vegetable plot (prevailing wind from the west in this region) to avoid contamination of crops by dust and odours.
\item \textbf{Vegetables (Plot C)} occupy the moist, relatively fertile lowland patch where water is accessible for irrigation in the dry season. Raised beds improve drainage for cabbage and carrots.
\item \textbf{Maize/Legume rotation (Plot A, upper slope)} benefits from contour ridges to control erosion on the steepest part.
\item \textbf{Cassava + fodder shrubs (Plot B, mid slope):} Cassava tolerates poorer soils and drought; \emph{Calliandra} hedgerows on the contour serve simultaneously as \cle{erosion barriers} and as \cle{protein-rich fodder} for the goats (cut-and-carry system).
\end{itemize}

\textbf{Task 3 — Three-Year Rotation (Plot A).}\par
\textbf{Year 1:} Maize (heavy feeder, high N demand) $\rightarrow$ \textbf{Year 2:} Groundnuts (\emph{Arachis hypogaea}) or cowpea (\emph{Vigna unguiculata}) (legume, N-fixing) $\rightarrow$ \textbf{Year 3:} Cassava (tolerates lower fertility, deep roots) or a green manure fallow (\emph{Mucuna pruriens}).

\textbf{Advantages:}
\begin{enumerate}
\item The legume fixes atmospheric nitrogen through \emph{Rhizobium} bacteria in its root nodules, adding $50$--$100\ \mathrm{kg\ N\ ha^{-1}}$ to the system and reducing the need for mineral nitrogen fertiliser for the subsequent crop.
\item Rotating crops with different botanical families, rooting depths, and pest hosts breaks \cle{pest and disease cycles} (e.g., maize streak virus, cassava mosaic virus, nematodes) and exploits nutrients from different soil layers (shallow maize roots vs. deep cassava roots).
\end{enumerate}

\textbf{Task 4 — Cropping Calendar.}\par
\begin{center}
\begin{tabularx}{\linewidth}{|l|X|X|X|}
\hline
\rowcolor{popblueL} \textbf{Crop} & \textbf{Land Preparation} & \textbf{Planting} & \textbf{Maintenance / Harvest} \\
\hline
Maize (1st season) & Feb -- early Mar: clearing, burning (if needed), contour ridging, basal fertiliser (NPK 20-10-10) & Mid-Mar (onset of rains): 2 seeds/hole, 80$\times$40 cm & Weed at 3 and 6 weeks; top-dress urea at 6 weeks (50 kg/ha); harvest Jul--Aug (green maize) or Aug--Sep (dry grain) \\
\hline
Maize (2nd season) & Late Aug: minimal tillage on residues & Early Sep & Weed + top-dress; harvest Dec--Jan \\
\hline
Cassava & Feb--Mar: mounds or ridges on contour & Mar--Apr: stem cuttings 25 cm, 1$\times$1 m & Weed at 1, 3, 6 months; harvest 10--12 months later (Jan--Mar Year 2) \\
\hline
Cabbage (Plot C) & Mar: raised beds, incorporate poultry manure (10 t/ha) & Nursery: Mar; Transplant: Apr (45$\times$45 cm) & Regular watering; side-dress urea at 3 weeks; harvest Jun--Jul. Dry-season crop: nursery Oct, transplant Nov, harvest Feb--Mar (with irrigation) \\
\hline
\end{tabularx}
\end{center}

\textbf{Task 5 — Pearson Square for Broiler Finisher Ration.}\par
\begin{methode}[Applying the Pearson Square]
\begin{enumerate}
\item Draw a square. Write the CP content of each feedstuff at the left corners (Maize = 9, Cottonseed cake = 41) and the target CP (20) in the centre.
\item Subtract diagonally (always larger minus smaller):
\begin{itemize}
\item $41 - 20 = 21$ parts of maize (energy feed)
\item $20 - 9 = 11$ parts of cottonseed cake (protein concentrate)
\end{itemize}
\item Total parts $= 21 + 11 = 32$.
\item Calculate the weight of each ingredient for \SI{100}{kg} of ration:
\begin{align*}
\text{Maize} &= \frac{21}{32} \times \SI{100}{kg} = \SI{65.625}{kg} \approx \SI{65.6}{kg} \\
\text{Cottonseed cake} &= \frac{11}{32} \times \SI{100}{kg} = \SI{34.375}{kg} \approx \SI{34.4}{kg}
\end{align*}
\end{enumerate}
\end{methode}

\textbf{Verification:}
\[
(\SI{65.6}{kg} \times 0.09) + (\SI{34.4}{kg} \times 0.41) = \SI{5.904}{kg} + \SI{14.104}{kg} = \SI{20.008}{kg} \text{ CP per } \SI{100}{kg} \text{ ration} = 20.0\,\% \text{ CP.} \quad \checkmark
\]

\begin{attention}[Practical Note]
This ration provides energy and protein only. A complete commercial broiler finisher ration must also include: a vitamin-mineral premix (\SIrange{0.3}{0.5}{\%}), salt (\SI{0.25}{\%}), lysine/methionine if needed, and a coccidiostat. The cottonseed cake must be from a source with low gossypol (or treated) to avoid toxicity.
\end{attention}

\textbf{Task 6 — Animal Husbandry Plan.}\par
\begin{itemize}
\item \textbf{Housing:}
\begin{itemize}
\item \emph{Poultry:} The house must be \textbf{well ventilated but draught-free}, with a deep litter of dry wood shavings or rice husks (\SI{10}{cm} deep) and an east--west orientation. \textbf{Reason:} Damp litter favours coccidiosis (sporulation of oocysts) and respiratory diseases (ammonia from uric acid decomposition); good ventilation removes moisture and ammonia while preventing cold draughts that stress chicks.
\item \emph{Goats:} The shed must be \textbf{raised \SI{60}{cm} above ground} with a slatted floor (bamboo or timber, gaps \SI{1.5}{cm}) and a roof overhang of \SI{50}{cm}. \textbf{Reason:} Droppings fall through, keeping the floor dry and clean; this drastically reduces re-infection by gastrointestinal nematodes (worm larvae develop in moist faeces on the ground) and prevents foot rot.
\end{itemize}
\item \textbf{Health Calendar:}
\begin{itemize}
\item \emph{Poultry (Newcastle disease):} Vaccinate day-old chicks (intra-ocular or spray), booster at 4 weeks (drinking water or eye drop), then every 3--4 months for layers. Use a thermostable vaccine (e.g., I-2) if cold chain is unreliable.
\item \emph{Goats (Strategic deworming):} Treat with a broad-spectrum anthelmintic (e.g., albendazole or levamisole) at the \textbf{start of the rains} (March) when infective larvae on pasture surge, and again at \textbf{mid-rains} (July). Rotate drug classes annually to delay resistance. Check faecal egg counts if possible.
\item \emph{General:} Isolate and treat any sick animal immediately; disinfect houses between batches; control rodents and wild birds.
\end{itemize}
\item \textbf{Records (three essential ones):}
\begin{enumerate}
\item \textbf{Production records:} Daily egg count per layer batch; weekly broiler weights; daily milk yield per goat. \emph{Why:} Early detection of disease, nutritional deficiency, or management problems; basis for culling decisions.
\item \textbf{Health records:} Dates of vaccination, deworming, treatments (drug, dose, withdrawal period), mortalities with suspected cause. \emph{Why:} Ensures timely boosters, avoids drug residues in meat/milk, tracks disease patterns.
\item \textbf{Financial records:} Costs of feed, chicks, drugs, equipment; revenue from sales of eggs, broilers, milk, manure, crops. \emph{Why:} Calculates gross margin, identifies profitable enterprises, supports loan applications.
\end{enumerate}
\end{itemize}

\textbf{Task 7 — Integration Diagram.}\par
\begin{popfigure}
\begin{tikzpicture}[font=\small, node distance=2.5cm]
\node[draw, rounded corners, fill=popgreenL, align=center, text width=3.5cm] (crops) at (0,0) {CROP COMPONENT\\ Maize, cassava, vegetables\\ Fodder shrubs (\emph{Calliandra})};
\node[draw, rounded corners, fill=popblueL, align=center, text width=3.5cm] (animals) at (8,0) {ANIMAL COMPONENT\\ 100 broilers + 50 layers\\ 2 dairy goats};
\node[draw, rounded corners, fill=popgold, align=center, text width=3.5cm] (soil) at (4,-2.5) {SOIL FERTILITY\\ Organic matter, nutrients, structure};
\draw[->, thick, poporange] (crops) to[bend left=15] node[above, align=center]{Crop residues (stover, peelings)\\ Fodder (shrubs, legume haulms)\\ Grain (maize for feed)} (animals);
\draw[->, thick, popgreen] (animals) to[bend left=15] node[below, align=center]{Manure (poultry litter, goat droppings)\\ Urine (N, K)\\ Compost} (crops);
\draw[->, thick, popgreen] (animals) -- node[right, align=center]{Manure \\ Compost} (soil);
\draw[->, thick, popgreen] (soil) -- node[left, align=center]{Nutrients \\ Water \\ Support} (crops);
\draw[->, thick, poppurple] (crops) to[bend right=20] node[below, align=center]{Human food \\ Income} (animals);
\draw[->, thick, poppurple] (animals) to[bend right=20] node[above, align=center]{Meat, eggs, milk \\ Income \\ Draught power (optional)} (crops);
\end{tikzpicture}
\legende{Nutrient and energy flows in the Nkeng family mixed farming system.}
\end{popfigure}

\textbf{Sustainability Statement:} Recycling manure and compost as fertiliser reduces the need to purchase mineral fertiliser, maintains soil organic matter and cation exchange capacity, and closes the nutrient loop. This makes the farm productive without degrading the soil resource, embodying the principle of \cle{sustainable intensification}.

\begin{exercice}[Extension Task]
The family later considers adding a fish pond on the marshy patch. Suggest one advantage and one risk of this addition, and state how pond mud could be used in the farm's nutrient cycle.
\end{exercice}

\begin{corrige}[Extension Task]
\emph{Advantage:} The pond converts an unproductive marshy patch into a source of high-quality animal protein (fish, e.g., \emph{Oreochromis niloticus}) and income. Pond water can also be used for supplementary irrigation of vegetables during dry spells, increasing crop yields.

\emph{Risk:} Stagnant water can favour breeding of \emph{Anopheles} mosquitoes (malaria vectors) and \emph{Simulium} (onchocerciasis); predators (birds, snakes, monitor lizards) may cause losses; poor embankment construction can lead to flooding and escape of fish. Proper design (deep water, steep banks, predator nets) and management (larvivorous fish, regular monitoring) are essential.

\emph{Pond mud:} During pond draining (every 1--2 years), the accumulated mud is rich in nitrogen, phosphorus, and organic matter from fish waste and uneaten feed. It can be spread on vegetable beds (Plot C) or composted with crop residues as an excellent organic fertiliser, further closing the nutrient cycle and reducing external input dependence.
\end{corrige}

\newpage

\coursec{Methods and worked examples}

\courssub{Method 1: Identifying a Soil Type from its Physical Properties}

\begin{methode}[How to Identify and Classify a Soil]
To determine the nature of a soil sample in the field or in the laboratory, proceed in order:
\begin{enumerate}
\item \textbf{Observe the colour}: dark soils are usually rich in \cle{humus} (decomposed organic matter); red or yellow colours indicate iron oxides (\cle{lateritic soils}, very common in the forest zone of Cameroon); grey-blue colours suggest poor drainage and waterlogging (\cle{gley soils}).
\item \textbf{Feel the texture} (the ``sausage test''): moisten a small sample and try to roll it into a cylinder. If it cannot be rolled, the soil is \cle{sandy}; if it rolls but cracks, it is \cle{loamy}; if it rolls into a long flexible sausage that can be bent into a ring without breaking, it is \cle{clayey}.
\item \textbf{Estimate particle sizes by sedimentation}: shake the soil thoroughly with water in a measuring cylinder and leave it to settle. \textbf{Sand} settles first (bottom layer), then \textbf{silt}, then \textbf{clay} (top mineral layer); humus floats or remains suspended. Measure the thickness of each layer and express it as a percentage of the total mineral depth.
\item \textbf{Test drainage and water retention}: pour a known volume of water through a known mass of soil and measure the volume collected. Sandy soils drain fast and retain little water; clay soils retain much water but may become waterlogged and poorly aerated.
\item \textbf{Conclude} by matching your results to the properties of the main soil types: sandy, clay, loamy, humus-rich, lateritic, alluvial. A \cle{loam} (a roughly balanced mixture of sand, silt and clay, with humus) is generally the best agricultural soil.
\end{enumerate}
\textbf{Reflex:} always link the soil type to its consequences for agriculture: aeration, drainage, water retention, ease of cultivation (tilth), and suitability for particular crops.
\end{methode}

\begin{exoresolu}[Identifying Soil Samples from a School Farm]
A student in Buea collected three soil samples (A, B, C). The results of simple tests are summarised below:

\begin{center}
\begin{tabularx}{\linewidth}{|l|X|X|X|}
\hline
\rowcolor{popblueL} \textbf{Test} & \textbf{Sample A} & \textbf{Sample B} & \textbf{Sample C} \\
\hline
Sausage test & Cannot be rolled & Rolls into a long flexible ring & Rolls but cracks easily \\
\hline
Water drained through 100 g of soil (out of 50 mL poured) & 45 mL & 5 mL & 25 mL \\
\hline
Colour & Pale yellow & Reddish & Dark brown \\
\hline
\end{tabularx}
\end{center}

Identify each soil type and state which is most suitable for growing maize, justifying your answer.
\tcblower
\textbf{Step 1 — Analyse Sample A.} It cannot be rolled into a sausage, so its particles do not stick together: they are coarse. 45 mL out of 50 mL of water drained through, i.e. very rapid drainage and poor water retention. The pale colour indicates little humus. \textbf{Sample A is a sandy soil.}

\textbf{Step 2 — Analyse Sample B.} It forms a long flexible ring, showing very fine, sticky particles. Only 5 mL of water passed through: drainage is very poor and water retention is very high, so the soil risks waterlogging and poor root aeration. The reddish colour indicates iron oxides. \textbf{Sample B is a clay soil (lateritic clay).}

\textbf{Step 3 — Analyse Sample C.} It rolls but cracks: an intermediate texture, i.e. a balanced mixture of sand, silt and clay. Drainage is moderate (25 mL out of 50 mL), so it retains enough water without waterlogging. The dark brown colour shows it is rich in humus. \textbf{Sample C is a loamy soil (humus-rich loam).}

\textbf{Step 4 — Conclusion.} \textbf{Sample C (loam) is the most suitable for maize} because it combines good drainage and aeration (roots respire easily) with adequate water and nutrient retention; it is easy to till, and its humus releases mineral salts as it decomposes. Sample A would need frequent irrigation and manuring; Sample B would be heavy to work and poorly aerated.
\end{exoresolu}

\courssub{Method 2: Diagnosing Plant Mineral Deficiencies}

\begin{methode}[How to Diagnose a Mineral Deficiency in a Crop Plant]
\begin{enumerate}
\item \textbf{Locate the symptoms}: do they appear first on \emph{old (lower) leaves} or on \emph{young (upper) leaves}? Mobile elements (N, P, K, Mg) are withdrawn from old leaves and moved to young ones when scarce, so their deficiency shows on \textbf{old leaves first}. Immobile elements (Ca, Fe) cannot be relocated, so their deficiency shows on \textbf{young leaves first}.
\item \textbf{Describe the symptom precisely}: uniform yellowing (\cle{chlorosis}), yellowing between green veins (\cle{interveinal chlorosis}), purple coloration, scorched leaf edges, stunted growth, poor roots, etc.
\item \textbf{Match the symptom to the element}:
\begin{itemize}
\item \textbf{Nitrogen (N)}: uniform yellowing of old leaves, stunted growth (N is needed for proteins and chlorophyll).
\item \textbf{Phosphorus (P)}: dark green or purplish leaves, poor root development, poor seed/fruit formation.
\item \textbf{Potassium (K)}: yellowing then scorching (browning) of the margins of old leaves; weak stems.
\item \textbf{Magnesium (Mg)}: interveinal chlorosis of old leaves (Mg is the central atom of chlorophyll).
\item \textbf{Calcium (Ca)}: death of growing points, distorted young leaves, blossom-end rot of fruits.
\item \textbf{Iron (Fe)}: interveinal chlorosis of young leaves.
\end{itemize}
\item \textbf{Propose a remedy}: apply the appropriate fertiliser (e.g. NPK 20-10-10, urea for N, superphosphate for P, muriate of potash for K), compost or farmyard manure, and correct the soil pH (liming) if necessary.
\end{enumerate}
\textbf{Reflex:} in the exam, always justify the diagnosis with TWO arguments: (i) where the symptom appears and (ii) its precise form.
\end{methode}

\begin{exoresolu}[Deficiency in a Tomato Field]
A farmer near Bafoussam observes that the \emph{older, lower} leaves of his tomato plants turn yellow while the veins remain green; the young leaves at the top look normal. The soil is slightly acidic.
\begin{enumerate}
\item[a)] Identify the deficient element, justifying your answer.
\item[b)] Suggest two corrective measures.
\end{enumerate}
\tcblower
\textbf{a) Diagnosis.} Two observations must be combined.

\textbf{(i) Location of symptoms.} The symptoms appear on the \emph{old leaves first}, while young leaves remain healthy. The element concerned is therefore \textbf{mobile} in the plant: when it becomes scarce, the plant withdraws it from old leaves and transports it to the growing points. This eliminates calcium and iron, which are immobile and whose deficiency appears on young leaves.

\textbf{(ii) Form of the symptom.} The yellowing is \emph{interveinal}: the tissue between the veins turns yellow while the veins stay green. This is the classic sign of \textbf{magnesium deficiency}, because magnesium is the central atom of the chlorophyll molecule; without it, chlorophyll cannot be synthesised, while the veins (mostly vascular tissue with little chlorophyll) remain green.

\textbf{Conclusion:} the crop suffers from \textbf{magnesium (Mg) deficiency}.

\textbf{b) Corrective measures.}
\begin{itemize}
\item Apply a magnesium-containing fertiliser, e.g. magnesium sulphate (Epsom salt) or dolomitic limestone, which supplies both Mg and Ca.
\item Correct the soil acidity by \textbf{liming}: in acidic soils, Mg$^{2+}$ ions are easily leached and their uptake is hindered by excess H$^{+}$ ions; adding lime (CaCO$_{3}$) raises the pH towards neutrality and improves Mg availability.
\item (Also acceptable: adding well-decomposed compost or farmyard manure, which releases Mg and other minerals slowly and improves soil structure.)
\end{itemize}
\end{exoresolu}

\courssub{Method 3: Calculating Fertiliser Requirements (NPK)}

\begin{methode}[How to Calculate the Quantity of Fertiliser to Apply]
\begin{enumerate}
\item \textbf{Read the fertiliser grade}: a bag labelled NPK 15-15-15 contains 15 \% nitrogen (N), 15 \% phosphorus pentoxide (P$_{2}$O$_{5}$) and 15 \% potassium oxide (K$_{2}$O) \emph{by mass}.
\item \textbf{Write the key formula}:
\[ \text{Mass of fertiliser} = \frac{\text{mass of nutrient required}}{\text{percentage of nutrient in fertiliser}} \times 100 \]
\item \textbf{Scale to the field area}: nutrient requirement per hectare $\times$ number of hectares.
\item \textbf{Check units} throughout (kg, ha) and round sensibly to whole bags (usually 50 kg bags).
\end{enumerate}
\textbf{Reflex:} the percentage refers to the \emph{nutrient}, not the whole fertiliser. Never multiply when you should divide: the fertiliser mass must be \emph{larger} than the nutrient mass.
\end{methode}

\begin{exoresolu}[Fertilising a Maize Field]
An agricultural extension officer recommends that a maize crop receives 90 kg of nitrogen (N) per hectare. The farmer has two fertilisers available: urea (46 \% N) and NPK 20-10-10 (20 \% N). His field covers 2.5 ha.
\begin{enumerate}
\item[a)] Calculate the mass of urea needed for the whole field.
\item[b)] Calculate the mass of NPK 20-10-10 that would be needed to supply the same amount of nitrogen.
\item[c)] How many 50 kg bags of urea must the farmer buy?
\end{enumerate}
\tcblower
\textbf{a) Mass of urea.}

Nitrogen required per hectare: 90 kg. For 2.5 ha:
\[ 90 \times 2.5 = 225 \ \text{kg of N required.} \]

Urea contains 46 \% N, i.e. 46 kg of N per 100 kg of urea. Therefore:
\[ \text{Mass of urea} = \frac{225}{46} \times 100 = 489.1 \ \text{kg} \approx 489 \ \text{kg of urea.} \]

\textbf{b) Mass of NPK 20-10-10.}

This fertiliser contains only 20 \% N:
\[ \text{Mass of NPK} = \frac{225}{20} \times 100 = 1125 \ \text{kg of NPK 20-10-10.} \]

\textbf{Comment:} much more NPK is needed because it is less concentrated in nitrogen; however, it also supplies phosphorus and potassium, which urea does not. The choice depends on the soil analysis.

\textbf{c) Number of 50 kg bags of urea.}
\[ \frac{489.1}{50} = 9.78 \ \text{bags.} \]
Since bags cannot be bought in fractions, the farmer must buy \textbf{10 bags of 50 kg} of urea.

\textbf{Check:} 10 bags = 500 kg of urea, supplying $500 \times 0.46 = 230$ kg of N, which slightly exceeds the 225 kg required — acceptable.
\end{exoresolu}

\courssub{Method 4: Calculating Plant Density and Expected Yield}

\begin{methode}[How to Use Spacing to Compute Plant Population and Yield]
\begin{enumerate}
\item \textbf{Convert all measurements to the same unit} (metres and hectares: 1 ha $= 10\,000\ \mathrm{m^2}$).
\item \textbf{Area occupied by one plant} = row spacing $\times$ spacing within the row.
\item \textbf{Plant population}:
\[ \text{Population} = \frac{\text{area of field}}{\text{area per plant}} \]
\item \textbf{Expected yield} = population $\times$ mean yield per plant (or per stand, if several seeds are planted per hole).
\item \textbf{Discuss}: closer spacing increases population but increases competition for light, water and nutrients; wider spacing wastes land. Correct spacing is a key \cle{cultivation technique}.
\end{enumerate}
\end{methode}

\begin{exoresolu}[Groundnut Field in the North Region]
A farmer plants groundnuts at a spacing of 50 cm between rows and 25 cm within the row, on a field of 1.5 ha. Each stand produces on average 40 g of dry pods.
\begin{enumerate}
\item[a)] Calculate the plant population of the field.
\item[b)] Estimate the total expected harvest in kg and in 100 kg bags.
\item[c)] State two advantages of correct spacing.
\end{enumerate}
\tcblower
\textbf{a) Plant population.}

Convert spacings to metres: 50 cm $= 0.5$ m; 25 cm $= 0.25$ m.

Area per plant:
\[ 0.5 \times 0.25 = 0.125 \ \mathrm{m^2}. \]

Area of field:
\[ 1.5 \ \text{ha} = 1.5 \times 10\,000 = 15\,000 \ \mathrm{m^2}. \]

Population:
\[ \frac{15\,000}{0.125} = 120\,000 \ \text{plants.} \]

\textbf{b) Expected harvest.}

Total yield:
\[ 120\,000 \times 40 \ \text{g} = 4\,800\,000 \ \text{g} = 4800 \ \text{kg.} \]

Number of 100 kg bags:
\[ \frac{4800}{100} = 48 \ \text{bags.} \]

\textbf{c) Advantages of correct spacing.}
\begin{itemize}
\item Each plant receives adequate light, water and mineral salts, reducing competition and giving uniform, high yields.
\item It facilitates weeding, pest control and harvesting (passage between rows), and ensures good aeration, which limits fungal diseases.
\end{itemize}
\end{exoresolu}

\courssub{Method 5: Choosing and Justifying Cultivation Techniques}

\begin{methode}[How to Answer a Question on Crop Cultivation Techniques]
For any named crop (food crop: maize, cassava, plantain; cash/non-food crop: cocoa, oil palm, cotton, rubber), structure your answer along the production cycle:
\begin{enumerate}
\item \textbf{Choice of site and land preparation}: suitable soil and climate; clearing, ploughing/tilling, ridging or mounding.
\item \textbf{Planting material}: improved/certified seeds, cuttings (cassava), suckers (plantain), budded/grafted seedlings (cocoa, rubber); treatment of seeds against pests.
\item \textbf{Planting}: correct season (onset of the rains), depth, spacing.
\item \textbf{Maintenance (cultural practices)}: weeding, mulching, thinning, staking, pruning, earthing-up, irrigation/drainage, fertiliser/manure application, pest and disease control (integrated pest management).
\item \textbf{Harvesting and post-harvest}: at correct maturity; drying, storage in cribs/barns, processing.
\item \textbf{For perennial cash crops}: add the nursery stage, shading of young plants, and the immature period before production (e.g. cocoa 3--5 years, oil palm about 3--4 years, rubber about 6--7 years).
\end{enumerate}
\textbf{Reflex:} always link each technique to its \emph{purpose} (e.g. ``mulching conserves soil moisture and suppresses weeds'').
\end{methode}

\begin{exoresolu}[Comparing Cassava and Cocoa Cultivation]
\begin{enumerate}
\item[a)] Describe the propagation and planting of cassava.
\item[b)] Give three ways in which the cultivation of cocoa differs from that of cassava, explaining each difference.
\end{enumerate}
\tcblower
\textbf{a) Propagation and planting of cassava.}

Cassava is propagated \textbf{vegetatively} by stem cuttings. Mature stems (from plants 8--12 months old) are cut into pieces about 20--25 cm long, each bearing several nodes. The land is cleared and tilled, often into ridges or mounds. Cuttings are planted at the onset of the rains, inserted vertically or slanting with about two-thirds of their length buried, at a spacing of about 1 m $\times$ 1 m (roughly 10 000 stands/ha). The field is weeded regularly (2--3 times) and the tubers are ready for harvest 9--12 months after planting.

\textbf{b) Differences with cocoa cultivation.}
\begin{enumerate}
\item \textbf{Mode of propagation:} cocoa is grown from \emph{seeds} first raised in a nursery (or from budded/grafted seedlings), whereas cassava is planted directly as stem cuttings. Reason: cocoa seeds lose viability quickly and young seedlings are delicate, so they need nursery protection before transplanting.
\item \textbf{Life cycle and waiting period:} cocoa is a \emph{perennial tree crop} that starts bearing pods only after 3--5 years and produces for decades; cassava is an \emph{annual/semi-perennial food crop} harvested after less than a year. Cocoa therefore requires long-term investment and permanent management of the plantation.
\item \textbf{Shade requirement:} young cocoa plants require \emph{shade} (provided by plantain/banana or shade trees) because direct strong sunlight damages them; cassava is a sun-loving crop planted in open fields.
\end{enumerate}
(Also acceptable: cocoa requires a humid equatorial climate with well-distributed rainfall, while cassava tolerates poorer soils and drier conditions.)
\end{exoresolu}

\courssub{Method 6: Formulating a Balanced Animal Ration (Pearson's Square)}

\begin{methode}[How to Balance a Ration Using Pearson's Square]
To mix two feedstuffs (an energy feed and a protein supplement) to obtain a desired protein percentage:
\begin{enumerate}
\item Write the \textbf{target protein \%} in the centre of a square.
\item Write the protein \% of the \textbf{energy feed} (top left) and of the \textbf{protein supplement} (bottom left).
\item Subtract \emph{diagonally} (larger minus smaller) to obtain the \textbf{parts} of each feed on the right side.
\item Convert parts to percentages of the total, then to masses for the quantity of feed required.
\item \textbf{Check}: the protein contributed by each component must sum to the target.
\end{enumerate}
\textbf{Reflex:} an animal ration must supply energy (carbohydrates/fats), protein, minerals, vitamins and water, adjusted to the animal's age, mass and production (growth, lactation, egg-laying). Pearson's square balances \emph{one} nutrient (usually crude protein).
\end{methode}

\begin{exoresolu}[Ration for Broiler Chickens]
A poultry farmer in Bamenda wants to prepare 100 kg of a broiler feed containing 20 \% crude protein (CP) by mixing maize (9 \% CP) and soybean meal (45 \% CP).
\begin{enumerate}
\item[a)] Use Pearson's square to find the proportions of maize and soybean meal.
\item[b)] Calculate the mass of each ingredient in the 100 kg mixture.
\item[c)] Verify your answer.
\end{enumerate}
\tcblower
\textbf{a) Pearson's square.}

\begin{center}
\begin{tikzpicture}[scale=1.0]
\node (m) at (0,1) {Maize 9\%};
\node (s) at (0,-1) {Soybean 45\%};
\node (t) at (2.2,0) {\textbf{20\%}};
\node (pm) at (4.8,1) {$45-20=25$ parts maize};
\node (ps) at (4.8,-1) {$20-9=11$ parts soybean};
\draw[thick,popblue] (m.east) -- (ps.west);
\draw[thick,popblue] (s.east) -- (pm.west);
\end{tikzpicture}
\end{center}

Total parts: $25 + 11 = 36$ parts.

Proportions:
\[ \text{Maize} = \frac{25}{36} \times 100 = 69.4\ \% \qquad \text{Soybean meal} = \frac{11}{36} \times 100 = 30.6\ \%. \]

\textbf{b) Masses for 100 kg of feed.}
\[ \text{Maize} = \frac{25}{36} \times 100 = 69.4 \ \text{kg}; \qquad \text{Soybean meal} = \frac{11}{36} \times 100 = 30.6 \ \text{kg}. \]

\textbf{c) Verification.}

Protein supplied by maize: $69.4 \times 0.09 = 6.25$ kg.

Protein supplied by soybean: $30.6 \times 0.45 = 13.77$ kg.

Total protein: $6.25 + 13.77 = 20.02 \approx 20$ kg per 100 kg of feed, i.e. \textbf{20 \% CP}. The mixture is correct.

\textbf{Practical note:} in a real ration, a mineral--vitamin premix, salt and a source of calcium (e.g. oyster shell) would also be added, and the maize/soybean proportions slightly reduced to accommodate them.
\end{exoresolu}

\courssub{Method 7: Evaluating Animal Production Performance}

\begin{methode}[How to Calculate Growth Rate and Feed Conversion Ratio]
\begin{enumerate}
\item \textbf{Average daily gain (ADG)}:
\[ \text{ADG} = \frac{\text{final mass} - \text{initial mass}}{\text{number of days}} \]
\item \textbf{Feed conversion ratio (FCR)}:
\[ \text{FCR} = \frac{\text{total feed consumed (kg)}}{\text{total weight gain (kg)}} \]
A \emph{low} FCR means efficient conversion (broilers typically around 1.6--2; pigs around 3; cattle considerably higher).
\item \textbf{Cost of production per kg of gain} = FCR $\times$ cost of 1 kg of feed.
\item \textbf{Interpret and advise}: compare with standard values; suggest improvements (balanced ration, disease control, improved breeds, good housing and hygiene).
\end{enumerate}
\textbf{Reflex:} in production questions, always end with an \emph{interpretation} — a number alone earns only part of the marks.
\end{methode}

\begin{exoresolu}[Performance of a Piggery]
A farmer fattens a batch of pigs. At the start, the average mass per pig is 25 kg; after 90 days the average mass is 85 kg. Each pig consumed on average 180 kg of feed costing 250 FCFA per kg.
\begin{enumerate}
\item[a)] Calculate the average daily gain (ADG).
\item[b)] Calculate the feed conversion ratio (FCR) and comment.
\item[c)] Calculate the feed cost per kg of weight gain.
\item[d)] Suggest two ways of improving the FCR.
\end{enumerate}
\tcblower
\textbf{a) Average daily gain.}
\[ \text{ADG} = \frac{85 - 25}{90} = \frac{60}{90} = 0.67 \ \text{kg/day} \quad (\approx 670 \ \text{g/day}). \]

\textbf{b) Feed conversion ratio.}
\[ \text{FCR} = \frac{\text{feed consumed}}{\text{weight gain}} = \frac{180}{60} = 3.0. \]
\textbf{Comment:} the pig eats 3 kg of feed to produce 1 kg of weight gain. This is a typical, acceptable value for pigs, though well-managed improved breeds on balanced rations can do slightly better (around 2.5--2.8).

\textbf{c) Feed cost per kg of gain.}
\[ \text{Cost} = \text{FCR} \times \text{price per kg of feed} = 3.0 \times 250 = 750 \ \text{FCFA per kg of weight gain.} \]
The farmer can compare this with the selling price of pork per kg live weight to judge profitability.

\textbf{d) Improving the FCR.}
\begin{itemize}
\item Provide a \textbf{balanced ration} with adequate protein, energy, minerals and vitamins (e.g. correctly formulated using Pearson's square), and clean water \emph{ad libitum}.
\item Improve \textbf{health and husbandry}: regular deworming and vaccination, clean dry housing, reduced stress; and/or use \textbf{improved breeds} with better genetic potential for growth.
\end{itemize}
\end{exoresolu}

\courssub{Method 8: Interpreting an Animal Production System}

\begin{methode}[How to Analyse an Animal Production System (Extensive vs Intensive)]
\begin{enumerate}
\item \textbf{Identify the system} from the description: \cle{extensive} (large area, animals graze freely, low inputs, low output per animal — e.g. traditional cattle herding on the Adamawa Plateau) or \cle{intensive} (animals housed/confined, high inputs of feed, veterinary care and labour, high output per animal — e.g. modern poultry or piggery units near Yaoundé and Douala). \cle{Semi-intensive} systems are intermediate.
\item \textbf{List inputs and outputs}: land, labour, capital, feed, veterinary care $\rightarrow$ meat, milk, eggs, hides, manure.
\item \textbf{Compare} in a table: stocking rate, productivity per animal, disease risk, cost, environmental impact.
\item \textbf{Conclude with a justified judgement}: which system suits which context (land availability, capital, market).
\end{enumerate}
\end{methode}

\begin{exoresolu}[Choosing a Production System for Dairy Cattle]
A group of farmers on the Bamenda Highlands wants to increase milk production. They currently keep local zebu cattle on open range (extensive system), obtaining about 2 litres of milk per cow per day. An adviser proposes an intensive zero-grazing unit with improved (exotic) dairy cows.
\begin{enumerate}
\item[a)] Give three advantages the intensive system could bring.
\item[b)] Give three constraints or risks the farmers must consider.
\item[c)] Propose one intermediate (semi-intensive) measure.
\end{enumerate}
\tcblower
\textbf{a) Advantages of the intensive system.}
\begin{itemize}
\item \textbf{Higher productivity per animal}: improved dairy breeds, well fed, can produce several times more milk per day than local zebus on range.
\item \textbf{Better health control and management}: housed animals are easily inspected, treated, vaccinated and milked hygienically; breeding (e.g. artificial insemination) is controlled.
\item \textbf{Efficient land use}: zero grazing produces more milk per hectare, which is important where land is limited, and manure is easily collected for crop fields or biogas.
\end{itemize}

\textbf{b) Constraints and risks.}
\begin{itemize}
\item \textbf{High capital and running costs}: construction of housing, purchase of improved cows, concentrates and veterinary products; a market failure or disease outbreak can cause heavy losses.
\item \textbf{Feed supply must be guaranteed}: zero-grazed animals depend entirely on cut fodder and purchased feed; fodder must be grown or conserved (silage, hay) for the dry season.
\item \textbf{Exotic breeds are less hardy}: they are generally more susceptible to heat, parasites (e.g. ticks) and diseases than local zebus, requiring skilled management and reliable veterinary services.
\end{itemize}

\textbf{c) Semi-intensive measure.}

The farmers could keep their cattle in fenced paddocks or night kraals, \textbf{supplementing grazing with cultivated fodder (e.g. elephant grass, \emph{Calliandra}) and mineral licks}, while introducing improved bulls or artificial insemination to upgrade the herd gradually. This raises productivity at moderate cost and risk, allowing a progressive transition towards intensification.
\end{exoresolu}

\newpage

\coursec{Exercises}

\coursec{I check my knowledge}

\begin{exercice}[Multiple choice questions]
For each question, choose the single correct answer.
\begin{enumerate}
\item The soil horizon richest in humus and most important for crop growth is:
\begin{enumerate}
\item the A horizon \item the B horizon \item the C horizon \item the bedrock
\end{enumerate}
\item The element whose deficiency causes interveinal chlorosis in older leaves and which is a component of chlorophyll is:
\begin{enumerate}
\item nitrogen \item phosphorus \item magnesium \item potassium
\end{enumerate}
\item In the Pearson square method, a ration of 16\% crude protein is to be made from maize (9\% CP) and groundnut cake (45\% CP). The proportion of groundnut cake needed is approximately:
\begin{enumerate}
\item 19\% \item 36\% \item 64\% \item 81\%
\end{enumerate}
\item Newcastle disease in poultry is caused by:
\begin{enumerate}
\item a bacterium \item a virus \item a fungus \item a protozoan
\end{enumerate}
\end{enumerate}
\end{exercice}

\begin{exercice}[True or false — justify]
State whether each statement is true or false, and justify your answer in one or two sentences.
\begin{enumerate}
\item A clay soil is always more fertile than a sandy soil and therefore always better for farming.
\item Legumes such as groundnut and cowpea can grow well without nitrogen fertiliser because they fix atmospheric nitrogen through \emph{Rhizobium} bacteria in their root nodules.
\item Ruminants such as cattle can digest cellulose because they themselves secrete the enzyme cellulase.
\item In extensive animal production, output per animal is generally lower than in intensive production.
\end{enumerate}
\end{exercice}

\begin{exercice}[Definitions]
Define each of the following terms precisely, as expected at the GCE Advanced Level:
\begin{enumerate}
\item soil profile;
\item humus;
\item macronutrient (in plant nutrition);
\item ration (in animal nutrition);
\item biosecurity (in animal production).
\end{enumerate}
\end{exercice}

\begin{exercice}[Direct calculations]
\begin{enumerate}
\item A soil sample contains \SI{40}{\percent} sand, \SI{40}{\percent} silt and \SI{20}{\percent} clay. Calculate the sand-to-clay ratio and state, with a reason, whether this soil is likely to drain quickly or slowly.
\item A farmer applies urea (46\% N) at a rate of \SI{100}{kg} per hectare. Calculate the mass of nitrogen applied per hectare.
\item A broiler chick consumes \SI{150}{g} of feed per day containing 21\% crude protein. Calculate the mass of protein consumed per day.
\end{enumerate}
\end{exercice}

\coursec{I practise}

\begin{exercice}[Using the texture triangle]
A soil sample from a farm near Buea contains 70\% sand, 20\% silt and 10\% clay.
\begin{enumerate}
\item Use the texture triangle to identify the textural class of this soil.
\item State three physical properties of this soil type (drainage, water retention, ease of tillage).
\item Suggest two crops suited to this soil, giving a reason for each choice.
\item Propose two measures the farmer could take to improve the fertility of this soil.
\end{enumerate}
\end{exercice}

\begin{exercice}[Mineral elements in plants]
Copy and complete the following table by giving, for each mineral element, one function in the plant and one symptom of its deficiency.

\begin{center}
\begin{tabularx}{\linewidth}{|l|X|X|}
\hline
\rowcolor{popblueL} \textbf{Element} & \textbf{One function} & \textbf{One deficiency symptom} \\
\hline
Nitrogen (N) & & \\
\hline
Phosphorus (P) & & \\
\hline
Potassium (K) & & \\
\hline
Magnesium (Mg) & & \\
\hline
Iron (Fe) & & \\
\hline
\end{tabularx}
\end{center}
\end{exercice}

\begin{methode}[Balancing a ration using the Pearson square]
The Pearson square is a simple tool to calculate the proportions of two feeds needed to obtain a mixture with a desired nutrient concentration.
\begin{enumerate}
\item Write the desired nutrient percentage in the centre of the square.
\item Write the nutrient percentages of the two available feeds at the top left and bottom left corners.
\item Subtract diagonally (larger minus smaller) to get the parts of each feed required. The parts are always positive.
\item The sum of the parts gives the total parts of the mixture.
\item Calculate the percentage of each feed: (parts of feed / total parts) $\times$ 100\%.
\end{enumerate}
\end{methode}

\begin{exercice}[Balancing a ration with the Pearson square]
A poultry farmer in Bafoussam wants to prepare a broiler starter ration containing 20\% crude protein using maize (9\% crude protein) and fish meal (60\% crude protein).
\begin{enumerate}
\item Draw a Pearson square and calculate the proportions of maize and fish meal required.
\item Calculate the mass of each ingredient needed to prepare \SI{100}{kg} of the ration.
\item Comment on the result: why is fish meal used only in a limited proportion, apart from the protein requirement? (Think of cost and palatability.)
\end{enumerate}
\end{exercice}

\begin{exercice}[Cultivation of a food crop]
A farmer in the West Region wishes to establish a one-hectare maize farm.
\begin{enumerate}
\item List, in chronological order, the main cultivation operations from land preparation to storage of the harvest.
\item For each operation, state its purpose.
\item Indicate two common pests or diseases of maize in Cameroon and one control measure for each.
\end{enumerate}
\end{exercice}

\coursec{I prepare for the GCE}

\begin{exercice}[Structured question: soils and fertility (20 marks)]
The table below gives the results of the analysis of two soil samples, A and B, taken from two farms.

\begin{center}
\begin{tabularx}{\linewidth}{|l|X|X|X|X|}
\hline
\rowcolor{popblueL} \textbf{Sample} & \textbf{Sand (\%)} & \textbf{Silt (\%)} & \textbf{Clay (\%)} & \textbf{Humus content} \\
\hline
A & 85 & 10 & 5 & Low \\
\hline
B & 20 & 25 & 55 & Moderate \\
\hline
\end{tabularx}
\end{center}

\begin{enumerate}
\item Identify the textural class of each soil. \hfill (4 marks)
\item Compare the drainage, aeration and water-holding capacity of the two soils. \hfill (6 marks)
\item Explain why soil A warms up quickly in the dry season but loses nutrients easily by leaching. \hfill (4 marks)
\item For each soil, suggest one suitable crop and one method of improving its fertility, justifying your answers. \hfill (6 marks)
\end{enumerate}
\end{exercice}

\begin{exercice}[Essay: cultivation of cocoa (20 marks)]
Cocoa is one of the main export crops of Cameroon.\par
Describe the cultivation of cocoa from land preparation to the marketing of dried beans, indicating the importance of each operation. Your answer should include: choice of site and soil, land preparation, planting and shade management, maintenance operations (weeding, pruning, pest and disease control, fertiliser application), harvesting, fermentation, drying and marketing.\par
\emph{(Marks will be awarded for the correct sequence of operations, the technical accuracy of each step, and the stated importance of each operation.)}
\end{exercice}

\begin{exercice}[Problem situation and essay: animal production (20 marks)]
\textbf{Part A.} A poultry farmer near Yaoundé records high mortality among week-old chicks showing respiratory signs (gasping, nasal discharge) and twisted necks.
\begin{enumerate}
\item Identify the most likely disease and its causal agent. \hfill (2 marks)
\item Describe how the disease spreads within a flock. \hfill (3 marks)
\item Outline a vaccination and biosecurity programme that the farmer should adopt to protect the farm. \hfill (5 marks)
\end{enumerate}
\textbf{Part B.} Discuss the advantages and disadvantages of intensive and extensive systems of animal production, with reference to cattle rearing in the Adamawa Region of Cameroon. Conclude by suggesting how the two systems could be improved or combined to increase meat and milk supply in Cameroon. \hfill (10 marks)
\end{exercice}

\begin{savaistu}[Cameroon context: Agriculture]
Cameroon is often called "Africa in miniature" due to its diverse agro-ecological zones. The volcanic soils around Mount Cameroon and in the West Region (Bamenda Highlands) are highly fertile and support intensive cultivation of crops like tea, coffee, and market gardening. The Adamawa plateau is the heart of extensive cattle ranching, while the Sudano-Sahelian zone in the Far North relies on drought-tolerant crops like millet, sorghum, and cotton. Cocoa, coffee, cotton, and bananas are major export earners. The government's "Second Generation Agriculture" policy aims to modernise the sector through mechanisation, improved seeds, and value chain development.
\end{savaistu}

\begin{attention}[Common mistakes in GCE Agriculture]
\begin{itemize}
\item Confusing soil texture (mineral particle size) with soil structure (arrangement of particles into aggregates).
\item Stating that legumes "do not need nitrogen" — they fix their own but still require nitrogen for growth; the fixation supplies it.
\item Forgetting that ruminants rely on microbial cellulase, not their own enzyme.
\item In Pearson square, subtracting the wrong way (always subtract diagonally: desired minus feed, taking absolute values).
\item Listing cultivation operations out of chronological order (e.g., fertiliser before planting, or weeding after harvest).
\item Vague disease control measures: "use chemicals" instead of naming specific fungicides/insecticides and application timing.
\end{itemize}
\end{attention}

\newpage

\coursec{Solutions to the exercises}

\begin{corrige}[Exercise 1 — Multiple choice questions]
\begin{enumerate}
\item \textbf{Answer: a) the A horizon.}\par
The A horizon (topsoil) is the zone where humus accumulates and mixes with mineral particles. It is dark in colour, biologically very active (roots, earthworms, micro-organisms) and supplies most of the nutrients and water absorbed by crops. The B horizon is mainly a zone of accumulation of leached minerals, the C horizon is weathered parent rock, and the bedrock is unweathered rock.
\item \textbf{Answer: c) magnesium.}\par
Magnesium is the central atom of the chlorophyll molecule. Because Mg is mobile in the plant, it is withdrawn from older leaves to supply young ones, so deficiency appears first as \cle{interveinal chlorosis} (yellowing between the veins, which remain green) on the \textbf{older} leaves. Nitrogen deficiency also causes yellowing of older leaves, but it is general, not interveinal, and N is not a component of the chlorophyll ring itself.
\item \textbf{Answer: a) 19\%.}\par
Pearson square: desired CP = 16\%; feeds 9\% (maize) and 45\% (groundnut cake).
\begin{align*}
\text{Parts of maize} &= 45 - 16 = 29\\
\text{Parts of groundnut cake} &= 16 - 9 = 7\\
\text{Total parts} &= 29 + 7 = 36\\
\text{Groundnut cake} &= \frac{7}{36} \times 100 \approx 19.4\% \approx 19\%
\end{align*}
\item \textbf{Answer: b) a virus.}\par
Newcastle disease of poultry is caused by an avian paramyxovirus (a virus). This is why antibiotics are useless against it and control relies on \textbf{vaccination} and biosecurity.
\end{enumerate}
\end{corrige}

\begin{corrige}[Exercise 2 — True or false — justify]
\begin{enumerate}
\item \textbf{False.} Clay soils retain water and nutrients well, but they are heavy, poorly drained, poorly aerated and difficult to till, and they become waterlogged in the rains and hard in the dry season. A \textbf{loam} (balanced mixture of sand, silt and clay) is generally the most fertile and easiest soil to farm, so a clay soil is not \emph{always} better than a sandy one.
\item \textbf{True.} Legumes live in symbiosis with \emph{Rhizobium} bacteria housed in root nodules; the bacteria reduce atmospheric nitrogen (\ce{N2}) to ammonium compounds that the plant can use. The legume therefore obtains its nitrogen supply without mineral nitrogen fertiliser (though it still needs the other mineral elements, especially P and K).
\item \textbf{False.} Ruminants do \textbf{not} secrete cellulase themselves. Cellulose digestion is carried out by symbiotic micro-organisms (bacteria, protozoa, fungi) living in the rumen, which produce cellulase and ferment cellulose into volatile fatty acids that the animal absorbs.
\item \textbf{True.} In extensive production (e.g. free-range grazing on the Adamawa plateau) animals receive little supplementary feed and limited health care, so growth rate, milk yield and offtake per animal are lower than in intensive systems, where feed, housing and veterinary care are optimised. (However, output \emph{per hectare of labour/capital invested} may be lower in extensive systems, which are cheaper to run.)
\end{enumerate}
\end{corrige}

\begin{corrige}[Exercise 3 — Definitions]
\begin{enumerate}
\item \textbf{Soil profile:} a vertical section through the soil, from the surface down to the parent rock, showing the sequence of horizontal layers called \cle{horizons} (O, A, B, C) which differ in colour, texture, structure and composition.
\item \textbf{Humus:} the dark, colloidal, partially decomposed organic matter of the soil, formed by the action of soil organisms on dead plant and animal remains; it improves soil structure, water retention and nutrient supply.
\item \textbf{Macronutrient:} a mineral element required by plants in relatively \textbf{large} quantities for normal growth, e.g. nitrogen, phosphorus, potassium, calcium, magnesium and sulphur (as opposed to micronutrients/trace elements such as Fe, Mn, Zn, Cu, B, Mo, needed in very small amounts).
\item \textbf{Ration:} the total quantity of feed given to an animal in a specified period (usually 24 hours), formulated to supply all the nutrients (energy, protein, minerals, vitamins, water) needed for maintenance and production (growth, milk, eggs).
\item \textbf{Biosecurity:} the set of preventive measures (quarantine of new animals, disinfection, control of visitors and vehicles, vaccination, hygiene of housing and equipment, safe disposal of carcasses) designed to prevent the introduction and spread of disease-causing agents on a farm.
\end{enumerate}
\end{corrige}

\begin{corrige}[Exercise 4 — Direct calculations]
\begin{enumerate}
\item Sand-to-clay ratio:
\[ \frac{\text{sand}}{\text{clay}} = \frac{40}{20} = 2 \quad \text{(i.e. } 2:1\text{)}. \]
With 40\% sand and only 20\% clay, the soil is a \textbf{loam} tending to a sandy loam. It will drain \textbf{moderately to fairly quickly}: the high sand fraction creates large pores allowing water to percolate, while the 20\% clay still retains some water and nutrients.
\item Mass of nitrogen applied:
\[ m(\mathrm{N}) = 100\ \mathrm{kg} \times \frac{46}{100} = \SI{46}{kg} \text{ of N per hectare}. \]
\item Mass of protein consumed per day:
\[ m(\text{protein}) = 150\ \mathrm{g} \times \frac{21}{100} = \SI{31.5}{g} \text{ of crude protein per day}. \]
\end{enumerate}
\end{corrige}

\begin{corrige}[Exercise 5 — Using the texture triangle]
\begin{enumerate}
\item On the texture triangle, the point corresponding to 70\% sand, 20\% silt and 10\% clay falls in the class \textbf{sandy loam} (close to the boundary with loamy sand). The soil is therefore a sandy loam.
\item Physical properties of a sandy loam:
\begin{itemize}
\item \textbf{Drainage:} good to rapid — the large pores between sand particles let excess water percolate freely, so waterlogging is rare.
\item \textbf{Water retention:} low to moderate — little water is held against gravity, so the soil dries out quickly and crops may suffer in dry spells.
\item \textbf{Ease of tillage:} very easy — the soil is light, crumbly and workable even with simple hand tools or animal traction; it also warms up quickly.
\end{itemize}
\item Suitable crops (any two):
\begin{itemize}
\item \textbf{Groundnut:} needs a light, well-drained soil so that the pegs can penetrate and the pods develop without rotting.
\item \textbf{Cassava:} tolerates low-fertility, well-drained sandy soils and gives reasonable tuber yields where heavier crops would fail. (Accept also: sweet potato, watermelon, pineapple, early maize with fertiliser.)
\end{itemize}
\item Measures to improve fertility (any two):
\begin{itemize}
\item Add \textbf{organic matter} (farmyard manure, compost, mulch, green manure) to increase water and nutrient retention and feed soil organisms.
\item Apply \textbf{mineral fertilisers} (NPK) in split doses to compensate for leaching losses.
\item Practise \textbf{crop rotation with a legume} (cowpea, groundnut) to enrich the soil in nitrogen.
\end{itemize}
\end{enumerate}
\end{corrige}

\begin{corrige}[Exercise 6 — Mineral elements in plants]
\begin{center}
\begin{tabularx}{\linewidth}{|l|X|X|}
\hline
\rowcolor{popblueL} \textbf{Element} & \textbf{One function} & \textbf{One deficiency symptom} \\
\hline
Nitrogen (N) & Component of proteins, nucleic acids and chlorophyll; promotes vegetative (leafy) growth & General yellowing (chlorosis) of older leaves; stunted growth \\
\hline
Phosphorus (P) & Component of ATP, nucleic acids and membranes; promotes root development, flowering and seed formation & Poor root growth; dark green or purplish coloration of older leaves; delayed maturity \\
\hline
Potassium (K) & Activates enzymes; regulates stomatal opening and water balance; improves disease resistance and fruit quality & Yellowing and scorching (necrosis) of the margins of older leaves; weak stems \\
\hline
Magnesium (Mg) & Central atom of the chlorophyll molecule; enzyme activator & Interveinal chlorosis of older leaves (veins stay green) \\
\hline
Iron (Fe) & Needed for chlorophyll synthesis (component of electron carriers such as cytochromes) & Interveinal chlorosis of \textbf{young} leaves (Fe is immobile in the plant) \\
\hline
\end{tabularx}
\end{center}
\textbf{Examiner's note:} the key distinction is the \emph{position} of the symptoms: deficiencies of mobile elements (N, P, K, Mg) appear first on \textbf{old} leaves, whereas deficiencies of immobile elements (Fe, Ca) appear first on \textbf{young} leaves.
\end{corrige}

\begin{corrige}[Exercise 7 — Balancing a ration with the Pearson square]
\begin{enumerate}
\item Pearson square (desired CP = 20\%; maize = 9\%; fish meal = 60\%):
\begin{center}
\begin{tikzpicture}[scale=1.1]
\node at (0,2) {maize 9};
\node at (0,0) {fish meal 60};
\node at (2,1) {\textbf{20}};
\node[text=popblue] at (4,2) {$60-20=40$ parts maize};
\node[text=poporange] at (4,0) {$20-9=11$ parts fish meal};
\draw[thick] (0.9,2) -- (3.1,0.2);
\draw[thick] (1.1,0) -- (3.1,1.8);
\end{tikzpicture}
\end{center}
Subtractions are done diagonally, taking positive differences:
\begin{align*}
\text{Parts of maize} &= 60 - 20 = 40\\
\text{Parts of fish meal} &= 20 - 9 = 11\\
\text{Total parts} &= 40 + 11 = 51\\
\text{Maize} &= \frac{40}{51} \times 100 \approx 78.4\%\\
\text{Fish meal} &= \frac{11}{51} \times 100 \approx 21.6\%
\end{align*}
\item For \SI{100}{kg} of ration:
\begin{align*}
m(\text{maize}) &= \frac{40}{51} \times 100 \approx \SI{78.4}{kg}\\
m(\text{fish meal}) &= \frac{11}{51} \times 100 \approx \SI{21.6}{kg}
\end{align*}
Check: protein supplied $= 78.4 \times 0.09 + 21.6 \times 0.60 \approx 7.06 + 12.96 = 20.0\ \mathrm{kg}$, i.e. 20\% of \SI{100}{kg}. \checkmark
\item Fish meal is used in a limited proportion because:
\begin{itemize}
\item it is \textbf{expensive} — increasing its share would raise the cost of the ration and reduce profit;
\item in excess it can give an undesirable \textbf{fishy taste} to the meat and reduce palatability of the feed;
\item it also supplies minerals (Ca, P), so beyond the requirement it would unbalance the mineral content of the ration.
\end{itemize}
\end{enumerate}
\end{corrige}

\begin{corrige}[Exercise 8 — Cultivation of a food crop (maize)]
\begin{enumerate}
\item and \textbf{2.} Operations in chronological order, with their purposes:
\begin{center}
\begin{tabularx}{\linewidth}{|l|X|}
\hline
\rowcolor{popblueL} \textbf{Operation} & \textbf{Purpose} \\
\hline
1. Land clearing & Remove bushes and previous crop residues to expose the soil \\
\hline
2. Ploughing / tilling & Loosen the soil, bury residues, improve aeration and root penetration \\
\hline
3. Harrowing / levelling & Break clods and produce a fine seedbed for good seed--soil contact \\
\hline
4. Planting (at onset of rains; spacing about $75 \times 25\ \mathrm{cm}$, 2--3 seeds per hole) & Establish the crop at the right density and depth for good germination \\
\hline
5. Thinning and gap filling & Keep one strong plant per stand and replace missing plants to maintain optimum population \\
\hline
6. Fertiliser application (NPK at planting, urea at knee height) & Supply N, P, K for vigorous growth and good grain filling \\
\hline
7. Weeding (2--3 times) & Remove weeds that compete for light, water and nutrients \\
\hline
8. Pest and disease control & Protect the crop from stem borers, armyworms, etc. \\
\hline
9. Harvesting (when husks are dry) & Collect mature cobs before losses to pests, rot or shattering \\
\hline
10. Shelling, drying and storage (in dry cribs/bags, sometimes with protectant) & Reduce moisture to prevent mould and weevil attack during storage \\
\hline
\end{tabularx}
\end{center}
\item Two common pests/diseases of maize in Cameroon (any two):
\begin{itemize}
\item \textbf{Fall armyworm / stem borer (pests):} control by early planting, field hygiene and application of a recommended insecticide into the funnel.
\item \textbf{Maize streak virus (disease):} control by using resistant/tolerant varieties and controlling the leafhopper vector.
\item \textbf{Storage weevils:} control by proper drying and use of a grain protectant or airtight storage.
\end{itemize}
\end{enumerate}
\end{corrige}

\begin{corrige}[Exercise 9 — Structured question: soils and fertility (20 marks)]
\textbf{Indicative mark scheme and model answers.}
\begin{enumerate}
\item \textbf{Textural classes (4 marks — 2 per soil).} Using the texture triangle:
\begin{itemize}
\item Soil A (85\% sand, 10\% silt, 5\% clay): \textbf{sand} (sandy soil). \hfill (2)
\item Soil B (20\% sand, 25\% silt, 55\% clay): \textbf{clay} (clay soil). \hfill (2)
\end{itemize}
\item \textbf{Comparison (6 marks — 1 per valid comparative point).}
\begin{itemize}
\item \textbf{Drainage:} A drains rapidly (large pores between sand grains); B drains slowly and is easily waterlogged.
\item \textbf{Aeration:} A is well aerated; B is poorly aerated, especially when wet.
\item \textbf{Water-holding capacity:} A holds little water and dries quickly; B holds a large amount of water (small pores retain it by capillarity).
\end{itemize}
\item \textbf{Explanation (4 marks).}
\begin{itemize}
\item A warms quickly because it holds little water: little heat is wasted evaporating water, so solar energy rapidly raises the temperature of the dry sand particles. \hfill (2)
\item It loses nutrients by leaching because its large pores let rainwater percolate rapidly, and with very little clay and low humus there are few negatively charged colloids to retain nutrient cations (\ce{K+}, \ce{Ca^2+}, \ce{Mg^2+}, \ce{NH4+}), which are therefore washed down beyond the root zone. \hfill (2)
\end{itemize}
\item \textbf{Crop and improvement for each soil (6 marks — 1.5 per item, justification required).}
\begin{itemize}
\item \textbf{Soil A:} suitable crop: groundnut, cassava or watermelon, which tolerate light, well-drained, low-fertility soils (groundnut pods also need loose soil to develop). Improvement: regular addition of organic matter (compost/manure) to raise water and nutrient retention, plus split applications of NPK fertiliser to limit leaching losses.
\item \textbf{Soil B:} suitable crop: swamp rice (paddy), which tolerates — indeed requires — waterlogged conditions, or yam/maize on ridges. Improvement: drainage or ridging to improve aeration, addition of organic matter and liming if acid, to improve structure and workability.
\end{itemize}
\end{enumerate}
\textbf{Examiner expectations:} correct use of the texture triangle; comparative (not just descriptive) statements in (b); a causal link between low water content / low colloid content and the two properties in (c); justified choices in (d) — a crop or method without justification earns only half the marks.
\end{corrige}

\begin{corrige}[Exercise 10 — Essay: cultivation of cocoa (20 marks)]
\textbf{Indicative mark scheme:} sequence and coverage 8 marks; technical accuracy 8 marks; stated importance of each operation 4 marks.\par
\textbf{Model plan and content.}
\begin{enumerate}
\item \textbf{Choice of site and soil:} cocoa requires a hot, humid climate (about 1500--2000 mm of well-distributed rainfall) and a deep, well-drained, humus-rich soil (loamy or volcanic), slightly acidic. Importance: waterlogging causes root rot and black pod, while shallow soils dry out and stress the trees.
\item \textbf{Land preparation:} clear the bush but retain or plant some tall trees for shade; line and peg at about $3 \times 3\ \mathrm{m}$. Importance: shade protects young plants from sun scorch and reduces evaporation; correct spacing ensures good canopy development and easy maintenance.
\item \textbf{Planting and shade management:} plant seedlings (from nursery) or seeds at stake at the start of the rains, with temporary shade (plantain, banana, \emph{Gliricidia}); reduce shade progressively as the canopy closes. Importance: good establishment determines the future yield of the plantation.
\item \textbf{Maintenance:}
\begin{itemize}
\item \emph{Weeding} (2--3 times per year): removes competition for water and nutrients.
\item \emph{Pruning:} removes dead, diseased and chupon (vertical) shoots to shape the tree, aerate the canopy and ease harvesting.
\item \emph{Pest and disease control:} capsids (mirids) controlled with recommended insecticides; black pod disease (\emph{Phytophthora}) controlled by regular harvesting of ripe pods, removal of diseased pods, pruning for aeration and copper-based fungicides.
\item \emph{Fertiliser application:} NPK fertiliser maintains yield on old or exhausted soils.
\end{itemize}
\item \textbf{Harvesting:} ripe pods (yellow/orange) are cut with a sharp knife or hooked pole every 2--3 weeks, without damaging the cushion. Importance: regular, careful harvesting limits black pod losses and protects future flowers on the same cushion.
\item \textbf{Breaking and fermentation:} pods are broken and the beans (with mucilage) heaped in boxes or banana-leaf piles for about 5--6 days, with turning. Importance: fermentation kills the embryo, removes the mucilage and develops the chocolate flavour and brown colour — it determines the commercial quality (grade).
\item \textbf{Drying:} beans are dried in the sun on trays/mats (or in dryers) to about 7--8\% moisture, with regular turning. Importance: proper drying prevents mould and off-flavours during storage and transport.
\item \textbf{Marketing:} dried beans are sorted, bagged and sold to cooperatives or licensed buying agents, then exported (e.g. through the port of Douala). Importance: grading and proper storage preserve quality and secure the best price.
\end{enumerate}
\textbf{Examiner expectations:} operations in the correct chronological order; at least one precise technical detail per step (spacing, fermentation duration, moisture target, named pest/disease); the \emph{importance} of each operation explicitly stated.
\end{corrige}

\begin{corrige}[Exercise 11 — Problem situation and essay: animal production (20 marks)]
\textbf{Part A (10 marks).}
\begin{enumerate}
\item \textbf{Most likely disease (2 marks):} \textbf{Newcastle disease}, caused by an avian \textbf{paramyxovirus} (a virus). The combination of respiratory signs (gasping, nasal discharge), nervous signs (twisted neck / torticollis) and high mortality in young chicks is characteristic. \hfill (2)
\item \textbf{Spread within the flock (3 marks — any three points):}
\begin{itemize}
\item by direct contact between sick and healthy birds;
\item through contaminated air/droplets (respiratory secretions);
\item through contaminated feed, water, litter and equipment (faeces contain the virus);
\item by people (attendants, visitors), vehicles and introduction of infected birds from other farms or live-bird markets.
\end{itemize}
\item \textbf{Vaccination and biosecurity programme (5 marks):}
\begin{itemize}
\item \emph{Vaccination:} vaccinate all chicks against Newcastle disease (e.g. first dose in the first days of life by eye drop or drinking water, with boosters according to the veterinary schedule); vaccinate against other local diseases (Gumboro, fowl pox) as advised; keep vaccines cold (cold chain) and follow the calendar strictly.
\item \emph{Biosecurity:} quarantine new birds before introduction; restrict and disinfect visitors, vehicles and equipment (footbaths at the entrance); provide clean water and feed; clean and disinfect the house between batches (all-in/all-out system); isolate sick birds and burn/bury dead birds; avoid mixing ages and species.
\end{itemize}
\end{enumerate}
\textbf{Part B (10 marks).} Indicative content:
\begin{itemize}
\item \textbf{Extensive system} (e.g. Fulani cattle herds grazing the Adamawa plateau): \emph{advantages} — low capital and feed costs (natural pasture), hardy local breeds (e.g. Gudali) adapted to the environment, large areas usable, animals convert grass inedible to humans into meat and milk. \emph{Disadvantages} — low productivity per animal (slow growth, low milk yield, late maturity), losses to diseases (trypanosomiasis, tick-borne diseases) and parasites, overgrazing and soil degradation, seasonal weight loss in the dry season, farmer--grazer conflicts, difficulty of veterinary follow-up.
\item \textbf{Intensive system} (zero-grazing, feedlots, dairy units): \emph{advantages} — high output per animal (improved breeds, balanced rations), better disease control and record keeping, production all year round, efficient use of small land areas. \emph{Disadvantages} — high investment and running costs (feed, housing, veterinary care), need for skilled labour, rapid spread of disease in dense housing if biosecurity fails, dependence on purchased inputs.
\item \textbf{Improvement / combination (conclusion):} adopt \emph{semi-intensive} systems: keep herds on improved or sown pastures with supplementary feeding (crop residues, agro-industrial by-products, mineral licks) especially in the dry season; provide vaccination and dipping against ticks; improve breeds by crossing local hardy cattle with higher-yielding breeds; develop water points and fodder banks on the Adamawa; organise ranching with paddocks to prevent overgrazing; strengthen veterinary services and farmer training. This combines the low cost of extensive grazing with the productivity gains of intensive management, increasing meat and milk supply sustainably.
\end{itemize}
\textbf{Examiner expectations:} Part A — correct identification of the disease and agent, concrete modes of transmission, and a programme mentioning \emph{both} a vaccination schedule \emph{and} named biosecurity measures. Part B — balanced discussion (at least two advantages and two disadvantages per system), explicit reference to the Adamawa context, and a reasoned, realistic conclusion.
\end{corrige}

\newpage

\coursec{Summary sheet}

\begin{popfigure}
\begin{center}\tcbox[colback=popdark,colframe=popdark,coltext=white,arc=9pt,boxsep=4pt,left=6pt,right=6pt]{\bfseries\small AGRICULTURE (BIO08)}\end{center}
\vspace{-2pt}
\begin{tcbraster}[raster columns=3,raster equal height=rows,raster column skip=6pt,raster row skip=6pt,enhanced,blankest]
\begin{tcolorbox}[enhanced,colback=white,colframe=popblue,boxrule=0.9pt,arc=3pt,title={Soils: types, formation, fertility},fonttitle=\bfseries\scriptsize,coltitle=white,colbacktitle=popblue,left=4pt,right=4pt,top=2pt,bottom=2pt,boxsep=1pt,toptitle=1.5pt,bottomtitle=1.5pt,halign=left]
\begin{itemize}[nosep,leftmargin=*,itemsep=1pt,font=\scriptsize]
\item Sandy, clay, loam, humus
\item Weathering, profile horizons
\end{itemize}
\end{tcolorbox}
\begin{tcolorbox}[enhanced,colback=white,colframe=popgreen,boxrule=0.9pt,arc=3pt,title={Plant nutrition},fonttitle=\bfseries\scriptsize,coltitle=white,colbacktitle=popgreen,left=4pt,right=4pt,top=2pt,bottom=2pt,boxsep=1pt,toptitle=1.5pt,bottomtitle=1.5pt,halign=left]
\begin{itemize}[nosep,leftmargin=*,itemsep=1pt,font=\scriptsize]
\item Macro: N, P, K, Ca, Mg, S
\item Micro: Fe, Mn, B, Zn, Cu, Mo
\end{itemize}
\end{tcolorbox}
\begin{tcolorbox}[enhanced,colback=white,colframe=poporange,boxrule=0.9pt,arc=3pt,title={Cultivation techniques},fonttitle=\bfseries\scriptsize,coltitle=white,colbacktitle=poporange,left=4pt,right=4pt,top=2pt,bottom=2pt,boxsep=1pt,toptitle=1.5pt,bottomtitle=1.5pt,halign=left]
\begin{itemize}[nosep,leftmargin=*,itemsep=1pt,font=\scriptsize]
\item Food crops: maize, cassava
\item Cash crops: cocoa, coffee, cotton
\end{itemize}
\end{tcolorbox}
\begin{tcolorbox}[enhanced,colback=white,colframe=poppurple,boxrule=0.9pt,arc=3pt,title={Animal nutrition},fonttitle=\bfseries\scriptsize,coltitle=white,colbacktitle=poppurple,left=4pt,right=4pt,top=2pt,bottom=2pt,boxsep=1pt,toptitle=1.5pt,bottomtitle=1.5pt,halign=left]
\begin{itemize}[nosep,leftmargin=*,itemsep=1pt,font=\scriptsize]
\item Carbs, proteins, lipids, vitamins, minerals, water
\item Rations, deficiency signs
\end{itemize}
\end{tcolorbox}
\begin{tcolorbox}[enhanced,colback=white,colframe=popgold,boxrule=0.9pt,arc=3pt,title={Animal production},fonttitle=\bfseries\scriptsize,coltitle=white,colbacktitle=popgold,left=4pt,right=4pt,top=2pt,bottom=2pt,boxsep=1pt,toptitle=1.5pt,bottomtitle=1.5pt,halign=left]
\begin{itemize}[nosep,leftmargin=*,itemsep=1pt,font=\scriptsize]
\item Systems: extensive, intensive
\item Breeding, health, housing
\end{itemize}
\end{tcolorbox}
\end{tcbraster}

\legende{Mind map of Chapter BIO08 -- Agriculture: five interconnected branches from soil to animal production.}
\end{popfigure}

\begin{bilan}

\textbf{1. Key definitions}
\begin{itemize}
\item \cle{Soil}: the upper layer of the Earth's crust, formed by weathering of parent rock plus decomposition of organic matter, which supports plant growth.
\item \cle{Soil profile}: vertical section showing horizons O (organic), A (topsoil), B (subsoil), C (weathered parent material), R (bedrock).
\item \cle{Soil texture}: relative proportions of sand, silt and clay particles; determines drainage, aeration and water retention.
\item \cle{Loam}: balanced mixture of sand, silt and clay with humus -- the most fertile soil for most crops.
\item \cle{Soil fertility}: capacity of a soil to supply water, air, anchorage and mineral nutrients for good crop growth.
\item \cle{Macronutrients}: elements needed in large amounts (N, P, K, Ca, Mg, S); \cle{micronutrients} (trace elements) needed in tiny amounts (Fe, Mn, B, Zn, Cu, Mo, Cl).
\item \cle{Fertiliser}: material (organic -- manure, compost; or inorganic -- NPK, urea) added to soil to restore nutrients.
\item \cle{Food crop}: grown mainly for human consumption (maize, cassava, plantain, yam); \cle{cash (non-food) crop}: grown for sale/industry (cocoa, coffee, cotton, oil palm).
\item \cle{Balanced ration}: feed supplying all an animal's nutrients in correct proportions for maintenance and production.
\item \cle{Extensive system}: animals range freely with low inputs; \cle{intensive system}: animals confined, high inputs, high yield per head.
\end{itemize}

\textbf{2. Facts, comparisons and ``laws'' to know}
\begin{itemize}
\item Soil formation = \cle{weathering} (physical: temperature, rain, wind; chemical: hydrolysis, oxidation; biological: roots, microbes, earthworms) + humus accumulation.
\item Texture comparison: \textbf{sandy} soils drain fast, warm quickly, poor nutrient retention; \textbf{clay} soils retain water and nutrients but are heavy, poorly aerated, waterlog easily; \textbf{loam} is intermediate and best.
\item Soil pH: most crops prefer $6.0$--$7.0$; acidity corrected by \cle{liming} (adding \ce{CaCO3}).
\item Roles of key elements: \textbf{N} -- leaf/vegetative growth (deficiency: yellowing, chlorosis); \textbf{P} -- roots, energy transfer ATP, flowering (deficiency: purpling, poor roots); \textbf{K} -- enzyme activation, water regulation, disease resistance (deficiency: scorched leaf edges); \textbf{Mg} -- chlorophyll component (deficiency: interveinal chlorosis); \textbf{Fe} -- chlorophyll synthesis (deficiency: chlorosis of young leaves).
\item Legumes (groundnut, cowpea, soya) fix atmospheric nitrogen via \emph{Rhizobium} in root nodules -- basis of crop rotation.
\item Fertility maintenance: crop rotation, fallowing, mulching, manuring, mixed cropping, contour ploughing against erosion.
\item Cultivation steps: land clearing $\rightarrow$ ploughing/tilling $\rightarrow$ harrowing $\rightarrow$ ridging $\rightarrow$ planting (correct spacing/depth) $\rightarrow$ weeding, fertilising, pest control $\rightarrow$ harvesting $\rightarrow$ storage.
\item Animal feed classes: carbohydrates (energy), proteins (growth, repair, milk/egg production), lipids (concentrated energy), vitamins and minerals (health, bones: Ca and P), water (essential, often forgotten).
\item Ruminants (cattle, goats, sheep) digest cellulose via rumen microbes; non-ruminants (pigs, poultry) need higher-quality protein and cannot use roughage efficiently.
\item Animal production management: selection and breeding, housing/hygiene, vaccination and deworming, balanced feeding, record keeping.
\end{itemize}

\textbf{3. Methods to master}
\begin{itemize}
\item \textbf{Soil texture test (feel method)}: moisten a sample, rub between fingers -- gritty = sandy; smooth and sticky, moulds into a ball = clay; slightly gritty, holds shape loosely = loam.
\item \textbf{Sedimentation test}: shake soil with water in a jar, let settle -- sand sinks first (bottom layer), then silt, clay stays suspended longest; humus floats.
\item \textbf{Water retention/drainage experiment}: equal volumes of sand, clay, loam in funnels with filter paper; pour equal water; compare volume drained and time.
\item \textbf{Planning a crop enterprise}: choose site and crop suited to soil/climate $\rightarrow$ prepare land $\rightarrow$ select viable seed $\rightarrow$ plant at correct season $\rightarrow$ manage $\rightarrow$ harvest and store.
\end{itemize}

\textbf{4. Common mistakes to avoid}
\begin{itemize}
\item Confusing \cle{soil texture} (particle sizes) with \cle{soil structure} (arrangement into crumbs).
\item Saying fertilisers ``feed'' plants -- plants make their own food by photosynthesis; fertilisers supply \emph{mineral elements}.
\item Believing more fertiliser always means more yield -- excess causes toxicity, leaching and water pollution (eutrophication).
\item Confusing macronutrients with ``most important'' -- trace elements are equally essential, only needed in smaller quantities.
\item Mixing up extensive and intensive systems, or assuming intensive is always better (cost, disease risk, waste disposal).
\item Forgetting water as a nutrient for animals, or ignoring the special digestion of ruminants.
\end{itemize}

\textbf{5. GCE tips}
\begin{itemize}
\item Draw and label the \cle{soil profile} (O, A, B, C, R) neatly -- a frequent diagram question.
\item Learn deficiency symptoms of N, P, K, Mg, Fe in a table -- examiners love ``state and explain'' questions.
\item Use Cameroonian examples: cocoa and coffee in the South-West and Centre, cotton in the North, maize in the West highlands, cattle on the Adamawa plateau.
\item In essays on soil fertility, always mention \emph{both} organic methods (manure, rotation, fallow) and inorganic methods (NPK, liming).
\item For animal production questions, structure answers under: feeding, housing, health, breeding, records.
\item Define every technical term you use (e.g.\ ration, weeding, mulching) -- marks are awarded for precise vocabulary.
\end{itemize}

\end{bilan}

\findecours

\end{document}
